\chapter{二维拓扑场论中的 D 膜与 K 理论}\label{chap:2}

\SourceParagraph{ch02-intro-001}{37}
让我们从尽可能简单的背景，也就是二维拓扑场论，简称 TFT，开始研究 Dirichlet 膜。正如 \BookBibText{ref242}{MS1} 所解释，并且如我们将要回顾的那样，二维闭弦 TFT 就是一个有限维交换 Frobenius 代数 $\mathcal{C}$。本章的目标，是将这一结果推广到二维开闭弦拓扑场论。

\SourceParagraph{ch02-intro-002}{37}
从带边界的一维流形之间的配边所构成的几何范畴出发，我们将在\BookSectionRef{sec:2.1}{2.1}中说明，描述缝合关系及其解，是一个非平凡但可以处理的问题。粗略地说，结果是 D 膜，也就是边界条件，对应于闭弦代数 $\mathcal{C}$ 上的模。

\SourceParagraph{ch02-intro-003}{37}
这一结果的一个推论，是 D 膜与 K 理论之间的一般关系。这一关系将在全书中发挥核心作用，后文将给出各种数学和物理论证，特别是在\BookSectionRef{sec:5.1}{5.1}中。不过，在这里就能看出，它的起源远为基础：归根结底，它仅仅来自缝合约束。

\SourceParagraph{ch02-intro-004}{37}
如果 $\mathcal{C}$ 是半单的，就能进一步对 D 膜作完整分类，我们将在\BookSectionRef{sec:2.2}{2.2}和\BookSectionRef{sec:2.3}{2.3}中完成这件事。这刻画了零维靶空间的物理。

\SourceParagraph{ch02-intro-005}{37}
在\BookSectionRef{sec:2.4}{2.4}和\BookSectionRef{sec:2.5}{2.5}中，我们开始逐步处理更高维的例子，其中 $\mathcal{C}$ 不再半单。也许最简单的例子是“Landau--Ginzburg 模型”：$\mathcal{C}$ 是由函数模去梯度所生成的理想而得到的 Jacobian 代数。若要进一步发展，可以考察以靶空间的上同调为基础的自然代数。

\SourceParagraph{ch02-intro-006}{37}
在\BookSectionRef{sec:2.6}{2.6}中，我们将这些结果推广到等变情形：给定一个有限群 $G$，并让世界面成为带 $G$ 丛的曲面。这与轨形中 D 膜的分类有关。

\BookUnsectionedSubsection{2.0.1}{结果概述}\label{sec:2.0.1}

\SourceParagraph{ch02-2.0.1-001}{37}
我们的主要结果是下面两个定理；现在先陈述它们，其含义将在后文逐渐清楚。为了表述第一个定理，必须指出，半单交换 Frobenius 代数\footnote{Frobenius 代数 $\mathcal{C}$，无论是否交换，都是一个复数域上的代数，配备一个线性映射 $\theta\colon\mathcal{C}\to\mathbb{C}$，使配对 $(a,a')\mapsto\theta(aa')$ 是 $\mathcal{C}$ 上的非退化双线性形式。经典例子是有限群 $G$ 的复群代数 $\mathbb{C}[G]$：对于群元素的线性组合 $a=\sum\lambda_g g$，定义 $\theta(a)=\lambda_1$。另一个例子是紧定向流形的复系数上同调代数。关于 Frobenius 代数的基础材料，例如可参见文献\BookCite{ref104}的第9章，或文献\BookCite{ref142}。}$\mathcal{C}$ 自动就是有限集合 $X=\operatorname{Spec}(\mathcal{C})$ 上的复值函数代数，这个集合由从 $\mathcal{C}$ 到 $\mathbb{C}$ 的代数同态组成。对于有限维代数，这与代数几何中通常的 $\operatorname{Spec}$ 定义一致。我们把 $X$ 看成一个“时空”，它带有“体积形式”或“伸缩子场”$\theta$，为每个点 $x\in X$ 指定测度 $\theta_x$。

\begin{lettertheorem}{A}\label{thm:2.A}
\SourceParagraph{ch02-thm-A-001}{38}
对于与有限时空 $(X,\theta)$ 对应的半单二维 TFT，选定一个极大的 D 膜范畴，就为 $X$ 的每个点 $x$ 固定了 $\theta_x$ 的一个平方根。边界条件范畴等价于 $X$ 上有限维复向量丛的范畴 $\operatorname{Vect}(X)$。不过，这一对应不是典范的；它可以任意地与一个等价 $\operatorname{Vect}(X)\to\operatorname{Vect}(X)$ 复合，其中这一等价由把每个向量丛与一个固定线丛作张量积给出；“固定”意味着，该线丛不依赖于所考虑的具体 D 膜。

\SourceParagraph{ch02-thm-A-002}{38}
反过来，每个半单 Frobenius 范畴 $\mathcal{B}$ 都是一个典范二维 TFT 的边界条件范畴；这个 TFT 所对应的交换 Frobenius 代数，是 $\mathcal{B}$ 的恒等函子的自同态环。
\end{lettertheorem}

\SourceParagraph{ch02-2.0.1-002}{38}
下一节将解释边界条件在何种意义下构成范畴。这个定理将在\BookSectionRef{sec:2.2}{2.2}中证明。\BookSectionRef{sec:2.2.4}{2.2.4}将描述它对自旋理论的类似版本。

\SourceParagraph{ch02-2.0.1-003}{38}
第二个定理涉及“$G$ 等变”或“$G$ 规范化”的 TFT，其中 $G$ 为有限群。Turaev 已证明，在二维情形下，半单 $G$ 等变 TFT 对应于一个有限时空 $X$，其上有给定的 $G$ 作用，并带有 $G$ 不变的伸缩子场 $\theta$，以及代表等变上同调群 $H_G^3(X,\mathbb{Z})$ 中一个元素的“B 场”$B$。

\begin{lettertheorem}{B}\label{thm:2.B}
\SourceParagraph{ch02-thm-B-001}{38}
对于与有限时空 $(X,\theta,B)$ 对应的半单 $G$ 等变 TFT，选定一个极大的 D 膜范畴，就像前面一样固定了平方根 $\sqrt{\theta_x}$ 的一个 $G$ 不变选择；此时，这个范畴等价于 $X$ 上有限维的、经 B 场扭曲的 $G$ 等变向量丛范畴，其不确定性仅在于整体与一个 $G$ 等变线丛作张量积。

\SourceParagraph{ch02-thm-B-002}{38}
在这一情形下，D 膜范畴等价于对规范场积分后，从规范化理论得到的“轨形”理论的 D 膜范畴，并且它不再记住产生这个轨形理论的等变理论。不过，存在一种自然的充实结构，能够记住该等变理论。
\end{lettertheorem}

\SourceParagraph{ch02-2.0.1-004}{39}
\BookSectionRef{sec:2.6}{2.6}将解释并证明这一点。

\SourceParagraph{ch02-2.0.1-005}{39}
当这些定理适用时，它们为我们的主要问题提供了完整答案。不过，正如将要看到的那样，半单情形这一限制，使它们的适用范围十分有限，实质上只适用于零维靶空间。尽管如此，仍然值得详细讨论，因为这是迄今理解理论基本结构最简单的方法。在后续章节中，我们将发展超越半单情形所需的形式体系。这里先在\BookSectionRef{sec:2.1}{2.1}和\BookSectionRef{sec:2.5}{2.5}中预示这一发展，说明边界条件范畴如何自然地成为 Fukaya、Kontsevich 等人意义下的 $A_\infty$ 范畴。

\section{缝合定理}\label{sec:2.1}
\subsection{二维开闭弦 TFT 的定义}\label{sec:2.1.1}

\SourceParagraph{ch02-2.1.1-001}{39}
粗略地说，$d$ 维量子场论是一种特殊的规则：它为每个具有适当结构的闭 $d$ 维流形指定一个数，称为理论的配分函数。使这一规则成为量子场论的，是闭流形被分割时配分函数的变化方式。在数学上，可以方便地把这种结构形式化为从几何范畴到线性范畴的函子。最简单的例子是拓扑场论：这里选取的几何范畴，以闭、定向的 $(d-1)$ 维流形为对象，以定向配边为态射；如果两个配边之间存在一个在入边界和出边界上均为恒等映射的微分同胚，就将它们等同起来。此时，线性范畴就是复向量空间与线性映射的范畴；我们对函子只要求一条性质：无论对对象还是态射，它都把不交并映为张量积。闭 $d$ 维流形可以看作从空 $(d-1)$ 维流形到自身的配边，而张量积公理意味着，赋给空 $(d-1)$ 维流形的向量空间就是复数域。因此，理论为每个闭 $d$ 维流形指定一个 $1\times1$ 矩阵，也就是一个复数；这就是理论的配分函数。当然，$d=2$ 的情形尤其为人熟知，并且已经得到了很好的理解。

\SourceParagraph{ch02-2.1.1-002}{39}
推广几何范畴有几种自然方式。例如，可以考虑带某种额外结构的流形，如 Riemann 度量或自旋结构。引入度量，很快就会使我们面对量子场论的全部复杂性，所以这一点留到\BookChapterRef{3}。不过，通过加入自旋结构，可以先作部分推广，我们将在\BookSectionRef{sec:2.1.6}{2.1.6}中这样做。

\SourceParagraph{ch02-2.1.1-003}{40}
本章主要关注另一种推广：配分函数不仅为闭流形定义，也为带边界的定向 $d$ 维流形定义，而且边界的每个连通分支都用某个固定集合 $\mathcal{B}_0$ 中的元素标记，$\mathcal{B}_0$ 称为边界条件集合。将其形式化时，取几何范畴的对象为带边界的定向 $(d-1)$ 维流形，并用 $\mathcal{B}_0$ 的元素标记每个边界分支。此时，从 $Y_0$ 到 $Y_1$ 的配边，是一个边界由三部分组成的 $d$ 维流形 $X$：
\[
\partial X=Y_0\cup Y_1\cup\partial_{\mathrm{cstr}}X,
\]
其中“受约束边界”$\partial_{\mathrm{cstr}}X$ 是从 $\partial Y_0$ 到 $\partial Y_1$ 的配边。此外，还要求用 $\mathcal{B}_0$ 的元素标记 $\partial_{\mathrm{cstr}}X$ 的各个连通分支，并使这些标记与 $\partial Y_0$ 和 $\partial Y_1$ 的标记一致。

\SourceParagraph{ch02-2.1.1-004}{40}
因此，在 $d=2$ 时，几何范畴的对象是圆周和端点带标记的定向区间的不交并。从这个范畴到复向量空间范畴、且将不交并映为张量积的函子，称为开闭弦拓扑场论。这类理论将为我们提供 D 膜理论的一个“初级”模型。始终用 $\mathcal{C}$ 表示赋给标准圆周 $S^1$ 的向量空间，用 $\mathcal{O}_{ab}$ 表示赋给端点用 $a,b\in\mathcal{B}_0$ 标记、从 $b$ 指向 $a$ 的区间的向量空间；也就是说，它是从点 $b$ 到点 $a$ 的配边，而\textbf{不是}相反方向。

\begin{figure}[H]
\centering
\BookImage{01_Numbered_Figures/Chapter_02/Ch02_Fig01_PDFp040_Bookp030.pdf}
\caption{开弦的基本配边。在本图及后续所有开弦图示中，都用点划线表示受约束边界。}\label{fig:2.1}
\end{figure}

\SourceParagraph{ch02-2.1.1-005}{40}
\BookRef{fig:2.1}{图1}中的配边给出一个线性映射 $\mathcal{O}_{ab}\otimes\mathcal{O}_{bc}\to\mathcal{O}_{ac}$，或者等价地，给出双线性映射
\begin{equation}\label{eq:2.1}
\mathcal{O}_{ab}\times\mathcal{O}_{bc}\longrightarrow\mathcal{O}_{ac},
\end{equation}
我们把它看作复合律。事实上，由此得到一个 $\mathbb{C}$ 线性范畴 $\mathcal{B}$，其对象是 $\mathcal{B}_0$ 的元素，从 $b$ 到 $a$ 的态射集合是向量空间 $\mathcal{O}_{ab}$，态射的复合由\BookEqRef{eq:2.1}{2.1}给出。说 $\mathcal{B}$ 是 $\mathbb{C}$ 线性范畴，无非就是说，双线性复合\BookEqRef{eq:2.1}{2.1}在显然的意义下满足结合律，而且对每个 $a\in\mathcal{B}_0$ 都有一个恒等元素 $1_a\in\mathcal{O}_{aa}$；稍后将解释这些性质为什么成立。

\SourceParagraph{ch02-2.1.1-006}{41}
对于任意开闭弦 TFT，由圆柱配边 $S^1\times[0,1]$ 定义一个映射 $e\colon\mathcal{C}\to\mathcal{C}$，由正方形 $[0,1]\times[0,1]$ 定义一个映射 $e_{ab}\colon\mathcal{O}_{ab}\to\mathcal{O}_{ab}$。显然，$e^2=e$ 且 $e_{ab}^2=e_{ab}$。如果这些映射全是恒等映射，就称理论是约化的。将讨论限制在约化理论上，并无损失；从现在起，我们都这样做。

\subsection{代数刻画}\label{sec:2.1.2}

\SourceParagraph{ch02-2.1.2-001}{41}
按上一小节的方式表述的最一般的二维开闭弦 TFT，由以下代数数据给出。

\SourceParagraph{ch02-2.1.2-002}{41}
\begin{enumerate}[label={},leftmargin=2.8em]
\item[1.] $(\mathcal{C},\theta_{\mathcal{C}},1_{\mathcal{C}})$ 是交换 Frobenius 代数。
\item[2a.] 对 $a,b\in\mathcal{B}_0$，$\mathcal{O}_{ab}$ 是一族向量空间，带有结合的双线性乘法
\begin{equation}\label{eq:2.2}
\mathcal{O}_{ab}\otimes\mathcal{O}_{bc}\longrightarrow\mathcal{O}_{ac}.
\end{equation}
\item[2b.] 各 $\mathcal{O}_{aa}$ 都带有非退化迹
\begin{equation}\label{eq:2.3}
\theta_a\colon\mathcal{O}_{aa}\longrightarrow\mathbb{C}.
\end{equation}
特别地，每个 $\mathcal{O}_{aa}$ 都是 Frobenius 代数，但未必交换。
\item[2c.] 此外，
\begin{equation}\label{eq:2.4}
\mathcal{O}_{ab}\otimes\mathcal{O}_{ba}\longrightarrow\mathcal{O}_{aa}\xrightarrow{\theta_a}\mathbb{C}
\end{equation}
是完美配对，而且对于 $\psi_1\in\mathcal{O}_{ab}$、$\psi_2\in\mathcal{O}_{ba}$，有
\begin{equation}\label{eq:2.5}
\theta_a(\psi_1\psi_2)=\theta_b(\psi_2\psi_1).
\end{equation}
\item[3.] 存在线性映射
\begin{equation}\label{eq:2.6}
\begin{aligned}
\iota_a&\colon\mathcal{C}\longrightarrow\mathcal{O}_{aa},\\
\iota^a&\colon\mathcal{O}_{aa}\longrightarrow\mathcal{C},
\end{aligned}
\end{equation}
满足以下条件。
\item[3a.]\phantomsection\label{cond:2.1.2.3a} $\iota_a$ 是代数同态：
\begin{equation}\label{eq:2.7}
\iota_a(\phi_1\phi_2)=\iota_a(\phi_1)\iota_a(\phi_2).
\end{equation}
\item[3b.]\phantomsection\label{cond:2.1.2.3b} 保持恒等元素：
\begin{equation}\label{eq:2.8}
\iota_a(1_{\mathcal{C}})=1_a.
\end{equation}
\item[3c.]\phantomsection\label{cond:2.1.2.3c} 此外，$\iota_a$ 是中心的，意即对所有 $\phi\in\mathcal{C}$ 和 $\psi\in\mathcal{O}_{ab}$，都有
\begin{equation}\label{eq:2.9}
\iota_a(\phi)\psi=\psi\iota_b(\phi).
\end{equation}
\item[3d.]\phantomsection\label{cond:2.1.2.3d} $\iota_a$ 与 $\iota^a$ 互为伴随：对所有 $\psi\in\mathcal{O}_{aa}$，有
\begin{equation}\label{eq:2.10}
\theta_{\mathcal{C}}\bigl(\iota^a(\psi)\phi\bigr)=\theta_a\bigl(\psi\iota_a(\phi)\bigr).
\end{equation}
\item[3e.]\phantomsection\label{cond:2.1.2.3e} “Cardy 条件”\footnote{实际上，这些条件是 Cardy 所陈述条件的推广。对它们取迹，就恢复 Cardy 的条件。当然，在闭弦通道中将双扭转图分解这一观察，可以追溯到弦理论发展的最初时期。}。如下定义 $\pi_b^a\colon\mathcal{O}_{aa}\to\mathcal{O}_{bb}$。由于 $\mathcal{O}_{ab}$ 和 $\mathcal{O}_{ba}$ 通过 $\theta_a$ 或 $\theta_b$ 互相对偶，若取 $\psi_\mu$ 为 $\mathcal{O}_{ba}$ 的一组基，就存在 $\mathcal{O}_{ab}$ 的一组对偶基 $\psi^\mu$。定义
\begin{equation}\label{eq:2.11}
\pi_b^a(\psi)=\sum_\mu\psi_\mu\psi\psi^\mu,
\end{equation}
于是有“Cardy 条件”
\begin{equation}\label{eq:2.12}
\pi_b^a=\iota_b\circ\iota^a.
\end{equation}
\end{enumerate}

\begin{figure}[H]
\centering
\BookImage{01_Numbered_Figures/Chapter_02/Ch02_Fig02_PDFp042_Bookp032.pdf}
\caption{定义二维闭弦 TFT 的 Frobenius 结构的四个图示。用第二行的平面图示表示这些态射，往往更方便。此时采用如下约定：若圆周的定向使曲面位于它的右侧，则这个圆周是入圆周。}\label{fig:2.2}
\end{figure}

\subsection{图示表示}\label{sec:2.1.3}

\SourceParagraph{ch02-2.1.3-001}{42}
现在解释这些代数条件的图示基础。二维闭弦 TFT 的情形已经非常熟知。Frobenius 结构的数据由\BookRef{fig:2.2}{图2}中的图示给出，一致性条件则来自\BookRef{fig:2.3}{图3}。

\SourceParagraph{ch02-2.1.3-002}{42}
在开弦情形下，完全类似的考察使我们在开弦部分构造出一个未必交换的 Frobenius 代数。基本数据概括于\BookRef{fig:2.4}{图4}。\BookEqRef{eq:2.4}{2.4}中的配对是对偶配对，这一点来自\BookRef{fig:2.5}{图5}。开闭弦理论中最重要的新因素，是开弦到闭弦及闭弦到开弦的转变。在二维 TFT 中，它们就是映射 $\iota^a$、$\iota_a$，由\BookRef{fig:2.6}{图6}表示。与开闭弦转变关联的新的相容性条件共有五条，它们展示在\BookRef{fig:2.7}{图7}至\BookRef{fig:2.12}{图12}中。检验这些图示的拓扑断言时，通常最好将曲面想象为被“摊平”的样子。例如，\BookRef{fig:2.9}{图9}中的两个曲面都是圆环，一个边界圆周是入闭圆周，而另一个边界圆周分成一个出区间和一个受约束边界区间。同样，Cardy 图示\BookRef{fig:2.12}{图12}中的两个曲面也是圆环，而\BookRef{fig:2.7}{图7}中的每个曲面都是带两个洞的圆盘，也就是从一个圆盘中挖去两个开子圆盘所得的曲面。

\begin{figure}[H]
\centering
\BookImage{01_Numbered_Figures/Chapter_02/Ch02_Fig03_PDFp043_Bookp033.pdf}
\caption{闭弦情形下的结合律、交换律和单位元约束。单位元约束要求一个自然的假设：圆柱对应于恒等映射 $\mathcal{C}\to\mathcal{C}$。}\label{fig:2.3}
\end{figure}

\subsection{缝合定理}\label{sec:2.1.4}

\SourceParagraph{ch02-2.1.4-001}{43}
从几何上说，任意定向曲面都可以分解为态射的复合，这些态射对应于定义 Frobenius 结构的基本数据。不过，一个给定曲面可以用许多不同方式分解。上面的缝合公理来自这些分解的一致性。缝合定理保证，缝合的一致性不会对代数数据施加任何进一步的关系。

\begin{figure}[H]
\centering
\BookImage{01_Numbered_Figures/Chapter_02/Ch02_Fig04_PDFp044_Bookp034.pdf}
\caption{开弦理论的基本数据。受约束边界以点划线表示，并带有边界条件 $a,b,c,\ldots\in\mathcal{B}_0$。}\label{fig:2.4}
\end{figure}

\SourceParagraph{ch02-thm-2.1}{44}
\BookStatementNumber{1}
\begin{theorem}\label{thm:2.1}
上面的条件1、2、3足以保证，这些代数数据给出一个良定义的开闭弦拓扑场论。
\end{theorem}

\SourceParagraph{ch02-2.1.4-002}{44}
证明见\BookSectionRef{sec:2.7}{2.7}。

\subsection{边界条件范畴}\label{sec:2.1.5}

\SourceParagraph{ch02-2.1.5-001}{44}
开闭弦 TFT 的边界条件范畴 $\mathcal{B}$ 是 $\mathbb{C}$ 线性范畴。我们可以用各种方式向其中加入新对象。

\begin{figure}[H]
\centering
\BookImage{01_Numbered_Figures/Chapter_02/Ch02_Fig05_PDFp045_Bookp035.pdf}
\caption{假设条带对应于恒等态射，则\BookEqRef{eq:2.4}{2.4}中的配对必然是完美配对。}\label{fig:2.5}
\end{figure}

\begin{figure}[H]
\centering
\BookImage{01_Numbered_Figures/Chapter_02/Ch02_Fig06_PDFp045_Bookp035.pdf}
\caption{表示开弦到闭弦及闭弦到开弦转变的两种方式。}\label{fig:2.6}
\end{figure}

\SourceParagraph{ch02-2.1.5-002}{45}
例如，如果范畴没有直和，那么对于任意两个对象 $a$ 和 $b$，可以定义一个新对象 $a\oplus b$，使
\begin{equation}\label{eq:2.13}
\mathcal{O}_{a\oplus b,c}\coloneqq\mathcal{O}_{ac}\oplus\mathcal{O}_{bc},
\end{equation}
\begin{equation}\label{eq:2.14}
\mathcal{O}_{c,a\oplus b}\coloneqq\mathcal{O}_{ca}\oplus\mathcal{O}_{cb},
\end{equation}
于是\footnote{这里的矩阵记号是为了帮助理解态射的复合。作为向量空间，$\mathcal{O}_{a\oplus b,a\oplus b}$ 就是四个空间 $\mathcal{O}_{aa}$、$\mathcal{O}_{ab}$、$\mathcal{O}_{ba}$、$\mathcal{O}_{bb}$ 的直和。}
\begin{equation}\label{eq:2.15}
\mathcal{O}_{a\oplus b,a\oplus b}\coloneqq
\begin{pmatrix}\mathcal{O}_{aa}&\mathcal{O}_{ab}\\\mathcal{O}_{ba}&\mathcal{O}_{bb}\end{pmatrix},
\end{equation}
配备显然的复合律，并定义
\begin{equation}\label{eq:2.16}
\theta_{a\oplus b}\colon\mathcal{O}_{a\oplus b,a\oplus b}\longrightarrow\mathbb{C},
\end{equation}
其表达式为
\begin{equation}\label{eq:2.17}
\theta_{a\oplus b}\begin{pmatrix}\psi_{aa}&\psi_{ab}\\\psi_{ba}&\psi_{bb}\end{pmatrix}
=\theta_a(\psi_{aa})+\theta_b(\psi_{bb}).
\end{equation}

\begin{figure}[H]
\centering
\BookImage{01_Numbered_Figures/Chapter_02/Ch02_Fig07_PDFp046_Bookp036.pdf}
\caption{$\iota_a$ 是同态。}\label{fig:2.7}
\end{figure}

\begin{figure}[H]
\centering
\BookImage{01_Numbered_Figures/Chapter_02/Ch02_Fig08_PDFp046_Bookp036.pdf}
\caption{$\iota_a$ 保持恒等元素。}\label{fig:2.8}
\end{figure}

\SourceParagraph{ch02-2.1.5-003}{46}
这个新对象，就是扩大的边界条件范畴中 $a$ 与 $b$ 的直和。如果范畴 $\mathcal{B}$ 中已经存在 $a$ 与 $b$ 的直和，那么新对象将与它典范同构。反过来，如果有一个边界条件 $a$ 和一个投影 $p\in\mathcal{O}_{aa}$，即满足 $p^2=p$ 的元素，那么也可以假设存在边界条件 $b=\operatorname{im}(p)$，使得对任意 $c$ 都有
\[
\mathcal{O}_{cb}=\{f\in\mathcal{O}_{ca}\mid fp=f\},\qquad
\mathcal{O}_{bc}=\{f\in\mathcal{O}_{ac}\mid pf=f\}.
\]
于是便有 $a\cong\operatorname{im}(p)\oplus\operatorname{im}(1-p)$\footnote{幂等元以这种方式分裂的线性范畴，常称为 Karoubi 范畴。另见\BookSectionRef{sec:8.3.4}{8.3.4}末尾的简短相关讨论。}。

\begin{figure}[H]
\centering
\BookImage{01_Numbered_Figures/Chapter_02/Ch02_Fig09_PDFp047_Bookp037.pdf}
\caption{$\iota_a$ 映入 $\mathcal{O}_{aa}$ 的中心。}\label{fig:2.9}
\end{figure}

\begin{figure}[H]
\centering
\BookImage{01_Numbered_Figures/Chapter_02/Ch02_Fig10_PDFp047_Bookp037.pdf}
\caption{$\iota^a$ 是 $\iota_a$ 的伴随。}\label{fig:2.10}
\end{figure}

\begin{figure}[H]
\centering
\BookImage{01_Numbered_Figures/Chapter_02/Ch02_Fig11_PDFp047_Bookp037.pdf}
\caption{双扭转图定义映射 $\pi_b^a\colon\mathcal{O}_{aa}\to\mathcal{O}_{bb}$。}\label{fig:2.11}
\end{figure}

\begin{figure}[H]
\centering
\BookImage{01_Numbered_Figures/Chapter_02/Ch02_Fig12_PDFp048_Bookp038.pdf}
\caption{广义 Cardy 条件，表示双扭转图在闭弦通道中的分解。}\label{fig:2.12}
\end{figure}

\SourceParagraph{ch02-2.1.5-004}{48}
范畴 $\mathcal{B}$ 具有一项非常特殊的性质：对于任意两个对象 $a$ 和 $b$，态射空间 $\mathcal{O}_{ab}$ 都通过一个配对与 $\mathcal{O}_{ba}$ 典范对偶，而且这个配对可以通过任一种次序的复合来分解。将具有这一性质的范畴称为 Frobenius 范畴，或者 Calabi--Yau 范畴\footnote{后一种名称来自紧 Kähler 流形上的凝聚层情形：对于两个层 $E$ 和 $F$，态射空间 $\operatorname{Ext}(E,F)$ 的对偶，一般是 $\operatorname{Ext}(F,E\otimes\omega)$，其中 $\omega$ 为典范丛。只有当 $\omega$ 平凡，也就是在 Calabi--Yau 情形下，它才与 $\operatorname{Ext}(F,E)$ 相同。细节见\BookSectionRef{sec:4.3}{4.3}。我们将在\BookSectionRef{sec:2.5}{2.5}中进一步讨论这个例子。}，是很自然的。这是对范畴的一个很强的限制：例如，有限维代数上的有限生成模范畴就没有这一性质，除非这个代数半单。

\SourceParagraph{ch02-example-2.2-001}{48}
\BookStatementNumber{2}
\begin{example}\label{ex:2.2}
我们所研究的这类开闭弦理论中，也许最简单的例子，是与有限群 $G$ 关联的理论。范畴 $\mathcal{B}$ 是 $G$ 的有限维复表示 $M$ 所构成的范畴，而迹 $\theta_M\colon\mathcal{O}_{MM}=\operatorname{End}(M)\to\mathbb{C}$ 将 $\psi\colon M\to M$ 映为 $\operatorname{trace}(\psi)/|G|$。闭弦代数 $\mathcal{C}$ 是群代数 $\mathbb{C}[G]$ 的中心，并以显然的方式映入每个 $\operatorname{End}(M)$。迹 $\theta_{\mathcal{C}}\colon\mathcal{C}\to\mathbb{C}$ 将群代数中的中心元素 $\sum\lambda_g g$ 映为 $\lambda_1/|G|$。

\SourceParagraph{ch02-example-2.2-002}{48}
在这个例子中，若曲面 $\Sigma$ 的受约束边界圆周 $C_1,\ldots,C_k$ 分别以 $M_1,\ldots,M_k$ 标记，那么理论的配分函数就是对 $\Sigma$ 上主 $G$ 丛 $P$ 的同构类，对下式作加权求和：
\[
\chi_{M_1}\bigl(h_P(C_1)\bigr)\cdots\chi_{M_k}\bigl(h_P(C_k)\bigr),
\]
其中 $\chi_M\colon G\to\mathbb{C}$ 是表示 $M$ 的特征标，$h_P(C)$ 表示 $P$ 绕边界圆周 $C$ 一周的平行移动元素。每个丛 $P$ 的权重，是其自同构群阶数的倒数。
\end{example}

\SourceParagraph{ch02-2.1.5-005}{49}
回到一般理论，现在可以提出三个基本问题。

\SourceParagraph{ch02-2.1.5-question-i}{49}
\begin{problem*}[(i)]
给定一个“闭弦”TFT，能否将它扩展为开闭弦理论？如果可以，这种扩展是否唯一？
\end{problem*}

\SourceParagraph{ch02-2.1.5-question-ii}{49}
\begin{problem*}[(ii)]
给定一个开闭弦理论的边界条件范畴 $\mathcal{B}$，以及定义 Frobenius 结构的线性映射 $\theta_a\colon\mathcal{O}_{aa}\to\mathbb{C}$，能否重建整个理论？也就是说，能否找出闭弦 Frobenius 代数 $\mathcal{C}$？
\end{problem*}

\SourceParagraph{ch02-2.1.5-question-iii}{49}
\begin{problem*}[(iii)]
任意 Frobenius 范畴，是否都是某个闭弦理论的边界条件范畴？
\end{problem*}

\SourceParagraph{ch02-2.1.5-006}{49}
要使第一个问题良定，应当假设边界条件范畴是极大的：如果 $\mathcal{B}'$ 是它的一个扩张，那么 $\mathcal{B}'$ 中的任意对象都同构于 $\mathcal{B}$ 中的某个对象。即使如此，我们也将看到，某些微妙之处仍然使这些问题都没有简单的肯定答案。

\subsection{推广：自旋理论}\label{sec:2.1.6}

\SourceParagraph{ch02-2.1.6-001}{49}
改变几何范畴或线性范畴，都可以使上述结构得到许多有趣的推广。现在将概览几类例子，以作说明；有时只概述其细节。

\SourceParagraph{ch02-2.1.6-002}{49}
可以考虑的最一般的靶范畴，是对称张量范畴。显然，我们需要张量积；而公理 $\mathcal{H}_{Y_1\sqcup Y_2}\cong\mathcal{H}_{Y_1}\otimes\mathcal{H}_{Y_2}$，只有在存在一个对合性的典范同构 $\mathcal{H}_{Y_1}\otimes\mathcal{H}_{Y_2}\cong\mathcal{H}_{Y_2}\otimes\mathcal{H}_{Y_1}$ 时才有意义。

\SourceParagraph{ch02-2.1.6-003}{49}
物理中一个很常见的选择，是超向量空间范畴，即带有模2分次 $V=V^0\oplus V^1$ 的向量空间 $V$；其中的典范同构 $V\otimes W\cong W\otimes V$ 为
\[
v\otimes w\longmapsto(-1)^{\deg v\deg w}w\otimes v.
\]
也可以考虑 $\mathbb{Z}$ 分次向量空间范畴，并对张量积采用相同的符号约定。

\SourceParagraph{ch02-2.1.6-004}{49}
在这两种情形中，闭弦代数都是带迹 $\theta\colon\mathcal{C}\to\mathbb{C}$ 的分次交换代数 $\mathcal{C}$。原则上，迹的次数应当为零；不过，实际常见的理论带有分次反常，使迹具有次数 $-n$，其中 $n$ 是某个整数\footnote{不难看出，容许对分次作整体平移后，最一般的反常会为 Euler 示性数为 $\chi$、具有 $i$ 个入边界圆周和 $o$ 个出边界圆周的配边，指定一个次数为 $\frac12n(i-o-\chi)$ 的算符。}。公式\BookEqRef{eq:2.5}{2.5}、\BookEqRef{eq:2.9}{2.9}和\BookEqRef{eq:2.11}{2.11}必须替换为各自的分次交换版本。特别地，如果选取齐次基 $\psi_\mu$ 及其对偶基 $\psi^\mu$，使
\begin{equation}\label{eq:2.18}
\theta_{\mathcal{C}}(\psi^\mu\psi_\nu)=\delta^\mu_\nu,
\end{equation}
那么
\begin{equation}\label{eq:2.19}
\pi_b^a(\psi)=\sum_\mu(-1)^{\deg\psi_\mu\deg\psi}\psi_\mu\psi\psi^\mu.
\end{equation}

\SourceParagraph{ch02-2.1.6-005}{50}
通过为流形及配边添加额外结构，改变它们所构成的几何范畴，也可以得到有趣的结构。

\SourceParagraph{ch02-example-2.3-001}{50}
\BookStatementNumber{3}
\begin{example}\label{ex:2.3}
将“流形”替换为“带自旋结构的流形”，就定义了拓扑自旋理论。

\SourceParagraph{ch02-example-2.3-002}{50}
曲面上的自旋结构，是其非零切向量空间的一个二重覆盖，并且在每个单独的切空间上都非平凡。在定向一维流形 $S$ 上，它是正向切向量空间的一个二重覆盖。为进行缝合，需要注意，这等同于可定向曲面中 $S$ 的一个带状邻域上的自旋结构。每个自旋结构都有一个交换两张覆盖片的自同构；它会在任何自然地赋给带自旋结构的一维流形的向量空间上，诱导一个对合 $T$，通过其 $\pm1$ 本征空间，为该向量空间给出模2分次。我们把拓扑自旋理论定义为一个函子，从带自旋结构的流形的配边范畴，映到带分次张量结构的超向量空间范畴。要求这个函子把不交并映为超张量积，还要求一维流形自旋结构的自同构，诱导超向量空间的分次自同构 $T=(-1)^{\mathrm{degree}}$。稍后将看到，如果希望允许半单边界条件范畴的存在，那么自旋结构的几何会迫使我们选取张量积的这种超对称性，而不能选取忽略分次的朴素对称性。带自旋结构的圆周有两种不同构的类型：带 Möbius 或“Neveu--Schwarz”结构的 $S^1_{\mathrm{ns}}$，以及带平凡或“Ramond”结构的 $S^1_{\mathrm{r}}$。拓扑自旋理论分别为它们赋予态空间 $\mathcal{C}_{\mathrm{ns}}$ 和 $\mathcal{C}_{\mathrm{r}}$。

\SourceParagraph{ch02-example-2.3-003}{50}
覆盖标准圆环的带自旋结构的配边共有四种。二重覆盖可以与其入端和区间 $[0,1]$ 的乘积等同，但随后将圆环上的二重覆盖的出端，与出边界圆周上选定的结构等同时，仍有两种选择。换言之，除了诱导 $\mathcal{C}_{\mathrm{ns},\mathrm{r}}$ 上恒等映射的圆柱 $A^+_{\mathrm{ns},\mathrm{r}}=S^1_{\mathrm{ns},\mathrm{r}}\times[0,1]$，还有圆柱 $A^-_{\mathrm{ns},\mathrm{r}}$，它们把 $S^1_{\mathrm{ns},\mathrm{r}}$ 连到自身，同时交换覆盖片。这些圆柱 $A^-_{\mathrm{ns},\mathrm{r}}$ 在态空间上诱导分次自同构。不过，通过一个在边界圆周上为恒等映射的同构，即将“圆柱的一端旋转 $2\pi$”的 Dehn 扭转，有 $A^-_{\mathrm{ns}}\cong A^+_{\mathrm{ns}}$，因此 $\mathcal{C}_{\mathrm{ns}}$ 的分次必须纯为偶数。空间 $\mathcal{C}_{\mathrm{r}}$ 则可以同时具有偶和奇两部分。

\SourceParagraph{ch02-example-2.3-004}{51}
对于“U 形”配边，也就是带两个入边界圆周或两个出边界圆周的圆柱，情形稍复杂。如果边界是 $S^1_{\mathrm{ns}}$，就只有一种可能；如果边界是 $S^1_{\mathrm{r}}$，则有两种，分别对应于 $A^\pm_{\mathrm{r}}$。复杂之处在于，似乎没有特殊理由把某一种自旋结构选为“正”。我们干脆选定一种，记为 $P$，其入边界为 $S^1_{\mathrm{r}}\sqcup S^1_{\mathrm{r}}$，并用它定义配对 $\mathcal{C}_{\mathrm{r}}\otimes\mathcal{C}_{\mathrm{r}}\to\mathbb{C}$。随后，在相反方向选取一个优先的配边 $Q$，使得把 $Q$ 右侧的出 $S^1_{\mathrm{r}}$ 与 $P$ 左侧的入 $S^1_{\mathrm{r}}$ 缝合后，得到的 S 形弯折是“平凡”圆柱 $A^+_{\mathrm{r}}$。不过，需要知道，复合 $P\circ Q$ 所形成的闭环面，具有偶自旋结构。注意，$\mathcal{C}$ 上的 Frobenius 结构 $\theta$ 限制到 $\mathcal{C}_{\mathrm{r}}$ 上为零。

\SourceParagraph{ch02-example-2.3-005}{51}
\BookRef{fig:2.2}{图2}的裤形配边上，有唯一的自旋结构在每个边界圆周上限制为 $S^1_{\mathrm{ns}}$；它以通常的方式使 $\mathcal{C}_{\mathrm{ns}}$ 成为交换 Frobenius 代数。如果一个入圆周是 $S^1_{\mathrm{ns}}$，另一个是 $S^1_{\mathrm{r}}$，那么出圆周就是 $S^1_{\mathrm{r}}$，此时有两种可能的自旋结构；其中，从圆柱 $A^+_{\mathrm{r}}$ 挖去一个圆盘所得的那一种是优先选取的，它使 $\mathcal{C}_{\mathrm{r}}$ 成为 $\mathcal{C}_{\mathrm{ns}}$ 上的分次模。对具有两个入圆周 $S^1_{\mathrm{r}}$ 的已选 U 形配边 $P$，挖去一个圆盘，就得到一个出圆周为 $S^1_{\mathrm{ns}}$ 的裤形曲面。它诱导一个分次双线性映射 $\mathcal{C}_{\mathrm{r}}\times\mathcal{C}_{\mathrm{r}}\to\mathcal{C}_{\mathrm{ns}}$；将它与 $\mathcal{C}_{\mathrm{ns}}$ 上的迹复合，就得到 $\mathcal{C}_{\mathrm{r}}$ 上的非退化内积。此时，对张量积对称性的选择变得重要。考虑裤形曲面的那个微分同胚，在通常情形下，正是它说明 Frobenius 代数满足交换律。将它提升到自旋结构时，它在一个入圆周上诱导恒等映射，却交换另一个入圆周上的覆盖片；这证明，将输入从 $S(\phi_1\otimes\phi_2)$ 改为 $T(\phi_1)\otimes\phi_2$ 时，配边必须给出相同的输出，其中 $T$ 是分次对合，$S\colon\mathcal{C}_{\mathrm{r}}\otimes\mathcal{C}_{\mathrm{r}}\to\mathcal{C}_{\mathrm{r}}\otimes\mathcal{C}_{\mathrm{r}}$ 是张量范畴的对称性。如果取 $S$ 为忽略分次的向量空间张量范畴的对称性，那么这说明，分次向量空间 $\mathcal{C}_{\mathrm{r}}$ 上的乘法按通常的符号规则分次对称；如果 $S$ 是分次对称性，则可以看出，$\mathcal{C}_{\mathrm{r}}$ 上的乘法在朴素意义下对称。这里必须记住，若 $\psi_1$ 与 $\psi_2$ 的奇偶性不同，那么它们的乘积无论如何都为零，因为已经看到 $\mathcal{C}_{\mathrm{ns}}$ 纯为偶数。

\SourceParagraph{ch02-example-2.3-006}{51}
说明二维拓扑场论“就是”交换 Frobenius 代数的那个定理，在自旋理论中也有类似版本。它断言，自旋拓扑理论“就是”一个 Frobenius 代数 $\mathcal{C}=(\mathcal{C}_{\mathrm{ns}}\oplus\mathcal{C}_{\mathrm{r}},\theta_{\mathcal{C}})$，具有刚刚提到的性质，以及如下额外性质。取 $\{\phi_k\}$ 为 $\mathcal{C}_{\mathrm{ns}}$ 的一组基，$\{\phi^k\}$ 为其对偶基，使 $\theta_{\mathcal{C}}(\phi^k\phi_m)=\delta^k_m$，并取 $\beta_k$ 和 $\beta^k$ 为 $\mathcal{C}_{\mathrm{r}}$ 中类似的对偶基。那么，Euler 元素
\[
\chi_{\mathrm{ns}}\coloneqq\sum_k\phi_k\phi^k,\qquad
\chi_{\mathrm{r}}\coloneqq\sum_k\beta_k\beta^k
\]
不依赖于基的选择；我们对代数 $\mathcal{C}$ 所要求的条件，是 $\chi_{\mathrm{ns}}=\chi_{\mathrm{r}}$。特别地，这一条件意味着，向量空间 $\mathcal{C}_{\mathrm{ns}}$ 与 $\mathcal{C}_{\mathrm{r}}$ 的维数相同\footnote{因此，在某种意义上，理论具有“时空超对称”。}。实际上，Euler 元素可以通过从环面中挖去一个洞来得到。环面上共有四种自旋结构，输出态必定位于 $\mathcal{C}_{\mathrm{ns}}$ 中。三个偶自旋结构的 Euler 元素相同，均为 $\chi_{\mathrm{e}}=\chi_{\mathrm{ns}}=\chi_{\mathrm{r}}$；另一方面，奇自旋结构对应的 Euler 元素为
\[
\chi_{\mathrm{o}}=\sum_k(-1)^{\deg\beta_k}\beta_k\beta^k.
\]

\SourceParagraph{ch02-example-2.3-007}{52}
我们略去一般自旋理论恰好是上述结构的证明；不过，它与\BookSectionRef{sec:2.7}{2.7}中将给出的 Turaev 关于 $G$ 等变理论的定理的证明，几乎完全相同，这里只需取群 $G=\mathbb{Z}/2$ 这一简单情形。事实上，自旋理论与 $\mathbb{Z}/2$ 等变理论非常相似，但并不相同；后者是让曲面带主 $\mathbb{Z}/2$ 丛，即二重覆盖，而非自旋结构时得到的结构。\BookSectionRef{sec:2.6}{2.6}将讨论等变理论。两者的一个区别是：在等变情形下，$\mathbb{Z}/2$ 作用在 $\mathcal{C}_1$ 部分非平凡，在 $\mathcal{C}_g$ 部分平凡，恰与自旋情形中的发现相反。与等变理论比较，就可以把 $\mathcal{C}_{\mathrm{r}}$ 上乘法朴素对称这一令人惊讶的结果，理解为扭曲的反交换性。

\SourceParagraph{ch02-example-2.3-008}{52}
如果代数 $\mathcal{C}_{\mathrm{ns}}$ 半单，即是某个有限集合 $X$ 上的函数代数，那么把这个自旋理论称为半单的，似乎是合理的。此时，$\mathcal{C}_{\mathrm{r}}$ 是 $X$ 上某个向量丛 $E$ 的截面空间，而条件 $\chi_{\mathrm{ns}}=\chi_{\mathrm{r}}$ 意味着每一点处的纤维都必须是一维的。因此，整个结构由 Frobenius 代数 $\mathcal{C}_{\mathrm{ns}}$，以及在每个点 $x\in X$ 处对线丛 $E$ 的纤维 $E_x$ 的分次所作的二选一选择确定。

\SourceParagraph{ch02-example-2.3-009}{52}
现在可以看到，如果在定义张量范畴时没有使用分次对称性，就会迫使 $\mathcal{C}_{\mathrm{r}}$ 的分次纯为偶数。因为在奇部分上，内积必须反对称，而这在一维空间上不可能。如果 $\mathcal{C}_{\mathrm{ns}}$ 和 $\mathcal{C}_{\mathrm{r}}$ 都纯为偶数，那么理论实际上完全不依赖于曲面上的自旋结构。

\SourceParagraph{ch02-example-2.3-010}{52}
二维拓扑自旋理论的一个具体例子，由 $\mathcal{C}=\mathbb{C}\oplus\mathbb{C}\eta$ 给出，其中 $\eta^2=1$，而 $\eta$ 为奇。Euler 元素为 $\chi_{\mathrm{e}}=1$ 和 $\chi_{\mathrm{o}}=-1$。因此，带自旋结构的闭曲面的配分函数，随自旋结构为偶或奇，分别为 $+1$ 或 $-1$。证明这一点时，可以计算与自旋结构关联的交叉乘积的二次精化的 Arf 不变量，并注意到它对添加柄具有乘法性。

\SourceParagraph{ch02-example-2.3-011}{53}
定义在带自旋结构曲面上的最常见的理论，并不是拓扑理论，而是具有 $N=1$ 超对称的二维共形场论。其结构的一般特征仍然如上所述，但应注意，如果理论不是拓扑的，就不应预期 $\mathcal{C}_{\mathrm{ns}}$ 的分次纯为偶数：态在旋转 $2\pi$ 后可以改变符号。如果曲面 $\Sigma$ 带共形结构，那么非零切向量的二重覆盖，就是 $\Sigma$ 上一个称为自旋丛的实二维向量丛 $L$ 中去掉零截面所得的空间。覆盖映射因而扩展为向量丛的对称配对 $L\otimes L\to T\Sigma$；若以自然方式将 $L$ 和 $T\Sigma$ 看作复线丛，它就诱导同构 $L\otimes_{\mathbb{C}}L\cong T\Sigma$。$N=1$ 超共形场论，是为带自旋丛 $L$ 的一维流形 $S$ 指定向量空间 $\mathcal{H}_{S,L}$ 的共形自旋理论，并配备额外映射
\begin{equation}\label{eq:2.20}
\Gamma(S,L)\otimes\mathcal{H}_{S,L}\longrightarrow\mathcal{H}_{S,L},
\end{equation}
\begin{equation}\label{eq:2.21}
(\sigma,\psi)\longmapsto G_\sigma\psi,
\end{equation}
其中 $\Gamma(S,L)$ 是 $L$ 的光滑截面空间，要求 $G_\sigma$ 对截面 $\sigma$ 为实线性的，并满足 $G_\sigma^2=D_{\sigma^2}$；这里，$D_{\sigma^2}$ 是向量场 $\sigma^2$ 的 Virasoro 作用，而 $\sigma^2$ 通过同构 $L\otimes_{\mathbb{C}}L\cong T\Sigma$ 与 $\sigma\otimes\sigma$ 相联系。此外，如果有从 $(S_0,L_0)$ 到 $(S_1,L_1)$ 的配边 $(\Sigma,L)$，并且 $L$ 的全纯截面 $\sigma$ 在 $S_i$ 上限制为 $\sigma_i$，那么有交织性质
\begin{equation}\label{eq:2.22}
G_{\sigma_1}\circ U_{\Sigma,L}=U_{\Sigma,L}\circ G_{\sigma_0}.
\end{equation}
\end{example}

\SourceParagraph{ch02-example-2.4-001}{53}
\BookStatementNumber{4}
\begin{example}\label{ex:2.4}
将“流形”替换为“带 $\operatorname{spin}^c$ 结构的流形”，就定义了拓扑 $\operatorname{spin}^c$ 理论，它们为具有 $N=2$ 超对称的二维理论提供模型。

\SourceParagraph{ch02-example-2.4-002}{53}
带共形结构的曲面上的 $\operatorname{spin}^c$ 结构，是一对全纯线丛 $L_1$、$L_2$，并配有全纯线丛同构 $L_1\otimes L_2\cong T\Sigma$。自旋结构就是 $L_1=L_2$ 的特殊情形。在一维流形 $S$ 上，$\operatorname{spin}^c$ 结构是指 $S$ 在带共形结构的曲面中的一个带状邻域上的 $\operatorname{spin}^c$ 结构。$N=2$ 超共形理论为每个带 $\operatorname{spin}^c$ 结构的一维流形 $S$ 指定向量空间 $\mathcal{H}_{S;L_1,L_2}$，并为每个从 $S_0$ 到 $S_1$ 的 $\operatorname{spin}^c$ 配边 $\Sigma$ 指定算符
\begin{equation}\label{eq:2.23}
U_{\Sigma;L_1,L_2}\colon\mathcal{H}_{S_0;L_1,L_2}\longrightarrow\mathcal{H}_{S_1;L_1,L_2}.
\end{equation}
为了说明其余结构，需要定义与 $\operatorname{spin}^c$ 一维流形 $(S;L_1,L_2)$ 关联的 $N=2$ Lie 超代数。令 $\mathcal{G}=\operatorname{Aut}(L_1)$，表示覆盖 $S$ 的微分同胚的丛同构 $L_1\to L_1$ 所构成的群，这个群也可以与 $\operatorname{Aut}(L_2)$ 等同。它具有到 $S$ 的保定向微分同胚群 $\operatorname{Diff}^+(S)$ 的满同态；其核是 $L_1$ 的逐纤维自同构群，可以与从 $S$ 到 $\mathbb{C}^\times$ 的光滑映射群等同。因此，Lie 代数 $\operatorname{Lie}(\mathcal{G})$ 是 $\operatorname{Diff}^+(S)$ 的 Lie 代数 $\operatorname{Vect}(S)$ 由 $S$ 上光滑实值函数构成的交换 Lie 代数 $\Omega^0(S)$ 的一个扩张。令 $\Lambda^0_{S;L_1,L_2}$ 表示将 $\operatorname{Lie}(\mathcal{G})$ 中的 $\operatorname{Vect}(S)$ 复化所得的复 Lie 代数。它是一个 Lie 超代数的偶部分，其奇部分为 $\Lambda^1_{S;L_1,L_2}=\Gamma(L_1)\oplus\Gamma(L_2)$。括号 $\Lambda^1\otimes\Lambda^1\to\Lambda^0$ 完全由以下性质确定：$\Gamma(L_1)$ 中的元素彼此反对易，$\Gamma(L_2)$ 中的元素也彼此反对易，而复合映射
\begin{equation}\label{eq:2.24}
\Gamma(L_1)\otimes\Gamma(L_2)\longrightarrow\Lambda^0\longrightarrow\operatorname{Vect}_{\mathbb{C}}(S)
\end{equation}
将 $(\lambda_1,\lambda_2)$ 映为 $\lambda_1\lambda_2\in\Gamma(TS)$。

\SourceParagraph{ch02-example-2.4-003}{54}
在 $N=2$ 理论中，我们要求超代数 $\Lambda(S;L_1,L_2)$ 作用于向量空间 $\mathcal{H}_{S;L_1,L_2}$，并与群 $\mathcal{G}$ 的作用相容，同时与配边算符具有类似于 $N=1$ 情形的交织性质。对于 $N=2$ 理论，态空间总带有圆周群的作用，它来自将圆周群嵌入 $\mathcal{G}$，作为 $L_1$ 与 $L_2$ 上互逆的逐纤维标量乘法。等价地，态空间总是 $\mathbb{Z}$ 分次的。

\SourceParagraph{ch02-example-2.4-004}{54}
为曲面 $\Sigma$ 分别配备 $\operatorname{spin}^c$ 结构 $(\mathbb{C},T\Sigma)$ 和 $(T\Sigma,\mathbb{C})$，每个 $N=2$ 理论都会给出两个通常的共形场论。它们称为与该 $N=2$ 理论关联的“A 模型”和“B 模型”。在这两种情形下，态空间都是上链复形，其微分是 $\operatorname{spin}^c$ 结构中平凡分量的常截面1的作用。
\end{example}

\subsection{推广：上链层面的理论}\label{sec:2.1.7}

\SourceParagraph{ch02-2.1.7-001}{54}
不过，对我们所描述的开闭弦拓扑场论，最重要的“推广”具有更为根本的性质。我们的拓扑理论，旨在成为弦理论中出现的共形场论的玩具模型。在闭弦理论中，核心对象是单条带参数的弦的态向量空间 $\mathcal{C}=\mathcal{C}_{S^1}$。它按“鬼数”作整数分次，并带有称为“BRST 算符”的算符 $Q\colon\mathcal{C}\to\mathcal{C}$，使鬼数增加1，并满足 $Q^2=0$。换言之，$\mathcal{C}$ 是上链复形。如果把弦看作在时空 $M$ 中运动，那么 $\mathcal{C}$ 大致是自由环路空间 $\mathcal{L}M$ 上，沿重参数化群 $\operatorname{Diff}^+(S^1)$ 的作用轨道定义的微分形式空间；更精确地说，是半无限次数的平方可积形式。同样，刚才描述的拓扑扭曲 $N=2$ 超对称理论的空间 $\mathcal{C}$，也是上链复形，为 Kähler 流形的环路空间上的半无限微分形式空间提供模型；此时包括全部平方可积微分形式，而不只是沿 $\operatorname{Diff}^+(S^1)$ 轨道的形式。在这两类例子中，从 $p$ 个圆周到 $q$ 个圆周的配边 $\Sigma$ 给出算符 $U_{\Sigma,\mu}\colon\mathcal{C}^{\otimes p}\to\mathcal{C}^{\otimes q}$，它依赖于 $\Sigma$ 上的共形结构 $\mu$。这个算符是上链映射，但其关键性质是：改变 $\Sigma$ 上的共形结构 $\mu$，只会使 $U_{\Sigma,\mu}$ 改变一个上链同伦。因此，上同调 $H(\mathcal{C})=\ker(Q)/\operatorname{im}(Q)$，也就是传统弦理论中的“物理态空间”，是一个拓扑场论的态空间。在通常的弦理论情形下，所得拓扑场论并不很有趣，因为 BRST 上同调集中在一个或两个次数中，而且存在“分次反常”，使与配边 $\Sigma$ 关联的算符，将次数改变 $\Sigma$ 的 Euler 示性数 $\chi(\Sigma)$ 的某个倍数。但在 $N=2$ 超对称模型中，没有分次反常，完整结构都能显现出来。

\SourceParagraph{ch02-2.1.7-002}{55}
可以用如下合适的方式，描述算符 $U_{\Sigma,\mu}$ 如何随 $\mu$ 变化。

\SourceParagraph{ch02-2.1.7-003}{55}
令 $\mathcal{M}_\Sigma$ 为配边 $\Sigma$ 上共形结构的模空间，其中模去在边界圆周上为恒等映射的 $\Sigma$ 的微分同胚。那么，有上链映射
\begin{equation}\label{eq:2.25}
U_\Sigma\colon\mathcal{C}^{\otimes p}\longrightarrow\Omega^*(\mathcal{M}_\Sigma,\mathcal{C}^{\otimes q}),
\end{equation}
其中右侧是 $\mathcal{M}_\Sigma$ 上取值于 $\mathcal{C}^{\otimes q}$ 的形式的 de Rham 复形。将 $U_\Sigma$ 从 $\mathcal{M}_\Sigma$ 限制到 $\{\mu\}$，就得到算符 $U_{\Sigma,\mu}$。两个配边 $\Sigma_1$ 和 $\Sigma_2$ 接合时，复合性质就是下图交换：
\begin{figure}[H]
\centering
\begin{equation}\label{eq:2.26}
\vcenter{\hbox{\BookImage{02_Unnumbered_Diagrams/Chapter_02/Ch02_Diagram_PDFp055_01_Bookp045.pdf}}}
\end{equation}
\end{figure}
其中，下方的水平箭头由映射 $\mathcal{M}_{\Sigma_1}\times\mathcal{M}_{\Sigma_2}\to\mathcal{M}_{\Sigma_2\circ\Sigma_1}$ 诱导，这个映射表示共形结构的接合。

\SourceParagraph{ch02-2.1.7-004}{55}
这一表述可以有许多变体。例如，也许更愿意给出上链映射
\begin{equation}\label{eq:2.27}
U_\Sigma\colon C_\bullet(\mathcal{M}_\Sigma)\longrightarrow(\mathcal{C}^{\otimes p})^*\otimes\mathcal{C}^{\otimes q},
\end{equation}
其中 $C_\bullet(\mathcal{M}_\Sigma)$ 例如可以取为 $\mathcal{M}_\Sigma$ 的光滑奇异链复形。也可以倾向于使用 Riemann 结构的模空间，而非共形结构的模空间。

\SourceParagraph{ch02-2.1.7-005}{55}
从刚才给出的闭弦图景推广到开闭弦理论，没有困难。我们不会在本书中深入讨论这些上链层面的理论，但重要的是要认识到，它们才是真正的目标。现在将指出它们的一些基本特点。Costello 的工作\BookCite{ref100}提供了详细得多的讨论。

\SourceParagraph{ch02-2.1.7-006}{56}
对于每一对边界条件 $a,b$，仍然有一个向量空间 $\mathcal{O}_{ab}$，实际上是一个上链复形，但它不再是某个范畴中从 $b$ 到 $a$ 的态射空间。相反，用 Fukaya、Kontsevich 等人的术语来说，我们拥有的是 $A_\infty$ 范畴。$A_\infty$ 范畴的概念将在全书中多次出现，而\BookChapterRef{8}讨论得最为详细。这里简要概括一下：这意味着，不再有一个单独的复合律 $\mathcal{O}_{ab}\times\mathcal{O}_{bc}\to\mathcal{O}_{ac}$，而是有一族复合方式，其参数空间是\BookRef{fig:2.1}{图1}中曲面上的共形结构的可缩空间。特别地，从这一族中任取两个复合律，它们都上链同伦。复合在如下意义下满足结合律：存在三重复合 $\mathcal{O}_{ab}\times\mathcal{O}_{bc}\times\mathcal{O}_{cd}\to\mathcal{O}_{ad}$ 的一个可缩族；从给定族中选择二重复合律，并以两种可能方式之一为三重乘积加括号所得的所有映射，都包含在这个族中。

\SourceParagraph{ch02-note-2.5}{56}
\BookStatementNumber{5}
\begin{note}[2.5]\label{note:2.5}
这不是定义 $A_\infty$ 结构的通常方式。按照 Stasheff 最初的定义，空间 $X$ 上的 $A_\infty$ 结构由一系列选择组成：首先，选择一个复合律 $m_2\colon X\times X\to X$；随后，选择映射
\[
m_3\colon[0,1]\times X\times X\times X\longrightarrow X,
\]
它是如下两个映射之间的同伦：
\[
(x,y,z)\longmapsto m_2(m_2(x,y),z)\quad\text{与}\quad
(x,y,z)\longmapsto m_2(x,m_2(y,z));
\]
接着，选择映射
\[
m_4\colon\overline{\mathcal{S}}_4\times X^4\longrightarrow X,
\]
其中 $\overline{\mathcal{S}}_4$ 是一个凸平面多边形，其顶点由为四重乘积加括号的五种方式标记，而 $m_4|_{(\partial\overline{\mathcal{S}}_4)\times X^4}$ 由 $m_3$ 确定；依此类推。对于上链复形而非空间，也有类似的定义，事实上还稍简单。精确定义见\BookSectionRef{sec:8.1}{8.1}。不过，这些定义本质上都等价于上面来自二维场论的定义；唯一重要的一点，是对每个 $k$ 都有一个 $k$ 重复合的可缩族。文献\BookCite{ref412}讨论了这些定义之间的关系。
\end{note}

\SourceParagraph{ch02-2.1.7-007}{56}
除复合律以外，我们在这些理论中发现的基本代数性质，是非退化内积，以及闭弦代数 $\mathcal{C}$ 的交换性。对于后者，转到上链理论时，$\mathcal{C}$ 中的乘法当然会在上链同伦意义下交换；但与开弦复合的情形不同，闭弦乘法的模空间 $\mathcal{M}_\Sigma$，也就是裤形曲面 $\Sigma$ 上共形结构模去在边界圆周上为恒等映射的微分同胚所得的模空间，并不可缩。它与如下嵌入方式的空间具有相同的同伦型：将标准圆盘 $D^2$ 的两个副本互不相交地嵌入 $D^2$ 的内部；这个嵌入空间当然是 $\mathcal{M}_\Sigma$ 的一个子空间。特别地，它包含一个自然的乘法圆周，其中一个嵌入圆盘像行星一样绕另一个圆盘移动，而且乘法与反向乘法之间有两个不同的自然同伦。这可能暗示弦性时空与经典时空之间的一个重要差异。闭弦上链复形 $\mathcal{C}$ 在弦理论中取代时空的 de Rham 复形，而后者是一个乘法严格结合、严格分次交换的代数。根据 Sullivan\BookCite{ref434}和 Quillen\BookCite{ref400}的工作，人们知道，在有限型幂零空间的有理同伦理论中，这类交换微分分次代数可以为同伦范畴提供模型；通过扩张标量，还可以得到实数或复数系数的模型\footnote{这里仍需注意基本群及其对高阶同伦群的作用所带来的微妙之处；幂零性条件正是对这些数据的限制。}。对于非负分次、连通且幂零的有限型有理 Sullivan 极小模型 $\mathcal{C}$，可以关联一个空间 $X_{\mathcal{C}}$，以及从 $\mathcal{C}$ 到 $X_{\mathcal{C}}$ 上有理分片多项式微分形式复形的上链代数拟同构；作为上链复形，这也是上链同伦等价。一般的交换微分分次代数需先取 Sullivan 模型，并以拟同构来比较；流形的实系数模型则可与通常的 de Rham 复形比较。如果不希望忽略空间同调中的挠元，就不能再用严格交换的上链代数编码同伦型，而必须将交换代数替换为所谓 $E_\infty$ 代数。粗略地说，这是整数域上的上链复形 $\mathcal{C}$，配备一个乘法，使它在给定的任意高阶同伦意义下满足结合律与交换律。任意空间 $X$ 都有一个上链的 $E_\infty$ 代数 $\mathcal{C}_X$；反过来，也可以给每个 $E_\infty$ 代数 $\mathcal{C}$ 关联一个空间 $X_{\mathcal{C}}$。因此，正如有理同伦论中那样，我们得到一对伴随函子。代数拓扑中长期的发展，最终达到了 Mandell 最近的定理\BookCite{ref346}：有限型幂零空间弱等价，当且仅当它们的整数系数奇异上链作为 $E_\infty$ 代数拟同构。因此，这类空间的弱同伦型可以由其 $E_\infty$ 上链代数确定；相应的上链反变函子忠实，但不充满。闭弦理论的上链代数，其高阶交换性弱于 $E_\infty$ 代数；这可能表明我们正在处理 Connes 意义下的非交换空间。这与将弦背景的 B 场解释为时空上的矩阵代数丛相当吻合。与此同时，$\mathcal{C}$ 上对应于 Poincaré 对偶的非退化内积，似乎说明，我们关注的是流形，而不是更奇异的空间。

\SourceParagraph{ch02-2.1.7-008}{57}
对于不习惯使用上链复形的读者，也许值得略加说明，这样做能带来什么好处。取最简单的例子，考虑由有限生成自由 Abel 群的有界上链复形，以及上链映射的上链同伦类构成的范畴 $\mathcal{K}$。这称为有限生成 Abel 群范畴的有界导出范畴。导出范畴将在\BookChapterRef{4}中详细讨论。取上同调给出从 $\mathcal{K}$ 到 $\mathbb{Z}$ 分次有限生成 Abel 群范畴的函子。实际上，$\mathcal{K}$ 中由除次数0以外上同调都为零的复形构成的子范畴 $\mathcal{K}_0$，与有限生成 Abel 群范畴确实等价\footnote{对于 Abel 群 $A$，可以关联上链复形
\[
C_A=(\cdots\longrightarrow0\longrightarrow R_A\longrightarrow F_A\longrightarrow0\longrightarrow\cdots),
\]
其中 $F_A$ 是位于次数0的自由 Abel 群，并带有满映射 $F_A\to A$；$R_A$ 为该映射的核。$F_A$ 的选择远不唯一，不过，$C_A$ 的不同选择仍然是 $\mathcal{K}$ 中典范同构的对象。}。不过，范畴 $\mathcal{K}$ 从有限生成自由 Abel 群范畴继承了一个对偶函子，具有几乎理想的性质：每个对象都与其二次对偶同构，并且取对偶保持正合列。复形 $C$ 的对偶 $C^*$ 定义为 $(C^*)^i=\operatorname{Hom}(C^{-i},\mathbb{Z})$。有限生成 Abel 群范畴中并没有这样好的对偶。事实上，子范畴 $\mathcal{K}_0$ 在对偶下不封闭，因为对应于群 $A$ 的复形 $C_A$ 的对偶，一般有两个非零上同调群：次数0处是 $\operatorname{Hom}(A,\mathbb{Z})$，次数 $+1$ 处是有限群 $\operatorname{Ext}^1(A,\mathbb{Z})$，它与 $A$ 的挠子群 Pontryagin 对偶。这来自以下正合列，注意不要把它与上链复形混淆：
\begin{equation}\label{eq:2.28}
\begin{aligned}
0&\longrightarrow\operatorname{Hom}(A,\mathbb{Z})\longrightarrow\operatorname{Hom}(F_A,\mathbb{Z})\\
&\longrightarrow\operatorname{Hom}(R_A,\mathbb{Z})\longrightarrow\operatorname{Ext}^1(A,\mathbb{Z})\longrightarrow0,
\end{aligned}
\end{equation}
而这一正合列由
\[
0\longrightarrow R_A\longrightarrow F_A\longrightarrow A\longrightarrow0
\]
导出。

\SourceParagraph{ch02-2.1.7-009}{58}
范畴 $\mathcal{K}$ 还有张量积，其性质比不保持正合列的 Abel 群张量积更好。更好的是，从局部性质良好的紧空间到 $\mathcal{K}$，存在一个典范的上链函子，把笛卡尔乘积映为张量积。空间的单纯上链复形、Čech 上链复形及其他候选上链复形，在 $\mathcal{K}$ 中都是典范同构的。

\SourceParagraph{ch02-2.1.7-010}{58}
我们将在\BookSectionRef{sec:2.5}{2.5}中回到这一讨论。

\section{代数条件的解：半单情形}\label{sec:2.2}
\subsection{分类定理}\label{sec:2.2.1}

\SourceParagraph{ch02-2.2.1-001}{58}
现在转向这样一个问题：给定闭弦理论 $\mathcal{C}$，相应的边界条件范畴是什么？在我们的表述中，这就变成：给定交换 Frobenius 代数 $\mathcal{C}$，可能的 $\mathcal{O}_{ab}$ 有哪些？

\SourceParagraph{ch02-2.2.1-002}{58}
当 $\mathcal{C}$ 半单时，可以回答这个问题。取 $\mathcal{C}$ 为复数域上的代数，此时半单性最有用的刻画，是“融合规则”
\begin{equation}\label{eq:2.29}
\phi_\mu\phi_\nu=N_{\mu\nu}^{\lambda}\phi_\lambda
\end{equation}
可以对角化\footnote{结构常数 $N_{\mu\nu}^{\lambda}$ 未必为整数；不过，在许多有趣的例子中，可以为代数选取一组基，使它们都为整数。}。也就是说，左正则表示的矩阵 $L(\phi_\mu)$，其矩阵元素为 $N_{\mu\nu}^{\lambda}$，可以同时对角化。

\SourceParagraph{ch02-2.2.1-003}{59}
等价地，存在一组基本幂等元 $\varepsilon_x$，使
\begin{equation}\label{eq:2.30}
\begin{gathered}
\mathcal{C}=\bigoplus_x\mathbb{C}\varepsilon_x,\\
\varepsilon_x\varepsilon_y=\delta_{xy}\varepsilon_y.
\end{gathered}
\end{equation}
再换一种等价表述，$\mathcal{C}$ 是有限集合 $X=\operatorname{Spec}(\mathcal{C})$ 上的复值函数代数，这个集合由 $\mathcal{C}$ 的特征标构成。

\SourceParagraph{ch02-2.2.1-004}{59}
迹 $\theta_{\mathcal{C}}\colon\mathcal{C}\to\mathbb{C}$ 应当看作有限时空 $\operatorname{Spec}(\mathcal{C})$ 上的“伸缩子场”。它完全由下列非零复数的无序族描述，其中相同的值仍计入重数：
\begin{equation}\label{eq:2.31}
\theta_x\coloneqq\theta_{\mathcal{C}}(\varepsilon_x).
\end{equation}
这是有限维交换半单 Frobenius 代数的唯一不变量。

\SourceParagraph{ch02-2.2.1-005}{59}
应当指出，复数域上最一般的有限维交换代数，具有形式 $\mathcal{C}=\bigoplus_x\mathcal{C}_x$，其中 $x$ 遍历 $\operatorname{Spec}(\mathcal{C})$，$\mathcal{C}_x$ 是局部环，即 $\mathcal{C}_x=\mathbb{C}\varepsilon_x\oplus\mathfrak{m}_x$，其中 $\varepsilon_x$ 如\BookEqRef{eq:2.30}{2.30}，$\mathfrak{m}_x$ 为幂零理想。如果 $\mathcal{C}$ 是 Frobenius 代数，那么每个 $\mathcal{C}_x$ 也是，并且存在某个 $\nu_x$，使 $\theta_{\mathcal{C}}\colon\mathfrak{m}_x^{\nu_x}\to\mathbb{C}$ 为同构，而 $\mathfrak{m}_x^{\nu_x+1}=0$。用 $\omega_x\in\mathfrak{m}_x^{\nu_x}$ 表示满足 $\theta_{\mathcal{C}}(\omega_x)=1$ 的元素。以各 $\omega_x$ 为分量的 $\mathcal{C}$ 的元素 $\omega$，可以看作时空上的“体积形式”。这类局部 Frobenius 代数 $\mathcal{C}_x$ 的一个典型例子，是复 $n$ 维射影空间 $\mathbb{P}^n$ 的复系数上同调环。它由一个满足 $t^{n+1}=0$ 的二次类 $t$ 生成，迹由在 $\mathbb{P}^n$ 上积分给出：当 $k=n$ 时，将 $t^k$ 映为1；其他情况下映为0。因此，这里 $\omega_x=t^n$。

\SourceParagraph{ch02-2.2.1-006}{59}
关于具有对称迹的 Frobenius 代数，不要求代数交换，有一个有用的技术事实：用\BookEqRef{eq:2.11}{2.11}中的记号，“Euler”元素 $\chi=\sum_\mu\psi_\mu\psi^\mu$ 可逆，当且仅当代数半单\footnote{为说明这一点，注意对代数中任意元素 $\psi$，都有 $\theta(\psi\chi)=\operatorname{tr}(\psi)$，其中 $\operatorname{tr}(\psi)$ 表示 $\psi$ 在正则表示中的迹。这是因为，$\theta(\psi^\mu\psi\psi_\nu)=\theta(\psi\psi_\nu\psi^\mu)$ 是正则表示中代表 $\psi$ 的矩阵的 $(\mu,\nu)$ 元素。配对 $(\psi_1,\psi_2)\mapsto\theta(\psi_1\psi_2)$ 非退化，因此迹形式 $(\psi_1,\psi_2)\mapsto\operatorname{tr}(\psi_1\psi_2)$ 非退化，当且仅当 $\chi$ 可逆。迹形式非退化是有限维代数半单的一个熟知判据。半单性有几种定义，而它们之间的等价正是经典的 Wedderburn 定理。对本书而言，半单代数就是若干全矩阵代数的直和。}；在一般的非交换情形下，这意味着代数同构于若干全矩阵代数的直和。元素 $\chi$ 总在代数的中心中；在交换情形下，其分量为 $\dim(\mathcal{C}_x)\omega_x$。

\SourceParagraph{ch02-2.2.1-007}{60}
在半单情形下，与固定闭弦代数 $\mathcal{C}$ 相容的开弦代数 $\mathcal{O}_{aa}$，具有如下完整刻画。

\SourceParagraph{ch02-thm-2.6}{60}
\BookStatementNumber{6}
\begin{theorem}\label{thm:2.6}
如果 $\mathcal{C}$ 半单，那么对每个 $a$，$\mathcal{O}=\mathcal{O}_{aa}$ 都半单，且必定具有形式 $\mathcal{O}=\operatorname{End}_{\mathcal{C}}(W)$，其中 $W$ 是 $\mathcal{C}$ 的某个有限维表示。
\end{theorem}

\SourceParagraph{ch02-proof-2.6}{60}
\begin{proof}
像 $\iota_a(\varepsilon_x)=P_x$ 是中心的单幂等元。因此，$\mathcal{O}_x=P_x\mathcal{O}=P_x\mathcal{O}P_x$ 是 Frobenius 代数 $\mathcal{C}_x=\varepsilon_x\mathcal{C}\cong\mathbb{C}$ 上的代数，所以只需在单个时空点上讨论。此时，对某个 $\alpha\in\mathbb{C}$，有 $\iota^a(1_{\mathcal{O}_x})=\alpha1_{\mathcal{C}_x}$。由 Cardy 条件，
\begin{equation}\label{eq:2.32}
\alpha1_{\mathcal{O}_x}=\chi_{\mathcal{O}_x}=\sum_\mu\psi_\mu\psi^\mu.
\end{equation}
施加 $\theta$，得到 $\alpha\theta(1_{\mathcal{O}_x})=\dim\mathcal{O}_x$；因此，当 $\mathcal{O}_x\neq0$ 时，$\alpha\neq0$，$\chi_{\mathcal{O}_x}$ 可逆。于是，每个点 $x$ 处的 $\mathcal{O}_x$ 都半单，即是矩阵代数的直和 $\bigoplus_i\operatorname{End}(W_i)$。实际上，Cardy 条件说明，每一点至多只有一个非零直和项 $W_i$，也就是说，代数是单的。因为映射 $\pi\colon\mathcal{O}_x\to\mathcal{O}_x$ 必须把每个直和项 $\operatorname{End}(W_i)$ 映入自身；如果有多个非零 $W_i$，它的像便不止一维，从而不能经由一维的 $\mathcal{C}_x$ 分解。\BookQED
\end{proof}

\SourceParagraph{ch02-2.2.1-008}{60}
根据\BookRef{thm:2.A}{定理A}，最一般的 $\mathcal{O}_{aa}$ 通过如下方式得到：为每个基本幂等元 $\varepsilon_x$ 选取向量空间 $W_{x,a}$，也就是在有限时空 $X=\operatorname{Spec}(\mathcal{C})$ 上选取向量丛，并形成
\begin{equation}\label{eq:2.33}
\mathcal{O}_{aa}=\bigoplus_x\operatorname{End}(W_{x,a}).
\end{equation}
不过，应当注意，对于形如 $\operatorname{End}(W)$ 的代数，向量空间 $W$ 由代数确定，但仍允许与任意复线作张量积这一不确定性：代数的任意不可约表示，都可以选作 $W$。

\SourceParagraph{ch02-2.2.1-009}{60}
将 $\psi\in\mathcal{O}_{aa}$ 记为 $\psi=\bigoplus_x\psi_x$。前面已经看到，投影到第 $x$ 个直和项上的投影算符 $P_x$ 由下式给出：
\begin{equation}\label{eq:2.34}
\iota_a(\varepsilon_x)=P_x.
\end{equation}
从伴随关系与 Cardy 条件，容易推出
\begin{equation}\label{eq:2.35}
\theta_a(\psi)=\sum_x\sqrt{\theta_x}\operatorname{Tr}_{W_{x,a}}(\psi_x),
\end{equation}
\begin{equation}\label{eq:2.36}
\iota^a(\psi)=\bigoplus_x\operatorname{Tr}_{W_{x,a}}(\psi_x)\frac{\varepsilon_x}{\sqrt{\theta_x}}.
\end{equation}
在 $\theta_{\mathcal{O}}$ 与 $\iota^a$ 的公式中，必须采用同一个平方根。注意
\[
\theta_{\mathcal{C}}\left(\frac{\varepsilon_x}{\sqrt{\theta_x}}\frac{\varepsilon_y}{\sqrt{\theta_y}}\right)=\delta_{x,y},
\]
也就是说，元素 $\varepsilon_x/\sqrt{\theta_x}$ 构成 $\mathcal{C}$ 的一组自然的标准正交基。因此，边界条件 $a$ 给出一组非负整数 $w_x=\dim W_x$，每个基本幂等元对应一个，并给出平方根 $\sqrt{\theta_x}$ 的选择。不过，关系\BookEqRef{eq:2.5}{2.5}说明，这些平方根是 Frobenius 范畴 $\mathcal{B}$ 的内在性质，不依赖于所考虑的具体对象。

\SourceParagraph{ch02-2.2.1-010}{61}
接下来确定与一对边界条件 $a,b$ 关联的 $\mathcal{O}_{aa}$--$\mathcal{O}_{bb}$ 双模 $\mathcal{O}_{ab}$。它们同样由 Cardy 条件确定。

\SourceParagraph{ch02-lem-2.7}{61}
\BookStatementNumber{7}
\begin{lemma}\label{lem:2.7}
当 $\mathcal{C}$ 半单时，有
\begin{equation}\label{eq:2.37}
\mathcal{O}_{ab}\cong\bigoplus_x\operatorname{Hom}(W_{x,b},W_{x,a}).
\end{equation}
\end{lemma}

\SourceParagraph{ch02-proof-2.7}{61}
\begin{proof}
限制到各个 $\mathcal{O}_{aa}$ 上，就可以使用\BookRef{thm:2.6}{定理2.6}。于是，
\[
\iota_a(\varepsilon_x)\mathcal{O}_{ab}=\mathcal{O}_{ab}\iota_b(\varepsilon_x)
\]
是单代数 $\mathcal{O}_{x,aa}$ 和 $\mathcal{O}_{x,bb}$ 的双模。我们限制到一个幂等元，并省去 $x$，也就是取 $\mathcal{C}=\mathbb{C}$。若 $W_a=0$ 或 $W_b=0$，单位元的作用便迫使 $\mathcal{O}_{ab}=0$，结论成立；以下设二者均非零。$\mathcal{O}_{aa}=\operatorname{End}(W_a)$ 唯一的不可约表示是 $W_a$ 本身，而唯一的不可约 $\mathcal{O}_{aa}$--$\mathcal{O}_{bb}$ 双模是 $W_a\otimes W_b^*$。因此，$\mathcal{O}_{ab}\cong n_{ab}W_a\otimes W_b^*$，其中 $n_{ab}$ 是非负整数。现在计算 Cardy 条件。如果 $v_m$ 是 $W_a$ 的一组基，$w_n$ 是 $W_b$ 的一组基，那么 $\mathcal{O}_{ab}$ 的一组基为 $v_{m,\alpha}\otimes w_{n,\alpha}^*$，其中 $\alpha=1,\ldots,n_{ab}$。于是，$\pi_b^a(\psi)=n_{ab}\operatorname{Tr}_{W_a}(\psi)P_b/\sqrt{\theta_x}$。与 $\iota_b\iota^a(\psi)$ 比较，得到 $n_{ab}=1$。\BookQED
\end{proof}

\SourceParagraph{ch02-2.2.1-011}{61}
这个引理的一个推论是
\begin{equation}\label{eq:2.38}
\pi_b^a(\psi)=\bigoplus_x\frac{1}{\sqrt{\theta_x}}\operatorname{Tr}_{W_{x,a}}(\psi_x)P_{x,b}.
\end{equation}

\SourceParagraph{ch02-2.2.1-012}{61}
现在可以描述边界条件的极大范畴 $\mathcal{B}$。首先，若 $p\in\mathcal{O}_{aa}$ 为投影，即 $p^2=p$，就可以假定 $\mathcal{B}$ 中有 $a=b\oplus c$，其中 $b$ 是 $p$ 的像。因为向加性范畴加入投影的像，与加入直和的方式非常相似。因此，如果闭弦代数 $\mathcal{C}$ 半单，就可以为每个时空点 $x$ 选取 $\mathcal{B}$ 的对象 $a_x$，使它的支集位于 $x$，即 $\iota_{a_x}(\varepsilon_x)\mathcal{O}_{a_xa_x}=\mathcal{O}_{a_xa_x}$，并且是单的，即 $\mathcal{O}_{a_xa_x}=\mathbb{C}$。对于 $\mathcal{B}$ 中的任意对象 $b$，有典范态射
\begin{equation}\label{eq:2.39}
\bigoplus_x\mathcal{O}_{ba_x}\otimes a_x\longrightarrow b,
\end{equation}
这里在左侧使用了线性范畴的任意对象可以与有限维向量空间作张量积这一事实。此外，由引理可知，态射\BookEqRef{eq:2.39}{2.39}是同构，因为两侧到任意其他对象 $c$ 的态射空间相同。最后，注意 $a_x$ 的选择只差与一条线 $L_x$ 作张量积：如果 $a'_x$ 是另一种选择，那么 $a'_x\cong a_x\otimes L_x$，其中 $L_x=\mathcal{O}_{a'_xa_x}$。

\SourceParagraph{ch02-thm-2.8}{62}
\BookStatementNumber{8}
\begin{theorem}\label{thm:2.8}
假设 $\mathcal{C}$ 半单，对应于时空 $X$。那么：
\begin{enumerate}[label=(\roman*)]
\item 边界条件的极大范畴 $\mathcal{B}$ 等价于 $X$ 上的向量丛范畴 $\operatorname{Vect}(X)$，互为逆的函子为
\begin{equation}\label{eq:2.40}
W\longmapsto\bigoplus_xW_x\otimes a_x,
\end{equation}
其中 $W$ 是 $X$ 上的向量丛，$W_x$ 是它在 $x\in X$ 处的纤维；以及
\begin{equation}\label{eq:2.41}
a\longmapsto\{\mathcal{O}_{aa_x}\},
\end{equation}
其中右侧表示 $X$ 上的向量丛，在 $x$ 处的纤维为 $\mathcal{O}_{aa_x}$。
\item $\mathcal{B}$ 与 $\operatorname{Vect}(X)$ 之间的等价，唯一到由与 $X$ 上线丛 $L=\{L_x\}$ 作张量积给出的变换 $\operatorname{Vect}(X)\to\operatorname{Vect}(X)$ 的差别。
\item $\mathcal{B}$ 上的 Frobenius 结构，由伸缩子场的平方根 $\{\sqrt{\theta_x}\}$ 的选择确定。因此，它唯一到与一个满足 $\sigma^2=1$ 的元素 $\sigma\in\mathcal{C}$ 相乘的差别。
\end{enumerate}
\end{theorem}

\SourceParagraph{ch02-note-2.9}{62}
\BookStatementNumber{9}
\begin{note}[2.9]\label{note:2.9}
\begin{enumerate}[label=(\arabic*)]
\item 边界条件 $a$ 具有支集
\begin{equation}\label{eq:2.42}
\operatorname{supp}(a)=\{x\in X\mid W_x\neq0\},
\end{equation}
它包含于 $X=\operatorname{Spec}(\mathcal{C})$。如果两个边界条件 $a$ 和 $b$ 具有相同的支集，那么 $\mathcal{O}_{ab}$ 是 $\mathcal{O}_{aa}$ 与 $\mathcal{O}_{bb}$ 之间的 Morita 等价双模。读者不妨将这一讨论与文献\BookCite{ref418}的第6.4节比较。注意，必须使用 Cardy 条件，才能得到这一结论。
\item 物理中的半单 Frobenius 代数例子包括：
\begin{enumerate}[label=(\alph*)]
\item 有理共形场论（RCFT）的融合规则代数，即 Verlinde 代数。
\item $N=2$ Landau--Ginzburg 理论在一般超势 $W$ 下的手征环，也就是 $W$ 的所有临界点都是 Morse 临界点的情形。这是红外（IR）理论有质量的情况。
\item 流形在一般参数下的量子上同调。
\end{enumerate}
\end{enumerate}
\end{note}

\subsection{关于 B 场的说明}\label{sec:2.2.2}

\SourceParagraph{ch02-2.2.2-001}{62}
从这一讨论可以准确看出，B 场的思想会在何处出现\footnote{B 场的标准物理定义，将在\BookSectionRef{sec:3.2.6}{3.2.6}、\BookSectionRef{sec:3.3.2}{3.3.2}和\BookSectionRef{sec:3.5.2.7}{3.5.2.7}中讨论。}；不过，实际上在零维时空上，任何 B 场都必须平凡。我们已经说明，与时空的每一点都关联着一个边界条件范畴，它与有限维向量空间范畴同构，但并非典范同构。更精确地说，它包含极小的、也就是不可约的对象，而任意其他对象都可以通过将这样的对象与有限维向量空间作张量积得到。

\SourceParagraph{ch02-2.2.2-002}{63}
本质上，B 场就是时空上的范畴丛，其中的纤维范畴彼此同构，但并非典范地同构。可以假定，每个纤维都同构于有限维向量空间范畴。关键之处是：将每个纤维与标准纤维等同起来时，其不确定性由一个等价的“群”描述；在这里，它实际上是一个范畴，其中的元素是复线，而复合由张量积给出。我们的边界条件范畴，恰好就是具有这个结构群的范畴丛的“截面”范畴。

\SourceParagraph{ch02-2.2.2-003}{63}
也许可以用如下方式帮助理解。电磁场是时空上带联络的线丛，可以看作时空结构的一部分，即使没有费米子也有意义。但是，在含有费米子的理论中，时空的每一点都有旋量空间，电磁场“实际上”就是关于这些旋量空间如何沿时空逐点连接起来的信息。在这一意义下，电磁场“就是”旋量丛及其联络。类似地，B 场“就是”边界条件的丛。

\SourceParagraph{ch02-2.2.2-004}{63}
在一般拓扑空间 $X$ 上，B 场的类由上同调群 $H^3(X,\mathbb{Z})$ 的元素分类；这个群也可以理解为 $H^1(X,\mathcal{G})$，其中 $\mathcal{G}$ 是线丛在张量积下构成的“群”，在代数拓扑中是 $K(\mathbb{Z},2)$ 型的 Eilenberg--MacLane 对象。我们将在\BookSectionRef{sec:2.6}{2.6}中回到这个主题。

\subsection{重建闭弦代数}\label{sec:2.2.3}

\SourceParagraph{ch02-2.2.3-001}{63}
在开闭弦 TFT 中，闭弦代数 $\mathcal{C}$ 的每个元素 $\xi$ 都为 $\mathcal{B}$ 的每个对象 $a$ 定义一个自同态 $\xi_a=\iota_a(\xi)\in\mathcal{O}_{aa}$；并且对每个从 $a$ 到 $b$ 的态射 $\eta\in\mathcal{O}_{ba}$，都有 $\eta\circ\xi_a=\xi_b\circ\eta$。因此，族 $\{\xi_a\}$ 构成从恒等函子 $1_{\mathcal{B}}\colon\mathcal{B}\to\mathcal{B}$ 到自身的自然变换。另见\BookRef{def:4.10}{定义4.10}。

\SourceParagraph{ch02-2.2.3-002}{63}
对于任意 $\mathbb{C}$ 线性范畴 $\mathcal{B}$，都可以考虑 $1_{\mathcal{B}}$ 的自然变换所构成的环 $\mathcal{E}$。它自动交换，因为如果 $\{\xi_a\},\{\eta_a\}\in\mathcal{E}$，那么由自然性的定义，$\xi_a\circ\eta_a=\eta_a\circ\xi_a$。从 $1_{\mathcal{B}}$ 到自身的自然变换，就是一族元素 $\{\xi_a\in\mathcal{O}_{aa}\}$，使得对每个从 $b$ 到 $a$ 的态射 $f\in\mathcal{O}_{ab}$，都有 $\xi_a\circ f=f\circ\xi_b$；这里可以取 $a=b$ 且 $f=\eta_a$。如果 $\mathcal{B}$ 是 Frobenius 范畴，那么每一对对象 $a,b$ 都有映射 $\pi_a^b\colon\mathcal{O}_{bb}\to\mathcal{O}_{aa}$，并可定义 $j^b\colon\mathcal{O}_{bb}\to\mathcal{E}$：对于 $\eta\in\mathcal{O}_{bb}$，令 $j^b(\eta)_a=\pi_a^b(\eta)$。换言之，$j^b$ 的定义使 Cardy 条件 $\iota_a\circ j^b=\pi_a^b$ 成立，其中 $\iota_a$ 对自然变换取其在 $a$ 上的分量。不过，问题是：能否定义迹 $\theta\colon\mathcal{E}\to\mathbb{C}$，使 $\mathcal{E}$ 成为 Frobenius 代数，并使对所有 $\xi\in\mathcal{E}$ 和 $\eta\in\mathcal{O}_{aa}$ 都有
\begin{equation}\label{eq:2.43}
\theta_a\bigl(\iota_a(\xi)\eta\bigr)=\theta\bigl(\xi j^a(\eta)\bigr)?
\end{equation}
如果 $\mathcal{B}$ 是只有有限多个单对象同构类的半单 Frobenius 范畴，这当然成立；因为此时 $\mathcal{E}$ 就是这些单对象的同构类集合上的复值函数环，可以直接定义 $\theta\colon\mathcal{E}\to\mathbb{C}$，使 $\theta(\varepsilon_a)=\theta_a(1_a)^2$，其中 $a$ 是单对象，$\varepsilon_a\in\mathcal{E}$ 为 $\mathcal{E}$ 的谱中点 $a$ 的特征函数。然而，Frobenius 范畴未必半单；遗憾的是，在一般情形下，不能将 $\mathcal{E}$ 作为闭弦代数。例如，如果 $\mathcal{B}$ 只有一个对象 $a$，且 $\mathcal{O}_{aa}$ 是维数大于1的交换局部环，那么 $\mathcal{E}=\mathcal{O}_{aa}$，因而 $\iota_a\colon\mathcal{E}\to\mathcal{O}_{aa}$ 是同构，它的伴随映射 $j^a$ 也应当是同构。但这与 Cardy 条件矛盾，因为 $\pi_a^a$ 是乘以 $\sum_i\psi_i\psi^i$ 的映射，而这个元素必为幂零的。\BookSectionRef{sec:2.6}{2.6}将给出一个例子，其中两个不同的闭弦 Frobenius 代数允许同一个开弦代数 $\mathcal{O}_{aa}$。

\SourceParagraph{ch02-2.2.3-003}{64}
线性范畴 $\mathcal{B}$ 的恒等函子的自然自同态所构成的交换代数 $\mathcal{E}$，称为 $\mathcal{B}$ 的零次 Hochschild 上同调 $\operatorname{HH}^0(\mathcal{B})$。当 $\mathcal{B}$ 半单时，群 $\operatorname{HH}^p(\mathcal{B})$ 对 $p>0$ 都为零，稍后将给出它们的定义；但在一般情形下，它们似乎与从 $\mathcal{B}$ 构造闭弦代数有关。与此同时，注意对任意 Frobenius 范畴 $\mathcal{B}$，都存在从其 Grothendieck 群\footnote{也就是由 $\mathcal{B}$ 的对象同构类在 $\oplus$ 下构成的半群，通过群化得到的群。}到零次 Hochschild 上同调的自然同态 $K(\mathcal{B})\to\operatorname{HH}^0(\mathcal{B})$：它将对象 $a$ 映为在 $b$ 上的值为 $\pi_b^a(1_a)\in\mathcal{O}_{bb}$ 的自然变换。在半单情形下，这个同态诱导同构 $K(\mathcal{B})\otimes\mathbb{C}\to\operatorname{HH}^0(\mathcal{B})$。

\SourceParagraph{ch02-2.2.3-004}{64}
对于任意加性范畴 $\mathcal{B}$，Hochschild 上同调定义为如下上链复形的上同调。一个 $k$ 上链 $F$ 是对各态射变量分别可加的法则，在线性范畴中取多重线性的法则：为每个可复合的 $k$ 元态射组
\begin{equation}\label{eq:2.44}
Y_0\xrightarrow{\phi_1}Y_1\xrightarrow{\phi_2}\cdots\xrightarrow{\phi_k}Y_k
\end{equation}
指定 $F(\phi_1,\ldots,\phi_k)\in\operatorname{Hom}(Y_0,Y_k)$。复形中的微分定义为
\begin{equation}\label{eq:2.45}
\begin{aligned}
(dF)(\phi_1,\ldots,\phi_{k+1})
&=F(\phi_2,\ldots,\phi_{k+1})\circ\phi_1\\
&\quad+\sum_{i=1}^k(-1)^iF(\phi_1,\ldots,\phi_{i+1}\circ\phi_i,\ldots,\phi_{k+1})\\
&\quad+(-1)^{k+1}\phi_{k+1}\circ F(\phi_1,\ldots,\phi_k).
\end{aligned}
\end{equation}
特别地，零上链为每个对象 $Y$ 指定一个自同态 $F_Y$，而当这些自同态构成自然变换时，它就是上闭链。同样，二上链 $F$ 为范畴的复合律 $(\phi_1,\phi_2)\mapsto\phi_2\circ\phi_1$ 给出一种可能的无穷小形变 $F(\phi_1,\phi_2)$；形变保持复合的结合律，当且仅当 $F$ 是上闭链。

\SourceParagraph{ch02-2.2.3-005}{65}
如果范畴 $\mathcal{B}$ 只有一个对象，其自同态代数为 $\mathcal{O}$，那么刚才描述的上同调通常称为代数 $\mathcal{O}$ 以自身这个 $\mathcal{O}$ 双模为系数的 Hochschild 上同调。必须仔细区分它与以对偶 $\mathcal{O}$ 双模 $\mathcal{O}^*$ 为系数的 Hochschild 上同调。不过，如果 $\mathcal{O}$ 是具有对称迹的 Frobenius 代数，那么它作为双模同构于 $\mathcal{O}^*$，因而不必再区分这两种 Hochschild 上同调概念。Frobenius 范畴 $\mathcal{B}$ 也一样：$\operatorname{Hom}(Y_k,Y_0)$ 是 $\operatorname{Hom}(Y_0,Y_k)$ 的对偶空间，所以可以把 $k$ 上链理解为这样的法则：为每个可复合的 $k$ 元态射组\BookEqRef{eq:2.44}{2.44}，指定关于元素 $\phi_0\in\operatorname{Hom}(Y_k,Y_0)$ 的一个线性函数。换言之，一个 $k$ 上链，就是为每个由 $k+1$ 个态射组成的“圆环”
\begin{equation}\label{eq:2.46}
\cdots\xrightarrow{\phi_0}Y_0\xrightarrow{\phi_1}Y_1\xrightarrow{\phi_2}\cdots\xrightarrow{\phi_k}Y_k\xrightarrow{\phi_0}\cdots
\end{equation}
指定一个复数 $F(\phi_0,\phi_1,\ldots,\phi_k)$ 的法则。

\begin{figure}[H]
\centering
\BookImage{01_Numbered_Figures/Chapter_02/Ch02_Fig13_PDFp065_Bookp055.pdf}
\caption{闭弦态 $\phi$ 与 $k+1$ 个开弦态的循环配对。}\label{fig:2.13}
\end{figure}

\SourceParagraph{ch02-2.2.3-006}{65}
在这一描述中，如果限制到按循环上链的符号约定，在旋转态射圆环 $(\phi_0,\phi_1,\ldots,\phi_k)$ 时循环不变的上链；在此非分次情形，这一条件写为
\[
F(\phi_0,\ldots,\phi_k)=(-1)^kF(\phi_k,\phi_0,\ldots,\phi_{k-1}),
\]
就得到 Hochschild 复形的一个上链子复形，其上同调称为范畴 $\mathcal{B}$ 的循环上同调 $\operatorname{HC}^*(\mathcal{B})$。循环上同调显然映入 Hochschild 上同调；与后者相比，它是与 $\mathcal{B}$ 关联的闭弦代数更自然的候选。例如，由态射组\BookEqRef{eq:2.46}{2.46}代表的态，按照\BookRef{fig:2.13}{图13}，以循环不变的方式与闭弦态配对，给出一个数。在我们的初级例子所用的零次部分，循环上同调与 Hochschild 上同调无法区分；但值得指出\footnote{这一点承蒙 Kontsevich 告知。}，正如前面说明的那样，$\operatorname{HH}^2(\mathcal{B})$ 是 $\mathcal{B}$ 作为范畴的无穷小形变空间，而群 $\operatorname{HC}^2(\mathcal{B})$ 是它作为 Frobenius 范畴的无穷小形变空间。

\SourceParagraph{ch02-2.2.3-007}{66}
可以用来检验这些思想的一个非常自然的 Frobenius 范畴例子，是紧 Calabi--Yau 流形上全纯向量丛的范畴。\BookSectionRef{sec:2.5}{2.5}将讨论这一例子。

\subsection{自旋理论与模2分次范畴}\label{sec:2.2.4}

\SourceParagraph{ch02-2.2.4-001}{66}
为了描述\BookSectionRef{sec:2.1.6}{2.1.6}所定义的拓扑自旋理论的边界条件范畴，需要对前面的讨论作一些修改。这里给出一个简短概述，不作证明。

\SourceParagraph{ch02-2.2.4-002}{66}
区间上只有一种自旋结构，其自同构群为 $\{\pm1\}$，所以对于每一对边界条件 $a,b$，向量空间 $\mathcal{O}_{ab}$ 都具有一个对合，也就是模2分次。双线性复合 $\mathcal{O}_{ab}\otimes\mathcal{O}_{bc}\to\mathcal{O}_{ac}$ 保持分次。还存在非退化迹 $\theta_a\colon\mathcal{O}_{aa}\to\mathbb{C}$，满足不带符号的交换性条件\BookEqRef{eq:2.5}{2.5}。

\SourceParagraph{ch02-2.2.4-003}{66}
如果闭弦理论如\BookSectionRef{sec:2.1.6}{2.1.6}所述，由 Frobenius 代数 $\mathcal{C}=\mathcal{C}_{\mathrm{ns}}\oplus\mathcal{C}_{\mathrm{r}}$ 描述，那么存在保持分次的伴随映射
\begin{equation}\label{eq:2.47}
\begin{aligned}
\iota_a^{\mathrm{ns}}&\colon\mathcal{C}_{\mathrm{ns}}\longrightarrow\mathcal{O}_{aa},\\
\iota_{\mathrm{ns}}^a&\colon\mathcal{O}_{aa}\longrightarrow\mathcal{C}_{\mathrm{ns}},\\
\iota_a^{\mathrm{r}}&\colon\mathcal{C}_{\mathrm{r}}\longrightarrow\mathcal{O}_{aa},\\
\iota_{\mathrm{r}}^a&\colon\mathcal{O}_{aa}\longrightarrow\mathcal{C}_{\mathrm{r}}.
\end{aligned}
\end{equation}
此外，映射 $\iota_a^{\mathrm{ns}}$ 和 $\iota_a^{\mathrm{r}}$ 合在一起，定义代数同态 $\mathcal{C}\to\mathcal{O}_{aa}$。中心性条件变为
\begin{equation}\label{eq:2.48}
\begin{aligned}
\iota_a^{\mathrm{ns}}(\phi)\psi&=\psi\iota_a^{\mathrm{ns}}(\phi),\\
\iota_a^{\mathrm{r}}(\phi)\psi&=(-1)^{\deg\phi\deg\psi+\deg\psi}\psi\iota_a^{\mathrm{r}}(\phi).
\end{aligned}
\end{equation}
因此，$\iota^{\mathrm{ns}}$ 映入代数 $\mathcal{O}_{aa}$ 的朴素中心。在这里得到朴素中心而非分次代数中心的原因，以及迹在朴素意义下满足交换性的原因，都与\BookSectionRef{sec:2.1.6}{2.1.6}中代数 $\mathcal{C}$ 的朴素交换性的原因相同。$\iota^{\mathrm{r}}$ 的符号，则通过仔细追踪在\BookRef{fig:2.9}{图9}的微分同胚下，对自旋丛截面所作的选择得到。

\SourceParagraph{ch02-2.2.4-004}{66}
有两个 Cardy 条件：
\begin{equation}\label{eq:2.49}
\begin{aligned}
\iota_b^{\mathrm{ns}}\iota_{\mathrm{ns}}^a(\psi)
&=\pi_b^a(\psi)\coloneqq\sum_\mu(-1)^{\deg\psi_\mu\deg\psi}\psi_\mu\psi\psi^\mu,\\
\iota_b^{\mathrm{r}}\iota_{\mathrm{r}}^a(\psi)
&=\widetilde{\pi}_b^a(\psi)\coloneqq\sum_\mu(-1)^{\deg\psi_\mu(\deg\psi+1)}\psi_\mu\psi\psi^\mu.
\end{aligned}
\end{equation}

\SourceParagraph{ch02-2.2.4-005}{67}
如果假设闭弦代数半单，那么与前面一样，可以假定 $\mathcal{C}_{\mathrm{ns}}$ 是有限集合 $X$ 上的函数代数，并逐点确定边界条件范畴。换言之，在单个点上可以取 $\mathcal{C}=\mathbb{C}[\eta]/(\eta^2-1)$，其中 $\mathcal{C}_{\mathrm{r}}$ 的生成元 $\eta$ 可以具有偶或奇次数。在这两种情形中，前面使用的论证通过第一个 Cardy 公式说明，开弦代数 $\mathcal{O}_{aa}$ 是某个向量空间 $W$ 的全矩阵代数。如果 $\eta$ 的次数为偶，那么 $\iota^{\mathrm{r}}(\eta)=P$，其中 $P$ 为偶、$P^2=1$，并满足 $P\psi P=(-1)^{\deg\psi}\psi$。此时，该点处的边界条件范畴等价于模2分次向量空间范畴。另一方面，如果 $\eta$ 的次数为奇，那么 $P$ 为奇、$P^2=1$，而 $P$ 位于朴素中心。取不可约分次表示空间 $W$，并令 $V=W^0$。代数 $\mathcal{O}_{aa}$ 的对合，对应于模 $W$ 的一个对合，$P$ 的作用则给出分次的两半之间的同构。$\mathcal{O}_{aa}$ 的偶子代数是全矩阵代数；整个分次代数的形式为 $\operatorname{End}(V)\otimes\mathbb{C}[\eta]/(\eta^2-1)$，其中 $\eta$ 为奇。因此，边界条件范畴等价于超代数 $\mathbb{C}[\eta]/(\eta^2-1)$ 的分次表示范畴，而后者又等价于通常的不分次向量空间范畴。与不分次情形一样，开弦代数的 Frobenius 结构，通过平方关系确定闭弦代数的 Frobenius 结构。$\deg\eta=0$ 与 $\deg\eta=1$ 这两种情形，大致类似于复数域上分别具有偶数个与奇数个生成元的 Clifford 代数的区别。

\SourceParagraph{ch02-2.2.4-006}{67}
反过来，假设给定一个任意的半单模2分次范畴 $\mathcal{B}$，带有相容的 Frobenius 结构，且只有有限多个单对象同构类。这里的模2分次范畴，是一个配备对合函子 $S$ 的线性范畴，$S$ 被理解为分次的翻转。这样的范畴有两类单对象 $P$：一类满足 $S(P)\cong P$，另一类则不满足。第一类对象生成 $\mathcal{B}$ 中一个同构于向量空间范畴的子范畴；第二类对象及其分次翻转，则生成一个同构于分次向量空间范畴的子范畴。因此，任意这样的半单分次 Frobenius 范畴 $\mathcal{B}$，都是某个唯一的拓扑自旋理论的边界条件范畴。

\section{向量丛、K 理论与“边界态”}\label{sec:2.3}

\SourceParagraph{ch02-2.3-001}{67}
在半单情形下，边界条件范畴 $\mathcal{B}$ 有一个很好的几何解释：其可能的对象，对应于与 $\mathcal{C}$ 关联的“时空”$X=\operatorname{Spec}(\mathcal{C})$ 上的向量丛；这个时空只是一个有限点集。点 $x$ 上方的纤维，就是向量空间 $W_x$。

\SourceParagraph{ch02-2.3-002}{67}
现在对“边界态”作一些说明。正如我们将在\BookSectionRef{sec:3.5}{3.5}中讨论的那样，在共形场论中，一个边界条件 $a$ 关联着闭弦态空间中相应的“态”$B_a$。（严格地说，$B_a$ 是代数对偶空间的一个元素。）

\begin{figure}[H]
\centering
\BookImage{01_Numbered_Figures/Chapter_02/Ch02_Fig14_PDFp068_Bookp058.pdf}
\caption{具有边界条件 $a$ 的上半平面上的关联函数，等于插入边界态 $B_a$ 后的闭弦振幅。}\label{fig:2.14}
\end{figure}

\SourceParagraph{ch02-2.3-003}{68}
在当前的情境下，$B_a\in\mathcal{C}$。边界态的定义性质如下：在边界条件为 $a$ 的圆盘上，算符的关联函数，等于相应闭弦理论在球面上的关联函数；这个球面通过用另一个圆盘封住原圆盘而得到，并在封口圆盘的中心插入态 $B_a$。\BookRef{fig:2.14}{图14}说明了这一点。

\SourceParagraph{ch02-2.3-004}{68}
用方程表示，就是对任意 $\phi_1,\ldots,\phi_n$ 都有
\begin{equation}\label{eq:2.50}
\theta_a\bigl(\iota_a(\phi_1)\cdots\iota_a(\phi_n)\bigr)
=\theta_{\mathcal{C}}(B_a\phi_1\cdots\phi_n).
\end{equation}
利用伴随关系与迹的非退化性，得到
\begin{equation}\label{eq:2.51}
B_a=\iota^a(1_{\mathcal{O}_{aa}}).
\end{equation}

\SourceParagraph{ch02-2.3-005}{68}
映射 $a\mapsto B_a$ 给出自然同态
\begin{equation}\label{eq:2.52}
K(\mathcal{B})\longrightarrow\mathcal{C}.
\end{equation}

\SourceParagraph{ch02-2.3-006}{68}
更一般地，如果 $\chi$ 是 $\mathcal{C}$ 的 Euler 元素，就可以把 $\mathcal{C}$ 中的元素 $\chi^g\prod_a(B_a)^{h_a}$ 看作“添加 $g$ 个柄和 $h=\sum_a h_a$ 个孔洞，其中有 $h_a$ 个孔洞的边界条件为 $a$”的算符。其含义是：在具有 $h$ 个孔洞、亏格为 $g$、并带有给定边界条件的曲面上，一组算符的关联函数，等于这些算符在球面上的关联函数，只需再插入一个 $\chi^g\prod_a(B_a)^{h_a}$。

\SourceParagraph{ch02-2.3-007}{68}
记录这些边界态的一个简单性质。首先，利用 Cardy 条件，有
\begin{equation}\label{eq:2.53}
\begin{aligned}
\theta_{\mathcal{C}}(B_aB_b)
&=\theta_a\bigl(\iota_a(B_b)\bigr)\\
&=\theta_a\bigl(\iota_a\iota^b(1_b)\bigr)\\
&=\theta_a\bigl(\pi_a^b(1_b)\bigr)\\
&=\dim\mathcal{O}_{ab}.
\end{aligned}
\end{equation}

\SourceParagraph{ch02-2.3-008}{69}
在半单情形下，公式\BookEqRef{eq:2.35}{2.35}给出用基本幂等元表示“边界态”的显式公式：
\begin{equation}\label{eq:2.54}
B_a=\iota^a(1_{\mathcal{O}_{aa}})
=\sum_x(\dim W_x)\frac{\varepsilon_x}{\sqrt{\theta_x}}.
\end{equation}

\SourceParagraph{ch02-2.3-009}{69}
这个公式说明，边界态在闭弦代数 $\mathcal{C}$ 中的幺模格 $\mathcal{L}_B$ 内构成一个正锥；这个格由正交归一基 $\varepsilon_x/\sqrt{\theta_x}$ 张成。特别地，由\BookEqRef{eq:2.54}{2.54}可知，边界态只能以非负整数系数相加。因此，它们并不像量子力学中的膜态。这一基本的整数结构，是 Cardy 条件的结果。

\SourceParagraph{ch02-2.3-010}{69}
人们自然会猜想，边界条件之间是否应当有一种“乘法”运算。支持与反对这一想法，都有理由。最初对 D 膜的看法，是将它们视为时空中的“闭链”，开弦可以在这些闭链上开始和终止；这提示应当存在一个乘法，对应于闭链的相交。本章建立的简化结构似乎没有产生任何乘法，因此，人们或许会怀疑，我们是否遗漏了某个重要的要素。不过，反对这一想法的考虑如下。我们的边界条件似乎更接近于时空上的向量丛，也就是 K 理论类，而非同调闭链；在\BookSectionRef{sec:2.6}{2.6}考虑等变情形时，这一点会更加清楚。对于一个环，只有在它交换时，模的张量积才自然给出其 K 理论类之间的乘积；而我们已经指出，作为闭弦时空模型一部分的 B 场，似乎编码了某种程度的非交换性。更准确地说，D 膜似乎定义了时空的扭曲 K 理论中的类，其中的扭曲由 B 场给出；一个空间的固定扭曲 K 理论不构成环：两个扭曲类的乘积，其扭曲是两个因子的扭曲之和。然而，在弦论中，并不存在通过“关掉”B 场，来得到一个潜在的无扭曲时空的概念。例如，一个具有非零 B 场的环面所对应的共形场论，可以通过“T 对偶”，同构于另一个不带 B 场的环面所给出的理论。

\SourceParagraph{ch02-2.3-011}{69}
不应期待 D 膜之间有乘法运算的另一个理由，也来自共形场论中的 T 对偶。那里，由黎曼环面 $T$ 及其对偶环面 $T^*$ 定义的闭弦理论是同构的，而且确实有 K 理论同构 $K(T)\cong K(T^*)$，但它与 K 理论中的乘法不相容。另一方面，在来自 $\mathcal{N}=2$ 超对称 $\sigma$ 模型的某些 TFT 例子中，边界条件范畴却似乎是一个张量范畴。

\SourceParagraph{ch02-2.3-012}{70}
边界态公式\BookEqRef{eq:2.54}{2.54}说明，由伸缩子场 $\theta$ 在 $\mathcal{C}$ 内选出的格 $\mathcal{L}_B$，通常并不对 $\mathcal{C}$ 中的乘法封闭；它对乘法封闭的条件，是对每个点 $x$ 都有 $1/\sqrt{\theta_x}\in\mathbb{Z}$，如果还要求这个格包含单位元，则必须对所有 $x$ 都有 $\theta_x=1$。但对应于不同伸缩子场的格，可以相乘并落入另一个格中，正如扭曲 K 类的情形一样。不过，在半单情形下，如果定义元素 $S\coloneqq\sum_x\sqrt{\theta_x}\varepsilon_x$，那么运算
\begin{equation}\label{eq:2.55}
(B_1,B_2)\longmapsto SB_1B_2
\end{equation}
确实定义了边界态之间的一个乘法，只是其意义尚不清楚。

\subsection{“Cardy 态”与“Ishibashi 态”}\label{sec:2.3.1}

\SourceParagraph{ch02-2.3.1-001}{70}
边界态公式\BookEqRef{eq:2.54}{2.54}，让人联想到边界共形场论中“Cardy 态”与“Ishibashi 态”之间的关系\BookCite{ref086,ref087}。对于熟悉这一关系的读者，我们简要说明这种相似之处，因为\BookEqRef{eq:2.54}{2.54}的两端，确实对应于这两组基。

\SourceParagraph{ch02-2.3.1-002}{70}
正如我们将在\BookSectionRef{sec:3.5.2}{3.5.2}中回顾的那样，Cardy 态与 Ishibashi 态，是边界态的两组自然基。Cardy 态是物理的边界态，对于这些态，所有关联函数都是单值的。另一方面，“Ishibashi 态”定义为相容性条件\BookEqRef{eq:3.136}{3.136}的最简单的解，并与闭弦初级场直接一一对应。二者的区别，在于边界态的左行与右行分量如何粘合在一起：Ishibashi 态采用对角粘合，而 Cardy 态的粘合满足 Cardy 关系。

\SourceParagraph{ch02-2.3.1-003}{70}
我们的类比是：$\varepsilon_i$，$i=1,\ldots,N$，对应于 Ishibashi 态；而基 $\phi_\mu$ 则类似于一组具有确定共形权的初级场基，由关系
\begin{equation}\label{eq:2.56}
\phi_\mu\phi_\nu=N_{\mu\nu}^{\lambda}\phi_\lambda
\end{equation}
刻画，其中 $N_{\mu\nu}^{\lambda}$ 为非负整数。

\SourceParagraph{ch02-2.3.1-004}{70}
不应把这个类比推得太远，因为在拓扑理论中没有手征代数，而且应当把 $\mathcal{C}$ 的每一个元素都视为 $(T-\overline{T})=0$ 及其推广的一个解。这里没有左行模或右行模。尽管如此，利用这些公式，我们基本上恢复了 Cardy 的公式，它用特征标边界态表示 Cardy 边界态。注意，这里无需使用与模群的任何关系。

\SourceParagraph{ch02-2.3.1-005}{70}
最后再作一个简短的说明。在当前情境下，能够看到“闭弦耦合常数是开弦耦合常数的平方”这一标准关系，颇为令人满意。如果作缩放 $\theta_{\mathcal{C}}\mapsto\lambda^{-2}\theta_{\mathcal{C}}$，则 $\chi_{\mathcal{C}}=\sum_\mu\phi_\mu\phi^\mu\mapsto\lambda^2\chi_{\mathcal{C}}$。因此，可以把 $\lambda^2$ 解释为闭弦耦合常数。另一方面，$B_{\mathcal{O}}$ 中的 $\theta_i$ 的平方根说明，$B_{\mathcal{O}}\mapsto\lambda B_{\mathcal{O}}$，因此 $\lambda$ 就是开弦耦合常数。事实上，具有 $g$ 个柄和 $h$ 个孔洞的曲面的配分函数，是 $Z(\Sigma)=\theta_{\mathcal{C}}\bigl((\chi_{\mathcal{C}})^g(B_{\mathcal{O}})^h\bigr)$，所以它按照预期缩放为 $Z(\Sigma)\mapsto\lambda^{-\chi(\Sigma)}Z(\Sigma)$，其中 $\chi(\Sigma)=2-2g-h$ 是 $\Sigma$ 的 Euler 数。

\section{Landau--Ginzburg 理论}\label{sec:2.4}

\SourceParagraph{ch02-2.4-001}{71}
可以在一般的二维 $\mathcal{N}=2$ Landau--Ginzburg 理论中定义 D 膜\BookCite{ref241}。这些理论可以作拓扑扭曲，得到拓扑 Landau--Ginzburg 理论。将我们的结果应用于所得到的闭弦拓扑理论，再把由此得到的 D 膜与上述 D 膜比较，是一件值得研究的事。这里仅作一些十分初步的说明。过去几年中，沿着 Kontsevich 最初提出的建议，人们发展了一个精细的拓扑 Landau--Ginzburg 膜的范畴理论。详情见\BookSectionRef{sec:3.6.8}{3.6.8}及其中的参考文献。人们认为，这些范畴能够捕捉 D 膜的更多物理信息。在超势的所有临界点都是 Morse 临界点的情形下，存在一个到我们所构造的膜范畴的函子。

\SourceParagraph{ch02-2.4-002}{71}
回顾拓扑 LG 理论的定义。首先取一个超势 $W(X_i)$，它是手征超场 $X_1,\ldots,X_n$ 的全纯函数。当 $W$ 是具有孤立临界点的多项式时，Frobenius 代数就是 Jacobian 代数
\begin{equation}\label{eq:2.57}
\mathcal{C}=\mathbb{C}[X_1,\ldots,X_n]/(\partial_1W,\ldots,\partial_nW).
\end{equation}
Frobenius 结构由留数公式定义。例如，在单变量情形下，定义
\begin{equation}\label{eq:2.58}
\theta(\phi)\coloneqq-\operatorname{Res}_{X=\infty}\frac{\phi(X)\,\mathrm{d}X}{W'(X)}.
\end{equation}

\SourceParagraph{ch02-2.4-003}{71}
如果 $W$ 的临界点全都是 Morse 临界点，那么代数\BookEqRef{eq:2.57}{2.57}是半单的。从物理上说，Morse 临界点对应于有质量的理论，而非 Morse 临界点在红外通过重整化流趋于非平凡的二维共形场论。

\SourceParagraph{ch02-2.4-004}{71}
如果所有临界点都是 Morse 临界点，那么很容易用临界点 $p_a$ 将迹写为
\begin{equation}\label{eq:2.59}
\theta(\phi)=\sum_{\mathrm{d}W(p_a)=0}\frac{\phi(p_a)}{\det\bigl(\partial_i\partial_jW|_{p_a}\bigr)}.
\end{equation}

\SourceParagraph{ch02-2.4-005}{71}
在半单的单变量情形下，可以按如下方式构造基本幂等元。设
\begin{equation}\label{eq:2.60}
W'=\prod_{\alpha=1}^n(X-r_\alpha),
\end{equation}
并假设所有根都互不相同。那么容易验证，
\begin{equation}\label{eq:2.61}
\varepsilon_\beta\coloneqq\prod_{\alpha\ne\beta}\frac{X-r_\alpha}{r_\beta-r_\alpha}
\end{equation}
是基本幂等元。（证明时，将 $X-r_\alpha$ 写为 $(X-r_\beta)+(r_\beta-r_\alpha)$ 即可。）

\SourceParagraph{ch02-2.4-006}{72}
\BookStatementNumber{10}
\begin{example}\label{ex:2.10}
$W=\frac13t^3-qt$。当 $n=2$ 且 $q\ne0$ 时，可以显式写出
\begin{equation}\label{eq:2.62}
\begin{aligned}
\varepsilon_1&=\frac{\sqrt{q}+t}{2\sqrt{q}},\\
\varepsilon_2&=\frac{\sqrt{q}-t}{2\sqrt{q}}.
\end{aligned}
\end{equation}
注意，$\theta_1=1/(2\sqrt{q})$，$\theta_2=-1/(2\sqrt{q})$。

\SourceParagraph{ch02-2.4-007}{72}
于是，从上面的一般结果，得到 $\mathcal{O}_{aa}=\operatorname{End}(W_1)\oplus\operatorname{End}(W_2)$，以及
\begin{equation}\label{eq:2.63}
\theta_{\mathcal{O}}(\Psi)=\sqrt{\theta_1}\operatorname{Tr}(\Psi_1)+\sqrt{\theta_2}\operatorname{Tr}(\Psi_2),
\end{equation}
\begin{equation}\label{eq:2.64}
\iota^a(\Psi)=\frac{\operatorname{Tr}(\Psi_1)}{\sqrt{\theta_1}}\varepsilon_1
+\frac{\operatorname{Tr}(\Psi_2)}{\sqrt{\theta_2}}\varepsilon_2.
\end{equation}

\SourceParagraph{ch02-2.4-008}{72}
因此，一般的边界态为
\begin{equation}\label{eq:2.65}
B=w_1\frac{\varepsilon_1}{\sqrt{\theta_1}}+w_2\frac{\varepsilon_2}{\sqrt{\theta_2}},
\end{equation}
其中，在群化后的边界态格中，$w_1,w_2$ 为整数；实际向量丛的秩则为非负整数。当 $q$ 沿曲线 $q(s)=|q|e^{2\pi\mathrm{i}s}$，$0\le s\le1$，绕原点逆时针转动时，求出边界态的单值化作用很有意思。在这一操作下，$\varepsilon_1$ 与 $\varepsilon_2$ 互换。此外，沿曲线连续选取平方根，有 $\sqrt{\theta_1}=e^{-\mathrm{i}\pi s/2}/(\sqrt{2}|q|^{1/4})$，$\sqrt{\theta_2}=\mathrm{i}e^{-\mathrm{i}\pi s/2}/(\sqrt{2}|q|^{1/4})$。因此，类型为 $(w_1,w_2)$ 的膜类，演化为类型 $(w_2,-w_1)$ 的膜类。
\end{example}

\SourceParagraph{ch02-2.4-009}{72}
显然，一般 Landau--Ginzburg 理论也会出现类似现象。在超势 $W$ 的空间中，有一个余维为1的“判别式轨迹”，其上的超势具有非 Morse 临界点。绕这一轨迹作解析延拓，会置换各个 $\varepsilon_i$；对于各个 $\sqrt{\theta_i}$，则还可能在置换时产生符号变化。利用这一初步图景，可以理解文献中许多地方讨论的膜的置换与产生现象。

\SourceParagraph{ch02-2.4-010}{72}
在1988—1991年间发展起来的、研究弦在低于一维的空间中运动的理论里，我们所找到的“时空上的向量丛”，可以完全按字面理解。（综述见文献\BookCite{ref180,ref122,ref118}。）在由 $n$ 个互不相交的点组成的时空中运动的弦，可以用矩阵链或拓扑场论建模。例如，文献\BookCite{ref122,ref118}描述了后一种观点。按照这一观点，考虑与拓扑物质耦合的拓扑引力。对于由 $n$ 个点组成的时空，拓扑物质可以取为与如下超势 $W$ 关联的 $\mathcal{N}=2$ Landau--Ginzburg 理论；这个超势来自 $A_n$ 奇点的展开：
\begin{equation}\label{eq:2.66}
W=\frac{x^{n+1}}{n+1}+a_nx^n+\cdots+a_0.
\end{equation}
对于一般的 $W$，得到的是 $n$ 个时空点上的向量丛。这当然就是我们对这类时空中的膜所作的预期！

\SourceParagraph{ch02-2.4-011}{73}
还应提到，在这些最简单的弦论，也就是“极小弦论”中，近年来，人们在理解 D 膜的完整谱方面取得了很大进展，已经超出了拓扑场论截断所包含的内容。综述见文献\BookCite{ref417}。

\section{超越半单 Frobenius 代数}\label{sec:2.5}

\SourceParagraph{ch02-2.5-001}{73}
来自 $\mathcal{N}=2$ 共形场论的拓扑场论例子，例如 Landau--Ginzburg 模型和 Calabi--Yau 流形的量子上同调环，说明有必要理解 $\mathcal{C}$ 不半单时，代数条件可能具有哪些解。本节将在这个问题上取得一些局部进展，并说明，或许应当如何在更广阔的情境中看待它。

\subsection{与流形上同调有关的例子}\label{sec:2.5.1}

\SourceParagraph{ch02-2.5.1-001}{73}
分次交换 Frobenius 代数的一个自然例子，是偶数维紧定向流形 $X$ 的复系数上同调。因此，可以取 $\mathcal{C}=H^*(X,\mathbb{C})$，并以 $\theta(\phi)=\int_X\phi$ 为迹。那么，对应的 $\mathcal{O}$ 是什么？

\SourceParagraph{ch02-2.5.1-002}{73}
一个自然的猜测，是对于某个 $N>0$，取 $\mathcal{O}=\mathcal{C}\otimes\operatorname{Mat}_N(\mathbb{C})=\operatorname{Mat}_N(\mathcal{C})$，并令
\begin{equation}\label{eq:2.67}
\theta_{\mathcal{O}}(\psi)=\int_X\operatorname{Tr}(\psi).
\end{equation}
这个猜测其实是错误的，但其出错的原因颇有意思。虽然 $\mathcal{O}$ 确实是 Frobenius 代数，但映射 $\iota_*$ 唯一自然的候选，是 $\iota_*(\phi)=\phi\otimes1_N$。然而，它不满足 Cardy 条件：由伴随关系计算得到 $\iota^*(\psi)=\operatorname{Tr}(\psi)$，从而 $\iota_*\iota^*(\psi)=\operatorname{Tr}(\psi)\otimes1_N$。另一方面，还可以得到
\begin{equation}\label{eq:2.68}
\begin{aligned}
\pi(\psi)
&=\sum_{i,l,m}(-1)^{\deg\omega_i(\deg\psi+\deg\omega^i)}
(\omega_i\otimes e_{lm})\psi(\omega^i\otimes e_{ml})\\
&=\bigl(\chi(TX)\operatorname{Tr}(\psi)\bigr)\otimes1_N.
\end{aligned}
\end{equation}
这里，$\omega_i$ 与 $\omega^i$ 是 $H^*(X,\mathbb{C})$ 关于 Poincaré 内积的齐次对偶基，$e_{ml}$ 是矩阵单位，$\chi(TX)\in H^{\mathrm{top}}(X,\mathbb{C})$ 是切丛 $TX$ 的 Euler 类，由 $\chi(TX)=\sum_i(-1)^{\deg\omega_i}\omega_i\omega^i$ 给出。最后，还使用了矩阵恒等式 $\sum_{l,m}e_{lm}\psi e_{ml}=\operatorname{Tr}(\psi)\otimes1_N$。映射 $\pi$ 将正次数的形式都映为零，因此不能与 $\iota_*\iota^*$ 一致。

\SourceParagraph{ch02-2.5.1-003}{74}
可以修改这个例子，令 $\mathcal{O}$ 与 $X$ 的某个子流形关联，从而得到一个开闭弦理论。毕竟，这正是 D 膜的标准图景！在 $\mathbb{Z}$ 分次向量空间的代数范畴中进行讨论，并继续取 $\mathcal{C}=H^*(X,\mathbb{C})$，其中 $X$ 是紧、连通、定向的 $n$ 维流形，迹与前面一样，为次数 $-n$ 的 $\theta_{\mathcal{C}}(\phi)=\int_X\phi$。寻找形如 $\mathcal{O}=\operatorname{Mat}_N(\mathcal{O}_0)$ 的开弦代数，其中 $\mathcal{O}_0$ 交换。那么，$\mathcal{O}_0$ 是 Frobenius 代数。考察它具有如下几何实现的情形：$\mathcal{O}_0=H^*(Y,\mathbb{C})$，其中 $Y$ 是某个 $m$ 维紧定向流形\footnote{为涵盖这类几何实现，事实上还需要允许 $Y$ 有轨形奇点。}，而 $\iota_*\colon\mathcal{C}\to\mathcal{O}_0$ 是某个映射 $f\colon Y\to X$ 的拉回 $f^*$。

\SourceParagraph{ch02-2.5.1-004}{74}
因此，$\mathcal{O}=H^*(Y,\mathbb{C})\otimes\operatorname{Mat}_N(\mathbb{C})$，其开弦迹为
\begin{equation}\label{eq:2.69}
\theta_{\mathcal{O}}(\Psi)=\theta_o\int_Y\operatorname{Tr}(\Psi),
\end{equation}
次数为 $-m$，其中 $\theta_o$ 是非零常数。这是一个非交换 Frobenius 代数。

\SourceParagraph{ch02-2.5.1-005}{74}
伴随关系确定了 $\iota^*$：
\begin{equation}\label{eq:2.70}
\iota^*(\Psi)=\theta_o f_*\bigl(\operatorname{Tr}(\Psi)\bigr),
\end{equation}
其中，$f_*$ 是环同态 $f^*\colon H^*(X)\to H^*(Y)$ 关于 Poincaré 对偶的伴随映射。因此，$\iota^*$ 的次数为 $n-m$。另一方面，立即可见，$\pi\colon\mathcal{O}\to\mathcal{O}$ 的次数为 $m$；所以，当 Cardy 条件两端的映射并非恒零时，次数匹配要求 $n=2m$。以下考察 $n=2m$ 的情形。此时，通过对 $f$ 作一般的小扰动，可以使它成为 $Y$ 到 $X$ 的浸入。下面用 $X$ 内的管状邻域来描述伴随映射时，进一步假定 $f$ 为嵌入。于是，$f_*$ 可以更明确地写为
\begin{equation}\label{eq:2.71}
f_*(\psi)=\operatorname{pr}^*(\psi)\wedge\Phi_{N_Y},
\end{equation}
其中，$\operatorname{pr}\colon N_Y\to Y$ 是法丛的投影，将法丛的零截面附近与 $Y$ 在 $X$ 中的管状邻域等同；$\Phi_{N_Y}$ 是该丛的 Thom 类，在管状邻域内具有紧支集，在邻域之外延拓为零，代表 $Y$ 在 $X$ 中的上同调类。容易得到
\begin{equation}\label{eq:2.72}
\iota_*\iota^*(\Psi)=\theta_o\chi(N_Y)\wedge\operatorname{Tr}(\Psi)\otimes1_N,
\end{equation}
其中，$\chi(N_Y)$ 是嵌入 $Y\hookrightarrow X$ 的法丛的 Euler 类，也就是 $Y$ 在 $X$ 中的同调自交类。

\SourceParagraph{ch02-2.5.1-006}{74}
另一方面，由\BookEqRef{eq:2.68}{2.68}与\BookEqRef{eq:2.70}{2.70}，有
\begin{equation}\label{eq:2.73}
\pi(\Psi)=\frac{1}{\theta_o}\operatorname{Tr}(\Psi)\chi(TY)\otimes1_N,
\end{equation}
其中，$\chi(TY)\in H^{\mathrm{top}}(Y,\mathbb{C})$ 是切丛 $TY$ 的 Euler 类，其积分是 $Y$ 的 Euler 数。

\SourceParagraph{ch02-2.5.1-007}{75}
显然，只要选取非零的 $\theta_o$，使 $\chi(TY)=\theta_o^2\chi(N_Y)$，Cardy 条件就成立。只要 $\chi(N_Y)$（即 $Y$ 在 $X$ 中的自交数）非零，这样的选择总是可能的。如果 $Y$ 是辛流形 $X$ 的 Lagrange 子流形，同样可以作出这一选择，因为此时 $N_Y\cong TY$。边界态为 $B=\theta_oN\Phi_{N_Y}$。

\SourceParagraph{ch02-2.5.1-008}{75}
这个讨论的一个直接结果是：如果从开弦代数 $\mathcal{O}=H^*(\mathbb{CP}^2,\mathbb{C})$ 出发，那么将 $Y$ 分别视为 $X=\mathbb{CP}^4$ 或 $X'=\mathbb{HP}^2$ 的子流形，就很容易找到两个与它相容的不同闭弦代数。

\SourceParagraph{ch02-2.5.1-009}{75}
遗憾的是，对于 $\mathcal{C}=H^*(X,\mathbb{C})$，我们还不知道如何描述其边界条件范畴。不过，无论如何，似乎都必须考虑上链层面的理论，才能得到有意义的结果。接下来转向这一情形。

\subsection{Chas--Sullivan 理论}\label{sec:2.5.2}

\SourceParagraph{ch02-2.5.2-001}{75}
Chas 与 Sullivan\BookCite{ref091}给出了一个有趣的例子：在上链层面，它的结构比我们的开闭弦理论略弱一些，却可以帮助说明使用上链理论的意义。从紧定向流形 $X$ 出发，假定它连通且单连通。可以定义一个范畴 $\mathcal{B}$，其对象是 $X$ 的定向子流形，从 $Y$ 到 $Z$ 的态射向量空间，在本小节按路径的方向记为 $\mathcal{O}_{YZ}=\operatorname{Ext}_{H^*(X)}^*\bigl(H^*(Y),H^*(Z)\bigr)$；与往常一样，上同调取复系数，并通过限制，将 $H^*(Y)$ 和 $H^*(Z)$ 视为 $H^*(X)$ 模。态射的复合，由 Ext 群的 Yoneda 复合给出。不过，按照这一定义，$\mathcal{O}_{YZ}$ 并不一定是 $\mathcal{O}_{ZY}$ 的对偶。（要看到这一点，只需考虑 $Y=Z$ 是 $X$ 中的一个点，而 $X$ 是若干奇数维球面的乘积的情形；此时，$\mathcal{O}_{YZ}$ 是一个对称代数，作为向量空间不是自对偶的。）

\SourceParagraph{ch02-2.5.2-002}{75}
可以作出更好的定义：将“态射”的上链复形 $\widehat{\mathcal{O}}_{YZ}$ 定义为
\begin{equation}\label{eq:2.74}
\widehat{\mathcal{O}}_{YZ}=\mathcal{B}_{\Omega(X)}\bigl(\Omega(Y),\Omega(Z)\bigr),
\end{equation}
其中，$\Omega(X)$ 表示流形 $X$ 通常的 de Rham 复形；对于微分分次代数 $A$ 以及微分分次 $A$ 模 $B,C$，$\mathcal{B}_A(B,C)$ 是通常的余 bar 型分解的总复形：
\begin{equation}\label{eq:2.75}
\operatorname{Hom}(B,C)\longrightarrow\operatorname{Hom}(A\otimes B,C)
\longrightarrow\operatorname{Hom}(A\otimes A\otimes B,C)\longrightarrow\cdots.
\end{equation}
在内部次数都为零时，其 bar 方向的微分为
\begin{equation}\label{eq:2.76}
\begin{aligned}
(df)(a_1\otimes\cdots\otimes a_k\otimes b)
&=a_1f(a_2\otimes\cdots\otimes a_k\otimes b)\\
&\quad+\sum_{i=1}^{k-1}(-1)^if(a_1\otimes\cdots\otimes a_ia_{i+1}\otimes\cdots\otimes a_k\otimes b)\\
&\quad+(-1)^kf(a_1\otimes\cdots\otimes a_{k-1}\otimes a_kb).
\end{aligned}
\end{equation}
在微分分次情形下，再加入内部微分，并使用相应的 Koszul 符号。这个总复形的上同调是 $\operatorname{Ext}_A(B,C)$。

\SourceParagraph{ch02-2.5.2-003}{76}
它与 $\mathcal{O}_{YZ}=\operatorname{Ext}_{H^*(X)}^*\bigl(H^*(Y),H^*(Z)\bigr)$ 不同，但二者通过一个谱序列关联；该谱序列的 $E_2$ 项为 $\mathcal{O}_{YZ}$，并收敛到 $H^*(\widehat{\mathcal{O}}_{YZ})=\operatorname{Ext}_{\Omega(X)}\bigl(\Omega(Y),\Omega(Z)\bigr)$。更重要的是，$H^*(\widehat{\mathcal{O}}_{YZ})$ 是路径空间 $\mathcal{P}_{YZ}$ 的同调，这个空间由 $X$ 中从 $Y$ 出发、在 $Z$ 终止的路径组成。在与弦拓扑比较时，采用同调式的总次数标号，它与通常 Ext 上链总次数的标号相反；准确地说，在这一约定下，$H^p(\widehat{\mathcal{O}}_{YZ})\cong H_{p+d_Z}(\mathcal{P}_{YZ})$，其中 $d_Z$ 是 $Z$ 的维数。在上链复形上，Yoneda 复合在上链同伦意义下满足结合律，并定义一个 $A_\infty$ 范畴 $\widehat{\mathcal{B}}$ 的结构。相应的同调群复合
\begin{equation}\label{eq:2.77}
H_i(\mathcal{P}_{YZ})\times H_j(\mathcal{P}_{ZW})\longrightarrow
H_{i+j-d_Z}(\mathcal{P}_{YW}),
\end{equation}
是两个映射的复合：先取与余维 $d_Z$ 的子流形 $\mathcal{M}\hookrightarrow\mathcal{P}_{YZ}\times\mathcal{P}_{ZW}$ 关联的 Gysin 映射，其中 $\mathcal{M}$ 由可连接的路径对组成，再取路径连接映射 $\mathcal{M}\to\mathcal{P}_{YW}$。

\SourceParagraph{ch02-2.5.2-004}{76}
尝试为这个 $A_\infty$ 范畴配上一个“闭弦”上链代数 $\mathcal{C}$。容易看出，$\mathcal{B}$ 的恒等函子的自同态代数，也就是\BookSectionRef{sec:2.2}{2.2}中记作 $\mathcal{E}$ 的代数，正是上同调代数 $H^*(X)$。我们已在\BookSectionRef{sec:2.2.3}{2.2.3}提到，这就是零次 Hochschild 上同调 $\operatorname{HH}^0(\mathcal{B})$。

\SourceParagraph{ch02-2.5.2-005}{76}
线性范畴 $\mathcal{B}$ 的 Hochschild 上同调的定义，已在\BookSectionRef{sec:2.2}{2.2}末尾给出。事实上，对于 $\widehat{\mathcal{B}}$ 这样的 $A_\infty$ 范畴，Hochschild 复形的定义仍然有意义，它是闭弦代数 $\mathcal{C}$ 的一个候选。

\SourceParagraph{ch02-2.5.2-006}{76}
在当前情形中，$\mathcal{C}$ 等价于微分分次代数 $\Omega(X)$ 通常的 Hochschild 复形，其上同调是自由环路空间 $\mathcal{L}X$ 的同调，并按照上述同调式次数约定，将次数向下平移 $X$ 的维数 $d_X$\footnote{因此，代数的单位元位于 $H^0(\mathcal{C})$ 中，它是 $X$ 的基本类；将 $X$ 的点视为 $\mathcal{L}X$ 中的常值环路，就可以把这个基本类看作 $H_{d_X}(\mathcal{L}X)$ 的元素。}，所以 $H^i(\mathcal{C})$ 可能非零的范围为 $-d_X\le i<\infty$。Chas 与 Sullivan 正是在当前这一情境下引入这个代数的：他们试图在经典代数拓扑的框架中，重现弦论的结构。有一个映射 $H^i(X)\to H^{-i}(\mathcal{C})$，将 $X$ 的普通上同调环嵌入 Chas--Sullivan 环；还存在环同态 $H^i(\mathcal{C})\to H_i(\mathcal{L}_0X)$，映入以 $X$ 中任意选定点为基点的基点环路空间 $\mathcal{L}_0X$ 的 Pontryagin 环。

\SourceParagraph{ch02-2.5.2-007}{76}
前面在\BookSectionRef{sec:2.2.3}{2.2.3}提到的 $\mathcal{C}$ 的另一个候选，是代数 $\Omega(X)$ 的循环上同调；众所周知\BookCite{ref176}，它对应于自由环路空间 $\mathcal{L}X$ 关于自然圆周作用的等变同调。与非等变同调相比，这也许是一个更好的选择。

\SourceParagraph{ch02-2.5.2-008}{76}
不过，我们得到的结构还不是上链层面的开闭弦理论，因为没有能够在 $H^*(\widehat{\mathcal{O}}_{YZ})$ 上诱导内积的迹映射。尝试定义与配边对应的算符时，会发现，只有当配边的每个连通分支都具有非空的出边界时，才能做到这一点。（定义在这一较小范畴上的理论，通常称为非紧理论。）按照我们的定义，与 Chas--Sullivan 理论最接近的，是为辛流形 $X$ 定义的所谓“A 模型”。那里，$A_\infty$ 范畴就是 Fukaya 范畴（更多细节见\BookChapterRef{8}），其对象是 $X$ 中配备带联络的丛的 Lagrange 子流形；从 $Y$ 到 $Z$ 的态射上链复形，是计算路径空间 $\mathcal{P}_{YZ}$ 的“半无限”上同调的 Floer 复形。在良好的情形下，这个 Floer 复形的上同调具有一组由 $Y$ 与 $Z$ 的交点标记的向量空间基，而相应闭弦复形的上同调，就是 $X$ 的普通上同调。从我们的观点看，Floer 理论的关键特点，是它实现了无限维流形 $\mathcal{L}X$ 上的 Poincaré 对偶。

\subsection{关于 B 模型的说明}\label{sec:2.5.3}

\SourceParagraph{ch02-2.5.3-001}{77}
设 $X$ 是复维数为 $d$ 的复簇，并给定其典范丛的一个平凡化。也就是说，假设存在一个处处非零的全纯 $d$ 形式 $\Omega$。B 模型\BookCite{ref468}是来自 $X$ 的 $\mathcal{N}=2$ 超对称 $\sigma$ 模型的一个 $\mathbb{Z}$ 分次拓扑场论。这个理论的自然边界条件，由 $X$ 上的全纯向量丛给出。

\SourceParagraph{ch02-2.5.3-002}{77}
全纯向量丛范畴不是 Frobenius 范畴。不过，与 $X$ 关联着一个十分自然的 $\mathbb{Z}$ 分次 Frobenius 范畴 $\mathcal{V}_X$：其对象是 $X$ 上的向量丛，而从 $E$ 到 $F$ 的态射空间，继续按照源对象在前的记号，取为
\begin{equation}\label{eq:2.78}
\mathcal{O}_{EF}=\operatorname{Ext}_X^*(E,F)=H^{0,*}(X,E^*\otimes F).
\end{equation}
次数为 $-d$ 的迹 $\theta_E\colon\mathcal{O}_{EE}\to\mathbb{C}$，定义为
\begin{equation}\label{eq:2.79}
\theta_E(\Psi)=\int_X\operatorname{Tr}(\Psi)\wedge\Omega.
\end{equation}
由 Serre 对偶，这个迹是非退化的；但该范畴仍然不是半单的。事实上，对于 $i>0$，群 $\operatorname{Ext}^i$ 的非零性，正是范畴非半单性的表现。（$\operatorname{Ext}_X^1(E,F)$ 中的一个非零元素，对应于不分裂的正合列 $0\to F\to G\to E\to0$；也就是一个向量丛 $G$，带有一个不存在补丛的子丛 $F$。）

\SourceParagraph{ch02-2.5.3-003}{77}
$\mathcal{V}_X$ 的恒等函子有哪些自同态？乘以 $H^{0,*}(X)$ 的任意一个元素，显然都定义了这样的自同态。$X$ 上的全纯向量场 $\xi$ 也定义了一个次数为1的自同态，因为每个丛 $E$ 都有一个“Atiyah 类”\footnote{它对应于向量丛的扩张 $0\to E\otimes T_X^*\to J^1E\to E\to0$，其中 $J^1E$ 是 $E$ 的全纯截面的1阶喷射丛。}$a_E\in\operatorname{Ext}_X^1(E,E\otimes T_X^*)$，代表其曲率；将它与 $\xi$ 缩并，就得到 $e_\xi=\iota_\xi a_E\in\operatorname{Ext}_X^1(E,E)$。更一般地，一个类
\[
\eta\in H^{0,q}\bigl(X,\textstyle\bigwedge^pT_X\bigr)
=\operatorname{Ext}_X^q\bigl(\textstyle\bigwedge^pT_X^*,\mathcal{O}_X\bigr)
\]
可以与 $(a_E)^p\in\operatorname{Ext}_X^p\bigl(E,E\otimes(T_X^*)^{\otimes p}\bigr)$ 缩并，得到
\[
e_\eta=\iota_\eta(a_E)^p\in\operatorname{Ext}_X^{p+q}(E,E).
\]

\SourceParagraph{ch02-2.5.3-004}{78}
Witten 在文献\BookCite{ref468}中证明，$H^{0,*}(X,\bigwedge^*T_X)$ 的确是 B 模型的闭弦代数。为了在我们的框架中理解这一点，必须再次转到上链层面的理论；上述 Ext 群是这个理论的上同调。一种很好的办法，是把全纯向量丛 $E$ 替换为它的 $\overline{\partial}$ 复形 $\widehat{E}=\Omega^{0,*}(X,E)$；它是微分分次代数 $A=\Omega^{0,*}(X)$ 的一个微分分次模。然后，将 $\widehat{\mathcal{O}}_{EF}$ 定义为上链复形 $\operatorname{Hom}_A(\widehat{E},\widehat{F})$，其上同调群为 $\operatorname{Ext}_X^*(E,F)$。

\SourceParagraph{ch02-2.5.3-005}{78}
既然这样做，自然也就应当允许更大的一类对象：底层由有界、有限生成投射的分次 $A^0$ 模诱导，并配备相容微分的 $A$ 模。$X$ 上的任意凝聚层 $E$ 都定义一个这样的模：先用向量丛复形 $\mathcal{E}$ 分解 $E$，必要时将局部分解以导出方式粘合，再取双复形 $\widehat{\mathcal{E}}$ 的总复形。所得到的扩大后的范畴，在取同伦范畴后，本质上就是 $X$ 上凝聚层范畴的有界导出范畴；这也是\BookChapterRef{4}的主要主题之一。

\SourceParagraph{ch02-2.5.3-006}{78}
在这一框架中，与上面的拓扑例子一样，可以得到恒等函子的导出自同态由整个微分分次范畴的 Hochschild 上链复形 $\widehat{\mathcal{C}}$ 给出。对于一个局部代数模型，Hochschild 上链的形状是
\[
\bigl\{A\longrightarrow\operatorname{Hom}_{\mathbb{C}}(A,A)
\longrightarrow\operatorname{Hom}_{\mathbb{C}}(A\otimes A,A)\longrightarrow\cdots\bigr\}.
\]
在整体几何中，需要使用局部多重微分上链的层及其导出整体截面，或者等价地，直接使用上述整个微分分次范畴的 Hochschild 复形。其上同调为 $H^*(X,\bigwedge^*T_X)$。不过，要理解 $\widehat{\mathcal{C}}$ 上的迹映射，以及伴随映射 $\iota_E$ 与 $\iota^E$，仍有工作需要完成。我们认为，文献尚未充分阐明这一点。关于这个问题的一些进展，见文献\BookCite{ref348,ref081,ref082}；更近期的进展见文献\BookCite{ref080}。

\section{等变二维拓扑开闭弦理论}\label{sec:2.6}

\SourceParagraph{ch02-2.6-001}{78}
弦论中一个重要的构造，是“轨形”构造。抽象地说，当闭弦背景有一个自同构群 $G$，并且相应对称性可以规范化时，就可以进行这一构造。定义轨形理论分为两步。首先，引入“外部”规范场，将理论扩展；这些规范场就是世界面上的 $G$ 丛及其联络。其次，对所有可能的 $G$ 丛及其联络求和，构造一个新理论。

\SourceParagraph{ch02-2.6-002}{78}
先仔细描述形成轨形理论的第一步。在 $G$ 为有限群的情形下，第二步，也就是对 $G$ 丛求和，就非常容易。

\subsection{等变闭弦理论}\label{sec:2.6.1}

\SourceParagraph{ch02-2.6.1-001}{79}
先作一些一般的说明。在 $d$ 维拓扑场论中，从一个范畴 $\mathcal{S}$ 出发，其对象是定向的 $(d-1)$ 维流形，态射是定向配边。如果群 $G$ 作用在与每个 $(d-1)$ 维流形关联的向量空间上，并且与每个配边关联的线性算符都是 $G$ 等变映射，物理学家就说，这个理论具有全局对称群 $G$。有了这样的“全局”对称群 $G$，就可以考察这个对称性是否能够“规范化”，也就是，$G$ 的元素是否可以在某种意义下，“独立地”作用于时空的每个点。从数学上说，“规范化”过程有一个非常优美的描述：它将场论函子从范畴 $\mathcal{S}$ 扩展到范畴 $\mathcal{S}_G$，后者的对象是配备主 $G$ 丛的 $(d-1)$ 维流形，态射则是带有 $G$ 丛的配边\footnote{这里假设群 $G$ 是离散的。如果 $G$ 是 Lie 群，则应将 $\mathcal{S}_G$ 定义为配备带联络的主 $G$ 丛的流形所构成的范畴。}。为每个 $(d-1)$ 维流形 $S$ 配上平凡 $G$ 丛 $S\times G$，就可以把 $\mathcal{S}$ 视为 $\mathcal{S}_G$ 的子范畴。在 $\mathcal{S}_G$ 中，平凡丛 $S\times G$ 的自同构群包含 $G$，因此，在这样扩展后的理论中，$G$ 作用于态空间 $\mathcal{H}(S)$；这个作用应当就是原来的“全局”$G$ 作用。不过，对于 $S$ 上的每个 $G$ 丛 $P$，扩展后的理论都有一个态空间 $\mathcal{H}(S,P)$。如果 $P$ 非平凡，就把 $\mathcal{H}(S,P)$ 称为理论的“扭曲扇区”。当 $d=2$、$S=S^1$ 时，可以用乘以 $g$ 的映射，将 $[0,2\pi]\times G$ 的两端粘合，得到丛 $P_g\to S^1$。（按照定义，$P_g$ 在 $S^1$ 的基点处的纤维就是群 $G$。）任意丛都同构于其中某个丛；$P_g$ 与 $P_{g'}$ 同构，当且仅当 $g'$ 与 $g$ 共轭。不过要注意，态空间依赖于丛本身，而不仅仅依赖于它的同构类；因此，得到的扭曲扇区态空间 $\mathcal{C}_g=\mathcal{H}(S,P_g)$，由群元素 $g$ 标记，而不是由共轭类标记。

\SourceParagraph{ch02-2.6.1-002}{79}
我们把定义在范畴 $\mathcal{S}_G$ 上的理论，称为 $G$ 等变 TFT。必须区分等变理论与下面将描述的相应“规范化理论”。在物理上，等变理论通过与不参与动力学的背景规范场耦合得到；规范化理论则通过在路径积分中对这些规范场“求和”得到。

\SourceParagraph{ch02-2.6.1-003}{79}
另一个等价的观点，在二维情形下尤其有用：将 $\mathcal{S}_G$ 视为如下范畴，其对象是配备映射 $p\colon S\to BG$ 的定向 $(d-1)$ 维流形 $S$，其中 $BG$ 是 $G$ 的分类空间。在这一观点下，空间 $\operatorname{Map}(S,BG)$ 上有一个丛，在 $p$ 处的纤维为 $\mathcal{H}_p$。说 $\mathcal{H}_p$ 仅依赖于 $S$ 上的 $G$ 丛 $p^*EG$，也就是通过 $p$ 从 $BG$ 上的普遍 $G$ 丛 $EG$ 拉回的丛，等价于说，$\operatorname{Map}(S,BG)$ 上的这个丛配备了平坦联络，可以通过平行移动，识别同一连通分支内各点处的纤维。因为，丛同构 $p_0^*EG\to p_1^*EG$ 的集合，与从 $p_0$ 到 $p_1$ 的路径的同伦类集合相同。当 $S=S^1$ 时，映射空间的连通分支对应于 $G$ 的共轭类：每个丛 $P_g$ 对应于映射空间中一个特定的点 $p_g$，而群元素 $h$ 则确定一条从 $p_g$ 到 $p_{hgh^{-1}}$ 的特定路径。

\SourceParagraph{ch02-2.6.1-004}{80}
第二种观点清楚地说明，$G$ 等变拓扑场论，是 Turaev\BookCite{ref452}意义下的“同伦拓扑场论”的例子。我们将使用他的两个主要结果：其一，是对“二维 TFT 就是交换 Frobenius 代数”这一定理的一个漂亮推广；其二，是对给定的半单 TFT 的全局 $G$ 对称性进行规范化的各种方式的分类。现在简要回顾他的工作。

\begin{figure}[H]
\centering
\BookImage{01_Numbered_Figures/Chapter_02/Ch02_Fig15_PDFp080_Bookp070.pdf}
\caption{$G$ 等变闭弦理论中乘积的定义。实心点是 $S^1$ 的基点。为了无歧义地指定态射，需要沿一组曲线给出相容的平行移动，使这些曲线的补集由单连通的部分组成。这些平行移动，总是在按定义纤维为 $G$ 的点之间的路径上进行。因此，乘积通常不交换。必须约定复合曲线上的平行移动采用左路径排序还是右路径排序；这里取 $h(\gamma_1\circ\gamma_2)=h(\gamma_1)\cdot h(\gamma_2)$。}\label{fig:2.15}
\end{figure}

\SourceParagraph{ch02-2.6.1-005}{80}
一个 $G$ 等变 TFT，为每个元素 $g\in G$ 给出一个向量空间 $\mathcal{C}_g$，它与配备丛 $P_g$ 的圆周关联，这个丛绕圆周的平行移动为 $g$。通常的裤面配边，配上一个自然的 $G$ 丛，使其在两个入圆周上的限制分别为 $P_{g_1}$ 与 $P_{g_2}$，而在出圆周上的限制为 $P_{g_1g_2}$，就诱导乘积
\begin{equation}\label{eq:2.80}
\mathcal{C}_{g_1}\otimes\mathcal{C}_{g_2}\longrightarrow\mathcal{C}_{g_1g_2},
\end{equation}
从而使 $\mathcal{C}\coloneqq\bigoplus_{g\in G}\mathcal{C}_g$ 成为 $G$ 分次代数，如\BookRef{fig:2.15}{图15}所示。

\begin{figure}[H]
\centering
\BookImage{01_Numbered_Figures/Chapter_02/Ch02_Fig16_PDFp081_Bookp071.pdf}
\caption{(a) $\alpha_h$ 对态 $\phi\in\mathcal{C}_g$ 的作用。(b) 也可以用圆柱表示这一作用。}\label{fig:2.16}
\end{figure}

\SourceParagraph{ch02-2.6.1-006}{81}
与通常的情形一样，一个具有入圆周的圆盘图示，定义了迹 $\theta\colon\mathcal{C}_1\to\mathbb{C}$。注意，绕圆盘边界的平行移动必须为1。作出“圆柱对应于恒等算符”这一标准假设，就得到非退化配对 $\mathcal{C}_g\otimes\mathcal{C}_{g^{-1}}\to\mathbb{C}$。

\SourceParagraph{ch02-2.6.1-007}{81}
等变理论中的一个新要素，是 $G$ 作为自同构群作用于 $\mathcal{C}$。也就是说，存在同态 $\alpha\colon G\to\operatorname{Aut}(\mathcal{C})$，使
\begin{equation}\label{eq:2.81}
\alpha_h\colon\mathcal{C}_g\longrightarrow\mathcal{C}_{hgh^{-1}}.
\end{equation}
在图示中，$\alpha_h$ 由\BookRef{fig:2.16}{图16}中的曲面定义。

\SourceParagraph{ch02-2.6.1-008}{81}
记录 $\alpha$ 的一些性质。首先，如果 $\phi\in\mathcal{C}_h$，那么 $\alpha_h(\phi)=\phi$。原因由\BookRef{fig:2.17}{图17}说明。

\begin{figure}[H]
\centering
\BookImage{01_Numbered_Figures/Chapter_02/Ch02_Fig17_PDFp082_Bookp072.pdf}
\caption{如果沿路径 $\mathcal{P}_2$ 的平行移动为 $h$，那么沿路径 $\mathcal{P}_1$ 的平行移动为1。不过，绕内圆周的一个 Dehn 扭转，将 $\mathcal{P}_1$ 映为 $\mathcal{P}_2$。因此，对于 $\phi\in\mathcal{C}_h$，有 $\alpha_h(\phi)=\alpha_1(\phi)=\phi$。}\label{fig:2.17}
\end{figure}

\SourceParagraph{ch02-2.6.1-009}{81}
其次，虽然 $\mathcal{C}$ 通常不交换，但它在如下意义下“扭曲交换”：如果 $\phi_1\in\mathcal{C}_{g_1}$，$\phi_2\in\mathcal{C}_{g_2}$，那么
\begin{equation}\label{eq:2.82}
\alpha_{g_2}(\phi_1)\phi_2=\phi_2\phi_1.
\end{equation}
\BookRef{fig:2.18}{图18}说明了这一条件的必要性。

\begin{figure}[H]
\centering
\BookImage{01_Numbered_Figures/Chapter_02/Ch02_Fig18_PDFp082_Bookp072.pdf}
\caption{扭曲中心性的图示。}\label{fig:2.18}
\end{figure}

\SourceParagraph{ch02-2.6.1-010}{81}
所需的最后一个性质稍复杂一些。$\mathcal{C}_g$ 的恒等映射的迹，等于理论在环面上的配分函数，其中所配备的丛沿两个生成圆周的平行移动为 $(g,1)$。沿另一个方向切开环面，可以看出，它就是 $\alpha_g$ 在 $\mathcal{C}_1$ 上的迹。类似地，考虑平行移动为 $(g,h)$ 的丛，其中 $g,h$ 是 $G$ 中两个交换的元素，就可以看出，$\alpha_g$ 在 $\mathcal{C}_h$ 上的迹，等于 $\alpha_h$ 在 $\mathcal{C}_{g^{-1}}$ 上的迹。不过，还需要加强这一性质。即使 $g$ 与 $h$ 不交换，也可以在有一个孔洞的环面上，构造平行移动为 $(g,h)$ 的丛；绕孔洞的平行移动为 $c=hgh^{-1}g^{-1}$。沿环面的任一生成圆周切开，就得到一个配边算符，或者从 $\mathcal{C}_c\otimes\mathcal{C}_h$ 到 $\mathcal{C}_h$，或者从 $\mathcal{C}_{g^{-1}}\otimes\mathcal{C}_c$ 到 $\mathcal{C}_{g^{-1}}$。对于 $\psi\in\mathcal{C}_{hgh^{-1}g^{-1}}$，引入两个线性变换 $L_\psi,R_\psi$，分别对应于左乘与右乘 $\psi$。一方面，$L_\psi\alpha_g\colon\phi\mapsto\psi\alpha_g(\phi)$ 是映射 $\mathcal{C}_h\to\mathcal{C}_h$；另一方面，$R_\psi\alpha_h\colon\phi\mapsto\alpha_h(\phi)\psi$ 是映射 $\mathcal{C}_{g^{-1}}\to\mathcal{C}_{g^{-1}}$。最后的缝合条件要求，这两个自同态具有相同的迹：
\begin{equation}\label{eq:2.83}
\operatorname{Tr}_{\mathcal{C}_h}(L_\psi\alpha_g)
=\operatorname{Tr}_{\mathcal{C}_{g^{-1}}}(R_\psi\alpha_h).
\end{equation}

\begin{figure}[H]
\centering
\BookImage{01_Numbered_Figures/Chapter_02/Ch02_Fig19_PDFp083_Bookp073.pdf}
\caption{将 (a) 的左边变形为时空演化图示，就得到 (b)，其值为 $\operatorname{Tr}_{\mathcal{C}_h}(L_\psi\alpha_g)$。类似地，将 (a) 的右边变形，所得图示的值为 $\operatorname{Tr}_{\mathcal{C}_{g^{-1}}}(R_\psi\alpha_h)$。}\label{fig:2.19}
\end{figure}

\SourceParagraph{ch02-2.6.1-011}{83}
仔细思考\BookRef{fig:2.19}{图19}中的图示，就可以推知这个条件的原因。

\begin{figure}[H]
\centering
\BookImage{01_Numbered_Figures/Chapter_02/Ch02_Fig20_PDFp083_Bookp073.pdf}
\caption{一个比 Turaev 的环面公理更简单的公理。}\label{fig:2.20}
\end{figure}

\SourceParagraph{ch02-2.6.1-012}{83}
Turaev 将方程\BookEqRef{eq:2.83}{2.83}作为公理之一。不过，可以把它改写成一种更便于我们使用的形式。设 $\Delta_g\in\mathcal{C}_g\otimes\mathcal{C}_{g^{-1}}$ 是“对偶”元素，对应于将两端都看作出边界的 $(S^1,P_g)$ 的恒等配边。于是，$\Delta_g=\sum_i\xi_i\otimes\xi^i$，其中 $\xi_i$ 与 $\xi^i$ 分别遍历 $\mathcal{C}_g$ 与 $\mathcal{C}_{g^{-1}}$ 的对偶基。再记
\[
\Delta_h=\sum_i\eta_i\otimes\eta^i\in\mathcal{C}_h\otimes\mathcal{C}_{h^{-1}}.
\]
那么，容易看出，\BookEqRef{eq:2.83}{2.83}等价于
\begin{equation}\label{eq:2.84}
\sum_i\alpha_h(\xi_i)\xi^i=\sum_i\eta_i\alpha_g(\eta^i),
\end{equation}
其两端都是 $\mathcal{C}_{hgh^{-1}g^{-1}}$ 中的元素。\BookRef{fig:2.20}{图20}中的同构配边说明了这一方程。

\SourceParagraph{ch02-2.6.1-013}{84}
总而言之，$G$ 等变二维拓扑场论的缝合定理，可以表述为下面的定理。

\SourceParagraph{ch02-2.6.1-014}{84}
\BookStatementNumber{11}
\begin{theorem}\label{thm:2.11}
\BookCite{ref452} 给定一个二维 $G$ 等变拓扑场论，等价于给定一个各分量有限维的 $G$ 分次结合含幺代数 $\mathcal{C}=\bigoplus_g\mathcal{C}_g$，以及一个群同态 $\alpha\colon G\to\operatorname{Aut}(\mathcal{C})$，满足分次相容性\BookEqRef{eq:2.81}{2.81}，并使以下条件成立：
\begin{enumerate}[label=(\arabic*)]
\item\label{ax:2.11.1} 配备一个 $G$ 不变迹 $\theta\colon\mathcal{C}_1\to\mathbb{C}$，诱导非退化配对 $\mathcal{C}_g\otimes\mathcal{C}_{g^{-1}}\to\mathbb{C}$。
\item\label{ax:2.11.2} $\alpha_h$ 在 $\mathcal{C}_h$ 上的限制是恒等映射。
\item\label{ax:2.11.3} 对于任意 $\phi\in\mathcal{C}_g$、$\phi'\in\mathcal{C}_h$，有 $\alpha_h(\phi)\phi'=\phi'\phi$。
\item\label{ax:2.11.4} 对于任意 $g,h\in G$，有
\begin{equation}\label{eq:2.85}
\sum_i\alpha_h(\xi_i)\xi^i=\sum_i\eta_i\alpha_g(\eta^i)\in\mathcal{C}_{hgh^{-1}g^{-1}},
\end{equation}
其中，$\Delta_g=\sum_i\xi_i\otimes\xi^i\in\mathcal{C}_g\otimes\mathcal{C}_{g^{-1}}$，$\Delta_h=\sum_i\eta_i\otimes\eta^i\in\mathcal{C}_h\otimes\mathcal{C}_{h^{-1}}$，与前面一样。
\end{enumerate}
\end{theorem}

\SourceParagraph{ch02-2.6.1-015}{84}
\BookStatementNumber{12}
\begin{note}[2.12]\label{note:2.12}
\begin{enumerate}[label=(\arabic*)]
\item 我们将在\BookSectionRef{sec:2.7.3}{2.7.3}中证明缝合定理。
\item 注意，Turaev 将上述代数称为交叉 $G$ Frobenius 代数，但它并不是 $C^*$ 代数意义下的交叉积代数，见下文。我们把满足\BookRef{thm:2.11}{定理2.11}条件的代数，称为 Turaev 代数。
\item \BookRef{ax:2.11.1}{公理1}与\BookRef{ax:2.11.3}{公理3}，在非等变理论中有对应物；而\BookRef{ax:2.11.2}{公理2}与\BookRef{ax:2.11.4}{公理4}则是新要素。
\end{enumerate}
\end{note}

\subsection{轨形理论}\label{sec:2.6.2}

\SourceParagraph{ch02-2.6.2-001}{84}
在继续之前，先描述如何从 Turaev 代数得到轨形理论。

\SourceParagraph{ch02-2.6.2-002}{84}
回到\BookSectionRef{sec:2.6.1}{2.6.1}开头的一般讨论，那里概述了等变理论的定义。粗略地说，规范化理论通过对规范场求和，从等变理论得到。更准确地说，对于这里考虑的有限群 $G$，规范化理论给 $(d-1)$ 维流形 $S$ 指定的态空间，由“波函数”$\psi$ 组成；它为 $S$ 上的每个 $G$ 丛 $P$ 指定一个等变理论态空间 $\mathcal{H}_{S,P}$ 中的元素 $\psi_P$。映射 $\psi$ 必须是“自然的”：如果 $\theta\colon P\to P'$ 是丛同构，那么所诱导的同构 $\mathcal{H}_{S,P}\to\mathcal{H}_{S,P'}$ 必须把 $\psi_P$ 映为 $\psi_{P'}$。这通常称为“Gauss 定律”。在二维情形中，Gauss 定律意味着，圆周的态空间 $\mathcal{C}_{\mathrm{orb}}$ 是 Turaev 代数 $\mathcal{C}=\bigoplus_g\mathcal{C}_g$ 的 $G$ 不变部分。换言之，作为向量空间，
\begin{equation}\label{eq:2.86}
\mathcal{C}_{\mathrm{orb}}=\bigoplus_{[g]}(\mathcal{C}_g)^{Z_g},
\end{equation}
其中，$g$ 遍历 $G$ 的共轭类的一组代表元，取 $\mathcal{C}_g$ 关于 $g$ 在 $G$ 中的中心化子 $Z_g$ 的不变部分。当 $G$ 非 Abel 时，代数 $\mathcal{C}_{\mathrm{orb}}$ 通常不再继承原来的 $G$ 分次。需要验证，$\mathcal{C}_{\mathrm{orb}}$ 中的乘积，就是 $\mathcal{C}$ 中乘积的限制。迹 $\mathcal{C}_{\mathrm{orb}}\to\mathbb{C}$，是迹 $\mathcal{C}\to\mathbb{C}$ 的限制；后者在 $\mathcal{C}_1$ 上等于给定的迹，在 $g\ne1$ 的 $\mathcal{C}_g$ 上为零。于是，$\mathcal{C}_{\mathrm{orb}}$ 是一个编码轨形理论的交换 Frobenius 代数。

\subsection{闭弦 \texorpdfstring{$G$}{G} 等变缝合条件的解}\label{sec:2.6.3}

\SourceParagraph{ch02-2.6.3-001}{85}
找到了 $G$ 等变情形下的缝合条件之后，就可以尝试对这种结构的例子进行分类。带有 $G$ 作用的 Frobenius 代数 $\mathcal{C}_1$，对应于一个具有全局 $G$ 对称性的拓扑场论。当 $\mathcal{C}_1$ 是半单 Frobenius 代数，也就是某个有限 $G$ 集 $X$ 上的函数代数时，Turaev 给出了一个很好的答案：对这一对称性进行规范化的方式，也就是将 $\mathcal{C}_1$ 扩展为 Turaev 代数的方式，对应于 $X$ 上的等变 B 场，即取值于 $\mathbb{C}^{\times}$ 的等变二上闭链。此外，要求同构在给定的 $\mathcal{C}_1$ 上为恒等映射时，两个这样的 B 场定义同构的 Turaev 代数，当且仅当它们代表 $H_G^2(X,\mathbb{C}^{\times})\cong H_G^3(X,\mathbb{Z})$ 中相同的类；这里使用了 $G$ 为有限群的假设。现在回顾这一结果，并借此介绍半单情形下 Turaev 代数 $\mathcal{C}$ 的一个更具几何性的图景。首先回顾一些十分一般的构造。

\subsubsection{一般构造}\label{sec:2.6.3.1}

\SourceParagraph{ch02-2.6.3.1-001}{85}
当群 $G$ 作用于集合 $X$ 时，可以构造范畴 $X//G$，其对象是 $X$ 的点 $x$，从 $x_0$ 到 $x_1$ 的态射为
\begin{equation}\label{eq:2.87}
\operatorname{Hom}(x_0,x_1)\coloneqq\{g\in G:gx_0=x_1\}.
\end{equation}

\begin{figure}[H]
\centering
\BookImage{01_Numbered_Figures/Chapter_02/Ch02_Fig21_PDFp085_Bookp075.pdf}
\caption{空间 $|X//G|$ 中的一个定向二单形 $\Delta_{x,g_1,g_2}$。}\label{fig:2.21}
\end{figure}

\begin{figure}[H]
\centering
\BookImage{01_Numbered_Figures/Chapter_02/Ch02_Fig22_PDFp086_Bookp076.pdf}
\caption{$|X//G|$ 中的一个定向三单形。}\label{fig:2.22}
\end{figure}

\SourceParagraph{ch02-2.6.3.1-002}{86}
其次，对于任意小范畴 $\mathcal{C}$，可以构造该范畴的空间，记为 $|\mathcal{C}|$。它是范畴的神经这个有向单纯集的几何实现，其中 $p$ 单形与范畴中可复合的 $p$ 元态射组一一对应。具体地说，顶点就是范畴的对象，边就是态射。满足 $f_3=f_2\circ f_1$ 的三元态射组 $(f_1,f_2,f_3)$ 对应于二单形，依此类推。在当前情形中，构造 $|X//G|$ 时，二单形就是\BookRef{fig:2.21}{图21}所示的三元组 $(g_1,g_2,x)$。三单形如\BookRef{fig:2.22}{图22}所示，等等。

\SourceParagraph{ch02-2.6.3.1-003}{86}
空间 $|X//G|$ 是 $(X\times EG)/G$ 的一个模型。因此，这个空间的胞腔上同调 $H^*(|X//G|,\mathbb{C}^{\times})$，就是等变上同调 $H_G^*(X,\mathbb{C}^{\times})$。

\SourceParagraph{ch02-2.6.3.1-004}{86}
可以给任意小范畴 $\mathcal{C}$ 关联的另一个对象，是复数域上的范畴代数 $A(\mathcal{C})$。它具有由 $\mathcal{C}$ 的态射标记的向量空间基 $\{\varepsilon_f\}$；当 $f_1,f_2$ 可复合时，乘积为 $\varepsilon_{f_1}\varepsilon_{f_2}=\varepsilon_{f_1\circ f_2}$，否则为零。对于范畴 $X//G$，当 $X$ 有限时，代数 $A(X//G)$ 就是算子代数理论意义下通常的交叉积代数 $A(X)\rtimes G$，其中 $A(X)$ 是集合 $X$ 上的复值函数代数\footnote{对于任意带有 $G$ 作用的交换代数 $A$，$A\rtimes G$ 由元素 $a\otimes g$ 张成，其中 $a\in A$，$g\in G$，其乘积为
\begin{equation}\label{eq:2.88}
(a_1\otimes g_1)(a_2\otimes g_2)=a_1g_1(a_2)\otimes g_1g_2.
\end{equation}
同构 $A(X//G)\to A(X)\rtimes G$，将 $\varepsilon_{g,x}$ 映为 $\chi_{gx}\otimes g$，其中 $\chi_x$ 是支集位于 $x$ 的特征函数。}。这里的拓扑场论应用中，$X$ 正是有限集合。

\SourceParagraph{ch02-2.6.3.1-005}{86}
范畴代数 $A(\mathcal{C})$ 的构造可以推广。范畴 $\mathcal{C}$ 上的一个 B 场，是这样的法则：为 $\mathcal{C}$ 的每个态射 $f$ 指定一个复直线 $L_f$，并为每一对可复合态射 $(f_1,f_2)$ 指定同构
\[
L_{f_1}\otimes L_{f_2}\longrightarrow L_{f_1\circ f_2},
\]
使这些同构满足结合律。具体地说，给定这样的乘积，等价于在空间 $|\mathcal{C}|$ 上给定一个二上闭链。事实上，选取基元素 $\ell_f\in L_f$，就必须有
\begin{equation}\label{eq:2.89}
\ell_{f_1}\cdot\ell_{f_2}=b(f_1,f_2,f_3)\ell_{f_3},
\end{equation}
其中 $f_3=f_1\circ f_2$，而 $b(f_1,f_2,f_3)\in\mathbb{C}^{\times}$ 定义了 $|\mathcal{C}|$ 上的一个二上链。（取值于 $\mathbb{C}^{\times}$ 而不是 $\mathbb{C}$，是为了保证乘积非退化。）\BookEqRef{eq:2.89}{2.89}满足结合律，当且仅当 $b$ 是二上闭链。改变 $L_f$ 的基，就会使 $b$ 改变一个上边界。因此，$\mathcal{C}$ 上 B 场的同构类，与上同调类 $[b]\in H^2(|\mathcal{C}|,\mathbb{C}^{\times})$ 一一对应。给定 $\mathcal{C}$ 上的 B 场 $b$，可以构造扭曲范畴代数 $A_b(\mathcal{C})$；作为向量空间，它是 $\bigoplus_fL_f$，乘法则由满足结合律的映射 $L_{f_1}\otimes L_{f_2}\to L_{f_1\circ f_2}$ 定义。

\SourceParagraph{ch02-2.6.3.1-006}{87}
将上述构造应用于范畴 $X//G$，直线 $L_{g,x}$ 上满足结合律的乘积，等价于 $H_G^2(X,\mathbb{C}^{\times})$ 的一个二上闭链代表。用直线 $L_{g,x}$ 的基元素 $\ell_{g,x}$，将乘法写为
\begin{equation}\label{eq:2.90}
\ell_{g_2,x_2}\ell_{g_1,x_1}=
\begin{cases}
b_{x_1}(g_2,g_1)\ell_{g_2g_1,x_1},&x_2=g_1x_1,\\
0,&\text{否则}.
\end{cases}
\end{equation}
这里，$b_{x_1}(g_2,g_1)=b(\Delta_{x_1,g_1,g_2})$ 是二上闭链在\BookRef{fig:2.21}{图21}所示定向二单形上的值。

\SourceParagraph{ch02-2.6.3.1-007}{87}
注意，如果 $G_x$ 是某点 $x\in X$ 的稳定子群，将\BookEqRef{eq:2.90}{2.90}限制到 $g\in G_x$ 的元素 $\ell_{g,x}$，就可以看出，$b_x$ 定义了群上同调 $H^2(G_x,\mathbb{C}^{\times})$ 的一个元素，对应于 $G_x$ 被 $\mathbb{C}^{\times}$ 的中心扩张；这个扩张的元素是二元组 $(g,\lambda)$，其中 $g\in G_x$，$\lambda\in L_{g,x}\setminus\{0\}$。稳定子群 $G_x$ 的这一中心扩张，一般不能延拓为整个群 $G$ 的中心扩张。不过，在一个特殊情形下可以做到：B 场 $b$ 通过映射 $X\to\{\text{点}\}$，从 $G$ 的某个二上闭链拉回；也就是说，$b_x(g_2,g_1)$ 与 $x$ 无关。一般地，经由下文所用的起点与终点坐标转换，二上闭链 $b\colon G\times G\times X\to\mathbb{C}^{\times}$ 可以视为群 $G$ 取值于带自然 $G$ 作用的 Abel 群 $A(X)^{\times}=\operatorname{Map}(X,\mathbb{C}^{\times})$ 的二上闭链。因此，它定义一个一般非中心的扩张
\[
1\longrightarrow A(X)^{\times}\longrightarrow\widetilde{G}\longrightarrow G\longrightarrow1.
\]

\SourceParagraph{ch02-2.6.3.1-008}{87}
需要注意一个技术细节：对于任意 B 场，当 $f$ 是恒等态射时，都有典范的等同 $L_f=\mathbb{C}$。因此，当 $g=1$ 时，$L_{g,x}=\mathbb{C}$。我们总是在 $g=1$ 时选取 $\ell_{g,x}=1$，由此将二上闭链归一化，使 $g_1$ 或 $g_2$ 中任意一个为1时，都有 $b_x(g_1,g_2)=1$。

\SourceParagraph{ch02-2.6.3.1-009}{87}
乘法由\BookEqRef{eq:2.90}{2.90}定义的代数 $A_b(X//G)=\bigoplus_{g\in G,\,x\in X}L_{g,x}$，可以通过
\[
\ell_{g,x}\longmapsto\chi_{gx}\otimes g
\]
与扭曲交叉积代数 $A(X)\rtimes_bG$ 等同，其中 $\chi_x$ 是支集位于 $x$ 的特征函数。扭曲交叉积定义为
\begin{equation}\label{eq:2.91}
(f_1\otimes g_1)(f_2\otimes g_2)
=\alpha_{g_1g_2}\bigl(b(g_1,g_2)\bigr)f_1\alpha_{g_1}(f_2)\otimes g_1g_2,
\end{equation}
其中，$b(g_1,g_2)$ 表示 $A(X)^{\times}$ 中的函数 $x\mapsto b_x(g_1,g_2)$，群 $G$ 以自然方式作用于 $A(X)$：
\[
\alpha_g(f)(x)=f(g^{-1}x),
\]
从而 $g\cdot\chi_x=\chi_{gx}$。公式中的因子 $\alpha_{g_1g_2}(b(g_1,g_2))$，就是将起点坐标表示的 $b$ 转换为通常左 $G$ 模群上闭链的坐标表达。

\SourceParagraph{ch02-2.6.3.1-010}{88}
如果要把这些考虑应用于\BookSectionRef{sec:2.1.6}{2.1.6}与\BookSectionRef{sec:2.2.4}{2.2.4}所描述的自旋情形，那么必须令直线 $L_f$ 具有 $\mathbb{Z}/2$ 分次。在这种情形下，还可以通过 $H^1(|\mathcal{C}|,\mathbb{Z}/2)$ 对理论作进一步扭曲。不过，这里不再讨论这一推广。

\subsubsection{与 \texorpdfstring{$G$}{G} 空间关联的 Turaev 代数}\label{sec:2.6.3.2}

\SourceParagraph{ch02-2.6.3.2-001}{88}
代数 $A_b(X//G)$ 一般不满足缝合条件，因而不是 Turaev 代数。特别地，交叉积代数通常不满足\BookEqRef{eq:2.82}{2.82}。不过，由起点对象与终点对象相同的态射定义的子范畴，对于 $X//G$ 上的任意 B 场 $b$，确实能够给出一个 Turaev 代数。我们称它为“小环路代数”。因此，定义 $\mathcal{C}=\bigoplus_g\mathcal{C}_g\subset A_b(X//G)$，其中
\begin{equation}\label{eq:2.92}
\mathcal{C}_g\coloneqq\bigoplus_{x:\,gx=x}L_{g,x},
\end{equation}
并将迹定义为
\begin{equation}\label{eq:2.93}
\theta(\ell_{g,x})=\delta_{g,1}\theta(\varepsilon_x),
\end{equation}
其中，右边的 $\theta$ 是 $\mathcal{C}_1$ 上给定的 $G$ 不变迹，$\varepsilon_x$ 是半单 Frobenius 代数 $\mathcal{C}_1=A(X)$ 中通常的幂等元；也就是说，$\varepsilon_x=1\in L_{1,x}=\mathbb{C}$。小环路代数可以用\BookRef{fig:2.23}{图23}直观表示。

\begin{figure}[H]
\centering
\BookImage{01_Numbered_Figures/Chapter_02/Ch02_Fig23_PDFp088_Bookp078.pdf}
\caption{$X=S_3/S_2$ 的小环路代数，其中 $S_n$ 是 $n$ 个字母的置换群。}\label{fig:2.23}
\end{figure}

\SourceParagraph{ch02-2.6.3.2-002}{88}
描述 $\mathcal{C}$ 的一种等价方式，是将它视为 $\mathcal{C}_1=A(X)$ 在 $A_b(X//G)=A(X)\rtimes_bG$ 中的中心化子代数。由于 $A(X)$ 位于 $\mathcal{C}$ 的中心，自然可以把 $\mathcal{C}$ 看作 $X$ 上一个代数丛的截面；这个丛在 $x\in X$ 处的纤维，是扭曲群代数 $\mathbb{C}_{b_x}[G_x]$，其中 $G_x$ 是 $x$ 的稳定子群。此外，这个代数丛具有自然的 $G$ 作用，覆盖 $G$ 在 $X$ 上的作用。要看到这一点，注意，由 B 场定义的群扩张 $\widetilde{G}=\{(f,g):f\in A(X)^{\times},\ g\in G\}$，是群 $G$ 被乘法群 $A(X)^{\times}$ 的扩张；它嵌入 $A(X)\rtimes_bG$ 的可逆元群，其共轭作用保持子代数 $A(X)$。由于 $A(X)$ 位于 $\mathcal{C}$ 的中心，这就意味着，共轭作用可以下降为 $G$ 在代数 $\mathcal{C}$ 上的作用。不过要注意，在更大的代数 $A_b(X//G)$ 上，作用的是 $\widetilde{G}$，而不能一般地下降为 $G$。

\SourceParagraph{ch02-2.6.3.2-003}{89}
用显式公式表示，$G$ 在代数 $\mathcal{C}$ 上的作用为
\begin{equation}\label{eq:2.94}
\alpha_{g_1}(\ell_{g_2,x})
=\ell_{g_1,x}\ell_{g_2,x}\ell_{g_1,x}^{-1}
=z_x(g_2,g_1)\ell_{g_1g_2g_1^{-1},\,g_1x},
\end{equation}
其中，逆元按照相应态射直线的逆态射来理解，并且
\begin{equation}\label{eq:2.95}
z_x(g_2,g_1)=
\frac{b_x(g_1,g_2)b_{g_1x}(g_1g_2,g_1^{-1})}{b_{g_1x}(g_1,g_1^{-1})}.
\end{equation}

\SourceParagraph{ch02-2.6.3.2-004}{89}
这样得到一个 Turaev 代数，记为 $\mathcal{C}=T(X,b,\theta)$。唯一需要认真验证的，是“环面”公理\BookEqRef{eq:2.83}{2.83}。事实上，容易看出，其等价形式\BookEqRef{eq:2.84}{2.84}的两端都等于
\[
\sum_{x:\,hx=gx=x}\frac{1}{\theta(\varepsilon_x)}
\ell_{h,x}\ell_{g,x}\ell_{h,x}^{-1}\ell_{g,x}^{-1},
\]
其中，$x$ 遍历集合 $\{x\in X:hx=gx=x\}$。

\SourceParagraph{ch02-2.6.3.2-005}{89}
Turaev 证明，上述构造是在半单情形下可能出现的最一般的构造。

\SourceParagraph{ch02-2.6.3.2-006}{89}
\BookStatementNumber{13}
\begin{theorem}\label{thm:2.13}
（文献\BookCite{ref452}，定理3.6）设 $\mathcal{C}$ 是 Turaev 代数。如果 $\mathcal{C}_1$ 半单，那么，对于某个二上闭链 $b\in Z_G^2(X,\mathbb{C}^{\times})$，$\mathcal{C}$ 就是 $X=\operatorname{Spec}(\mathcal{C}_1)$ 上的小环路扭曲代数 $T(X,b,\theta)$。
\end{theorem}

\SourceParagraph{ch02-2.6.3.2-007}{89}
\begin{proof}
如果 $\mathcal{C}_1$ 半单，就可以用基本幂等元 $\varepsilon_x$ 将它分解。$\mathcal{C}_g$ 是 $\mathcal{C}_1$ 模，因此，可以与有限集合 $X$ 上一个向量丛的截面空间等同；该丛在 $x$ 处的纤维为 $\mathcal{C}_{g,x}=\varepsilon_x\mathcal{C}_g$。（这是所谓 Serre--Swan 定理的一个平凡情形。）现在考虑 $h=1$ 时的环面公理\BookEqRef{eq:2.83}{2.83}。有 $\Delta_1=\sum_x\theta(\varepsilon_x)^{-1}\varepsilon_x\otimes\varepsilon_x$，因此，相关的一端为
\[
\sum_x\theta(\varepsilon_x)^{-1}\alpha_g(\varepsilon_x)\varepsilon_x
=\sum_{x:\,gx=x}\theta(\varepsilon_x)^{-1}\varepsilon_x,
\]
其中，第二个求和仅对满足 $gx=x$ 的 $x$ 进行。另一方面，如果 $\{a_{x,i}\}$ 是 $\mathcal{C}_{g,x}$ 的一组基，$\{a_{x,i}^*\}$ 是 $\mathcal{C}_{g^{-1},x}$ 中的对偶基，那么容易计算出 $a_{x,i}a_{x,i}^*=\theta(\varepsilon_x)^{-1}\varepsilon_x$，所以环面公理的另一端为
\[
\sum_x\theta(\varepsilon_x)^{-1}\dim(\mathcal{C}_{g,x})\varepsilon_x.
\]

\SourceParagraph{ch02-2.6.3.2-008}{90}
因此，这一公理说明，当 $gx=x$ 时，$\mathcal{C}_{g,x}$ 是一个一维空间 $L_{g,x}$，否则为零。$\mathcal{C}$ 中的乘法，使这些直线构成 $X//G$ 的小环路范畴上的一个相容的 $G$ 等变 B 场。最后，不难证明，$X//G$ 上的 B 场范畴，与小环路范畴上的这种相容的 $G$ 等变 B 场范畴等价；这里略去细节。\BookQED
\end{proof}

\SourceParagraph{ch02-2.6.3.2-009}{90}
现在考虑对 Turaev 代数 $\mathcal{C}=T(X,b,\theta)$ 所定义的等变理论进行规范化后，得到的轨形理论。在\BookSectionRef{sec:2.6.2}{2.6.2}中已经看到，它由交换 Frobenius 代数 $\mathcal{C}_{\mathrm{orb}}$ 定义，后者就是 $\mathcal{C}$ 的 $G$ 不变子代数。对于 $G$ 空间 $X$ 的 Turaev 代数，有下面的定理。

\SourceParagraph{ch02-2.6.3.2-010}{90}
\BookStatementNumber{14}
\begin{theorem}\label{thm:2.14}
轨形代数 $\mathcal{C}_{\mathrm{orb}}$ 是交叉积代数 $A(X)\rtimes_bG$ 的中心。它是有限集合 $(X/G)_{\mathrm{string}}$ 上的函数代数；这个集合是轨道空间 $X/G$ 的一个“增厚”：每个二元组 $(\xi,\rho)$ 对应一个点，其中 $\xi$ 是一个轨道，$\rho$ 是某点 $x\in\xi$ 的稳定子群 $G_x$ 的不可约射影表示，其射影二上闭链 $b_x$ 由 B 场定义。
\end{theorem}

\SourceParagraph{ch02-2.6.3.2-011}{90}
\begin{proof}
Turaev 代数 $\mathcal{C}$，由 $A(X)\rtimes_bG$ 中与 $A(X)$ 交换的元素组成。一个 $A(X)\rtimes_bG$ 的元素位于其中心，当且仅当它既与 $A(X)$ 交换，又与群元素 $G$ 的各个提升交换；也就是说，它是 $G$ 不变的。

\SourceParagraph{ch02-2.6.3.2-012}{90}
前面已经看到，$\mathcal{C}$ 是按 $X$ 的点 $x$ 取扭曲群代数 $\mathbb{C}_{b_x}[G_x]$ 的乘积。因此，其不变部分，是按轨道 $\xi$ 取 $\mathbb{C}_{b_x}[G_x]$ 的 $G_x$ 不变部分的乘积；这些不变部分，就是 $\mathbb{C}_{b_x}[G_x]$ 的中心，每个具有二上闭链 $b_x$ 的不可约表示 $\rho$，对应于中心中的一个 $\mathbb{C}$ 副本。\BookQED
\end{proof}

\SourceParagraph{ch02-2.6.3.2-013}{90}
Turaev 代数 $\mathcal{C}=T(X,b,\theta)$，位于 $\mathcal{C}_{\mathrm{orb}}$ 与 $A(X)\rtimes_bG$ 之间。将在\BookRef{thm:2.16}{定理2.16}中看到，$A(X)\rtimes_bG$ 是半单的，因此与其中心 $\mathcal{C}_{\mathrm{orb}}$ Morita 等价\footnote{这意味着，$A(X)\rtimes_bG$ 的表示范畴，与 $\mathcal{C}_{\mathrm{orb}}$ 的表示范畴等价。在固定中心作用之后，这种等价仅差张量上一个“线丛”；这里的线丛，是一个 $\mathcal{C}_{\mathrm{orb}}$ 的表示 $L$，存在另一个表示 $L'$，使 $L\otimes_{\mathcal{C}_{\mathrm{orb}}}L'\cong\mathcal{C}_{\mathrm{orb}}$。}。不过，Turaev 代数比轨形理论保留了更多信息：它编码了 $X$ 及其 $G$ 作用。当阶为 $n$ 的群 $G$ 在 $X$ 上自由作用时，这一区别最明显；此时，$A(X)\rtimes G$ 是按 $X$ 的每个 $G$ 轨道取一份 $n\times n$ 矩阵代数的乘积，因而无法从中区分 $X$ 的各个点。将在\BookSectionRef{sec:2.6.5}{2.6.5}中看到，至少在半单情形下，规范化理论 $\mathcal{C}$ 的边界条件范畴，是轨形理论的边界条件范畴的一个自然充实。

\SourceParagraph{ch02-2.6.3.2-014}{91}
鉴于交叉积代数在文献\BookCite{ref141,ref308,ref352}发展的轨形上 D 膜理论中处于核心地位，时空的交叉积代数 $A(X)\rtimes G$ 居然不是 $G$ 等变拓扑场论所需的 Frobenius 代数，也许会令人意外。事实上，这与本章的思想十分契合。Turaev 代数记住了 $X$ 的各个点，因此仅允许上面所说的“小环路”。由此，原本旨在形式化世界面局域性的缝合条件，也编码了一种初步的时空局域性。

\SourceParagraph{ch02-2.6.3.2-015}{91}
本小节最后，联系通常用路径积分表示的环面轨形配分函数。为此，计算 $\mathcal{C}$ 中投影到 $G$ 不变态的投影算符，从而求出 $\dim\mathcal{C}_{\mathrm{orb}}$。注意，只有在 $[g,h]=1$ 且 $gx=x$ 时，$\alpha_g(\ell_{h,x})$ 才是 $\ell_{h,x}$ 的倍数；此时，有
\begin{equation}\label{eq:2.96}
\alpha_g(\ell_{h,x})=\frac{b_x(g,h)}{b_x(h,g)}\ell_{h,x},
\end{equation}
这里将\BookEqRef{eq:2.95}{2.95}与二上闭链恒等式结合使用。因此，得到
\begin{equation}\label{eq:2.97}
\dim\mathcal{C}_{\mathrm{orb}}=
\frac{1}{|G|}\sum_{\substack{g,h\in G\\gh=hg}}\ 
\sum_{x:\,x=gx=hx}\frac{b_x(g,h)}{b_x(h,g)}.
\end{equation}

\subsection{等变开闭弦理论的缝合条件}\label{sec:2.6.4}

\SourceParagraph{ch02-2.6.4-001}{91}
现在转而考虑 $G$ 等变开闭弦理论。扩大范畴 $\mathcal{S}_G$，使其对象成为带有边界、各端点有标记、并配备主 $G$ 丛的定向一维流形。态射仍是非等变情形中的那些配边，但要配备 $G$ 丛。在同构意义下，区间上只有一个 $G$ 丛：它是平凡丛，且具有自同构群 $G$。因此，等变理论为每对标记 $a,b$ 给出一个带有 $G$ 作用的向量空间 $\mathcal{O}_{ab}$。$g\in G$ 在 $\mathcal{O}_{ab}$ 上的作用，可以看作来自“正方形”配边，其上所配的丛沿每条“受约束”边的平行移动均为 $g$。还有一个 $G$ 等变复合律 $\mathcal{O}_{ab}\times\mathcal{O}_{bc}\to\mathcal{O}_{ac}$。\BookRef{fig:2.24}{图24}与\BookRef{fig:2.25}{图25}说明了这些构造。

\begin{figure}[H]
\centering
\BookImage{01_Numbered_Figures/Chapter_02/Ch02_Fig24_PDFp091_Bookp081.pdf}
\caption{波浪线表示受约束边界。如果沿虚线路径 $\mathcal{P}$ 的平行移动为 $g$，那么这个态射就给出 $\mathcal{O}$ 上的 $G$ 作用。}\label{fig:2.24}
\end{figure}

\begin{figure}[H]
\centering
\BookImage{01_Numbered_Figures/Chapter_02/Ch02_Fig25_PDFp092_Bookp082.pdf}
\caption{$\mathcal{O}$ 中乘法的定义。沿所有虚线路径的平行移动均为1。注意乘法的顺序。}\label{fig:2.25}
\end{figure}

\SourceParagraph{ch02-2.6.4-002}{92}
在开闭弦情形中，与方程\BookEqRef{eq:2.2}{2.2}至\BookEqRef{eq:2.12}{2.12}相应的条件如下。

\SourceParagraph{ch02-2.6.4-003}{92}
先只考虑单个标记 $a$。有一个未必交换的 Frobenius 代数 $(\mathcal{O}_{aa}=\mathcal{O},\theta_{\mathcal{O}})$，以及一个 $G$ 作用 $\rho\colon G\to\operatorname{Aut}(\mathcal{O})$：
\begin{equation}\label{eq:2.98}
\rho_g(\psi_1\psi_2)=(\rho_g\psi_1)(\rho_g\psi_2),
\end{equation}
并且这个作用保持迹，即 $\theta_{\mathcal{O}}(\rho_g\psi)=\theta_{\mathcal{O}}(\psi)$。见\BookRef{fig:2.26}{图26}。

\begin{figure}[H]
\centering
\BookImage{01_Numbered_Figures/Chapter_02/Ch02_Fig26_PDFp092_Bookp082.pdf}
\caption{说明 $G$ 作为自同构群作用于 $\mathcal{O}$。}\label{fig:2.26}
\end{figure}

\SourceParagraph{ch02-2.6.4-004}{92}
还有带 $G$ 扭曲的开闭弦转变映射
\begin{equation}\label{eq:2.99}
\begin{aligned}
\iota_{g,a}=\iota_g&\colon\mathcal{C}_g\longrightarrow\mathcal{O}_{aa}=\mathcal{O},\\
\iota^{g,a}=\iota^g&\colon\mathcal{O}_{aa}=\mathcal{O}\longrightarrow\mathcal{C}_g,
\end{aligned}
\end{equation}
它们是等变的：
\begin{figure}[H]
\begin{equation}\label{eq:2.100}
\BookImage{02_Unnumbered_Diagrams/Chapter_02/Ch02_Diagram_PDFp093_01_Bookp083.pdf}
\end{equation}
\end{figure}
\begin{figure}[H]
\begin{equation}\label{eq:2.101}
\BookImage{02_Unnumbered_Diagrams/Chapter_02/Ch02_Diagram_PDFp093_02_Bookp083.pdf}
\end{equation}
\end{figure}

\begin{figure}[H]
\centering
\BookImage{01_Numbered_Figures/Chapter_02/Ch02_Fig27_PDFp093_Bookp083.pdf}
\caption{开闭弦转变 $\iota_g$ 与 $\iota^g$。}\label{fig:2.27}
\end{figure}

\SourceParagraph{ch02-2.6.4-005}{93}
这些映射如\BookRef{fig:2.27}{图27}所示。开闭弦映射必须满足\BookSectionRef{sec:2.1.2}{2.1.2}中\BookRef{cond:2.1.2.3a}{条件3a}至\BookRef{cond:2.1.2.3e}{3e}的 $G$ 扭曲版本。特别地，将各个 $\iota_g$ 合在一起所得的映射 $\iota\colon\mathcal{C}\to\mathcal{O}$，是一个环反同态，也就是从 $\mathcal{C}^{\mathrm{op}}$ 到 $\mathcal{O}$ 的环同态：
\begin{equation}\label{eq:2.102}
\iota_{g_1}(\phi_1)\iota_{g_2}(\phi_2)
=\iota_{g_2g_1}(\phi_2\phi_1),
\qquad\phi_1\in\mathcal{C}_{g_1},\ \phi_2\in\mathcal{C}_{g_2}.
\end{equation}
由于单位元位于 $\mathcal{C}_1$ 中，条件\BookEqRef{eq:2.8}{2.8}保持原样。$G$ 扭曲的中心性条件为
\begin{equation}\label{eq:2.103}
\iota_g(\phi)(\rho_g\psi)=\psi\iota_g(\phi),
\qquad\phi\in\mathcal{C}_g,\ \psi\in\mathcal{O},
\end{equation}
如\BookRef{fig:2.28}{图28}所示。

\begin{figure}[H]
\centering
\BookImage{01_Numbered_Figures/Chapter_02/Ch02_Fig28_PDFp094_Bookp084.pdf}
\caption{$G$ 扭曲的中心性公理。}\label{fig:2.28}
\end{figure}

\SourceParagraph{ch02-2.6.4-006}{93}
$G$ 扭曲的伴随条件为
\begin{equation}\label{eq:2.104}
\theta_{\mathcal{O}}\bigl(\psi\iota_{g^{-1}}(\phi)\bigr)
=\theta_{\mathcal{C}}\bigl(\iota^g(\psi)\phi\bigr),
\qquad\phi\in\mathcal{C}_{g^{-1}},
\end{equation}
如\BookRef{fig:2.29}{图29}所示。

\begin{figure}[H]
\centering
\BookImage{01_Numbered_Figures/Chapter_02/Ch02_Fig29_PDFp095_Bookp085.pdf}
\caption{$G$ 扭曲的伴随关系。上图是一个挖去两个圆盘的球面；外圈并不表示边界。}\label{fig:2.29}
\end{figure}

\SourceParagraph{ch02-2.6.4-007}{94}
$G$ 扭曲的 Cardy 条件，不仅限制各个代数 $\mathcal{O}_{aa}$，还限制任意 $a,b$ 之间的态射空间 $\mathcal{O}_{ab}$。对于每个 $g\in G$，都必须有
\begin{equation}\label{eq:2.105}
\pi_{g,b}^a=\iota_{g,b}\iota^{g,a}.
\end{equation}
这里，$\pi_{g,b}^a$ 定义为
\begin{equation}\label{eq:2.106}
\pi_{g,b}^a(\psi)=\sum_\mu\psi^\mu\psi(\rho_g\psi_\mu),
\end{equation}
其中，$\psi_\mu$ 遍历 $\mathcal{O}_{ab}$ 的一组基，$\psi^\mu$ 为 $\mathcal{O}_{ba}$ 中的对偶基。见\BookRef{fig:2.30}{图30}。



\SourceParagraph{ch02-2.6.4-008}{94}
现在可以表述下面的定理。

\SourceParagraph{ch02-2.6.4-009}{94}
\BookStatementNumber{15}
\begin{theorem}\label{thm:2.15}
上述条件构成 $G$ 等变二维开闭弦 TFT 的一组完整的缝合条件。
\end{theorem}

\SourceParagraph{ch02-2.6.4-010}{94}
这个定理将在\BookSectionRef{sec:2.7}{2.7}中证明。注意，上述公理稍有冗余，因为\BookEqRef{eq:2.100}{2.100}与\BookEqRef{eq:2.104}{2.104}已经蕴含\BookEqRef{eq:2.101}{2.101}。

\begin{figure}[H]
\centering
\BookImage{01_Numbered_Figures/Chapter_02/Ch02_Fig30_PDFp096_Bookp086.pdf}
\caption{$G$ 扭曲的 Cardy 条件。在双扭转图中，绕 $\mathcal{P}_1$ 的平行移动为1，绕 $\mathcal{P}_2$ 的平行移动为 $g$。}\label{fig:2.30}
\end{figure}

\subsection{半单 \texorpdfstring{$\mathcal{C}$}{C} 的缝合条件的解}\label{sec:2.6.5}

\SourceParagraph{ch02-2.6.5-001}{95}
现在证明，当 $\mathcal{C}_1$ 半单时，上述缝合条件的解，由 $X=\operatorname{Spec}(\mathcal{C}_1)$ 上的扭曲 $G$ 等变丛给出；按照上一小节的反同态约定，其协变表示的扭曲由 $\mathcal{C}$ 所定义的 B 场的逆场给出。

\SourceParagraph{ch02-2.6.5-002}{95}
先对这些丛作一些说明。给定交叉积代数 $A(X)\rtimes G$ 的一个有限维表示，等价于给定 $A(X)$ 的一个表示，也就是 $X$ 上的向量丛 $E$，以及一个与之相容的 $G$ 作用。因此，$A(X)\rtimes G$ 的表示，恰好就是 $X$ 上的 $G$ 等变向量丛。对于有限群 $G$，可以用许多等价方式定义扭曲 $G$ 等变向量丛，其扭曲由代表 $H_G^2(X,\mathbb{C}^{\times})$ 中一个元素的 B 场 $b_o$ 给出。对当前目的来说，最简单的定义，是将扭曲丛定义为 $A(X)\rtimes_{b_o}G$ 的表示；这里，相应闭弦二上闭链为 $b$ 时，取 $b_o=b^{-1}$。（遗憾的是，当 $G$ 不是有限群时，不能直接沿用这一有限维模的描述，所以文献\BookCite{ref026}使用了另一种描述。本节讨论结束时，将解释二者之间的关系。）

\SourceParagraph{ch02-2.6.5-003}{95}
问题很容易约化为单个 $G$ 轨道上的问题，因此可以假设，对于 $G$ 的某个子群 $H$，有 $X=G/H$。相应地，闭弦 Frobenius 数据，由二上闭链 $b$ 与定义迹的单个常数 $\theta_c\in\mathbb{C}^{\times}$ 指定：$\theta(\ell_{g,x})=\delta_{g,1}\theta_c$。与往常一样，在固定 $\theta_c$ 及给定轨道的条件下，同构类只依赖于 $[b]\in H^2(H,\mathbb{C}^{\times})$。

\SourceParagraph{ch02-2.6.5-004}{96}
\BookStatementNumber{16}
\begin{theorem}\label{thm:2.16}
设 $\mathcal{C}=T(X,b,\theta_c)$ 是一个 Turaev 代数，其中 $\mathcal{C}_1$ 半单，且 $X=G/H$。对于单个标记 $a$，缝合约束的一般解 $\mathcal{O}=\mathcal{O}_{aa}$，由平方根 $\theta_o=\sqrt{\theta_c}$ 的一个选择，以及 $H$ 的一个射影表示 $V$ 确定；其二上闭链 $b_o$ 是 $b$ 在 $H$ 上的限制的逆。

代数 $\mathcal{O}$，是 $X$ 上的 $G$ 等变代数丛的截面代数：
\begin{equation}\label{eq:2.107}
\mathcal{O}\coloneqq\Gamma\bigl(G\times_H\operatorname{End}(V)\bigr)
=\operatorname{Ind}_H^G\bigl(\operatorname{End}(V)\bigr),
\end{equation}
迹由 $\theta_o$ 确定：
\begin{equation}\label{eq:2.108}
\theta_{\mathcal{O}}(\Psi)=\theta_o\sum_{x\in G/H}\operatorname{Tr}_{V}\bigl(\Psi(x)\bigr).
\end{equation}
\end{theorem}

\SourceParagraph{ch02-2.6.5-005}{96}
\begin{proof}
假设给定一个 $\mathcal{C}_1$ 半单的 Turaev 代数 $\mathcal{C}$，以及满足缝合条件的 $\mathcal{O},\theta_{\mathcal{O}},\iota_g,\iota^g$。设 $X$ 是 $G$ 空间 $\operatorname{Spec}(\mathcal{C}_1)$。由非等变情形中的结果，知道对于 $X$ 上某个向量丛 $E$，有 $\mathcal{O}=\operatorname{End}_{\mathcal{C}_1}(\Gamma(E))=\Gamma(\operatorname{End}(E))$；$E$ 的选择仅差张量上一个线丛 $L\to X$。因此，$\mathcal{O}=\bigoplus_x\mathcal{O}_x$，其中 $\mathcal{O}_x=\operatorname{End}(E_x)$。还知道，$\mathcal{O}$ 上的迹必须由\BookEqRef{eq:2.35}{2.35}给出。为了使 $\theta\colon\mathcal{O}\to\mathbb{C}$ 是 $G$ 不变的，对同一轨道上的每个 $x$，必须选取 $\theta_x$ 的同一个平方根 $\theta_o$。

\SourceParagraph{ch02-2.6.5-006}{97}
$G$ 以相容的代数同构作用于 $\mathcal{C}_1$ 与 $\mathcal{O}$，因此，如果 $x=gx_0$，$g\in G$ 就通过代数同构将 $\mathcal{O}_{x_0}$ 映到 $\mathcal{O}_x$。这证明了\BookEqRef{eq:2.107}{2.107}，其中 $V=E_{x_0}$。最后，Turaev 代数 $\mathcal{C}$ 是各 $\mathcal{C}_x$ 的直和，其中 $\mathcal{C}_x$ 是稳定子群 $G_x$ 的扭曲群代数，扭曲由 $b_x$ 给出。反同态 $\mathcal{C}\to\mathcal{O}$，在 $E_x$ 上给出 $\mathcal{C}_x^{\mathrm{op}}$ 的作用。对于 $g,h\in G_x$，令 $R_g=\iota_g(\ell_{g,x})^{-1}$，则
\[
R_gR_h=b_x(g,h)^{-1}R_{gh},\qquad\rho_g=\operatorname{Ad}(R_g).
\]
因此，$V$ 是 $H=G_{x_0}$ 的一个射影表示，其二上闭链为 $b_{x_0}^{-1}$。

\SourceParagraph{ch02-2.6.5-007}{97}
这就证明，$\mathcal{O}$ 具有定理所述形式。还需要验证，定义\BookEqRef{eq:2.107}{2.107}确实给出缝合条件的一个解；这一验证不会遇到困难。\BookQED
\end{proof}

\SourceParagraph{ch02-2.6.5-008}{97}
\BookStatementNumber{17}
\begin{note}[2.17]\label{note:2.17}
虽然在定理的假设中，给定的是代表 $H_G^2(X,\mathbb{C}^{\times})$ 中一个元素的二上闭链 $b$，但结论只使用它在 $x_0$ 的稳定子群上的限制，以逆上闭链 $b_{x_0}^{-1}$ 表示开弦扭曲。这并不令人意外：上同调等价的二上闭链 $b$ 定义同构的 Turaev 代数；而当 $X=G/H$ 时，$H_G^2(X,\mathbb{C}^{\times})$ 与群上同调 $H_H^2(\{\text{点}\},\mathbb{C}^{\times})$ 典范同构。
\end{note}

\SourceParagraph{ch02-2.6.5-009}{97}
现在可以完整描述边界条件范畴；使用的论证，与由\BookRef{thm:2.6}{定理2.6}得到\BookRef{thm:2.8}{定理2.8}时完全相同。

\SourceParagraph{ch02-2.6.5-010}{97}
\BookStatementNumber{18}
\begin{theorem}\label{thm:2.18}
如果 $\mathcal{C}$ 是 $\mathcal{C}_1$ 半单的 Turaev 代数，对应于带 B 场 $b$ 的时空 $X$，那么，$\mathcal{C}$ 的极大边界条件范畴，等价于 $X$ 上经 $b_o=b^{-1}$ 扭曲的 $G$ 等变向量丛范畴；这一等价仅差张量上一个 $X$ 上的 $G$ 等变线丛。其 Frobenius 结构，由伸缩子场 $\theta$ 的一个 $G$ 不变平方根选择确定。
\end{theorem}

\SourceParagraph{ch02-2.6.5-011}{97}
需要解释这个定理的含义。给定 Turaev 代数的等变边界条件的线性范畴，是所谓“充实”范畴的一个例子：对于每对对象 $a,b$，向量空间 $\mathcal{O}_{ab}$ 都有群 $G$ 的作用。有限维 $G$ 表示的范畴 $\operatorname{Vect}_G$，是一个以 $\mathbb{C}$ 为单位对象的对称张量范畴。说一个范畴在 $\operatorname{Vect}_G$ 这样的张量范畴中充实，意味着给定
\begin{enumerate}[label=(\roman*)]
\item 一个对象集合；
\item 对于每对对象 $a,b$，一个 $\operatorname{Vect}_G$ 中的对象 $\mathcal{O}_{ab}$；
\item 对于每组三个对象 $a,b,c$，一个满足结合律的 $G$ 向量空间“复合”态射
\[
\mathcal{O}_{ab}\otimes\mathcal{O}_{bc}\longrightarrow\mathcal{O}_{ac}.
\]
\end{enumerate}

\SourceParagraph{ch02-2.6.5-012}{98}
连同与复合相容的单位态射，公理几乎与普通范畴的公理完全一样，只是态射空间带有额外结构。在这样的情形下，称该范畴是一个普通线性范畴的充实；后者从 $b$ 到 $a$ 的态射是 $F(\mathcal{O}_{ab})$，其中函子 $F\colon\operatorname{Vect}_G\to\operatorname{Vect}$ 定义为 $F(V)=\operatorname{Hom}_G(\mathbb{C},V)=V^G$。不过，还可以仅仅遗忘 $G$ 作用，得到与这一充实范畴关联的另一个普通范畴；其从 $b$ 到 $a$ 的态射，就是作为向量空间的 $\mathcal{O}_{ab}$。

\SourceParagraph{ch02-2.6.5-013}{98}
在 $\operatorname{Vect}_G$ 中充实的一个例子，是 $\widetilde{G}$ 的有限维表示范畴，其中 $\widetilde{G}$ 是由固定二上闭链给出的、$G$ 被圆周群的中心扩张，并要求中心圆周以标量乘法作用。事实上，对于两个这样的表示，$V_1^*\otimes V_2$ 是 $G$ 的表示。

\SourceParagraph{ch02-2.6.5-014}{98}
\BookRef{thm:2.18}{定理2.18}实际上应当作如下更完整的表述。$X$ 上经 $b_o$ 扭曲的 $G$ 等变向量丛范畴，在 $\operatorname{Vect}_G$ 中有自然的充实结构：其态射 $G$ 向量空间，由扭曲向量丛之间的同态组成，这些同态不必关于 $G$ 作用等变。这个充实范畴，与等变边界条件范畴等价。取 $G$ 不变态射得到的底层普通范畴，则是轨形理论的边界条件范畴。

\SourceParagraph{ch02-2.6.5-015}{98}
\BookRef{thm:2.18}{定理2.18}有一个逆命题，它是\BookSectionRef{sec:2.2.3}{2.2.3}讨论的 $G$ 等变推广。

\SourceParagraph{ch02-2.6.5-016}{98}
\BookStatementNumber{19}
\begin{theorem}\label{thm:2.19}
设 $\mathcal{B}$ 是在 $\operatorname{Vect}_G$ 中充实的线性范畴，具有使它成为 Frobenius 范畴的 $G$ 等变迹；假设遗忘 $G$ 作用所得的线性范畴半单，且只有有限多个单对象同构类。那么，$\mathcal{B}$ 等价于某个典范等变拓扑场论的等变边界条件范畴。定义这一理论的 Turaev 代数为 $\bigoplus_g\mathcal{C}_g$，其中，$\mathcal{C}_g$ 的元素是由 $\mathcal{B}$ 的对象 $a$ 标记的一族 $\phi_a\in\mathcal{O}_{aa}$，对每个 $f\in\mathcal{O}_{ab}$ 满足
\begin{equation}\label{eq:2.109}
\phi_a\circ f=(g^{-1}\cdot f)\circ\phi_b.
\end{equation}
各分量的乘法，以这些族的反向分量复合给出，从而 $\mathcal{C}_g\mathcal{C}_h\subset\mathcal{C}_{gh}$。
\end{theorem}

\SourceParagraph{ch02-2.6.5-017}{98}
要证明这一点，必须说明\BookEqRef{eq:2.109}{2.109}确实定义了一个 Turaev 代数。细节直接，这里略去。

\subsection{等变边界态}\label{sec:2.6.6}

\SourceParagraph{ch02-2.6.6-001}{98}
最后，考虑\BookSectionRef{sec:2.3}{2.3}讨论的“边界态”的等变版本。我们所定义的等变边界条件范畴，其底层普通范畴取 $G$ 不变态射；前面已经指出，它恰好与对规范场求和所得轨形理论的边界条件范畴相同。为了用边界态重新表述这一点，先给出定义。

\SourceParagraph{ch02-2.6.6-002}{99}
在与 Turaev 代数 $\mathcal{C}=T(X,b,\theta)$ 关联的规范化理论中，世界面各点处带扭曲扇区标记的可观测量，正是 $\mathcal{C}$ 的元素。与边界条件 $a$ 关联的边界态 $B_a\in\mathcal{C}$，由如下性质刻画：在边界为 $S^1$、边界条件为 $a$ 的曲面 $\Sigma$ 上，任取绕边界的平行移动，可观测量 $\phi_1,\ldots,\phi_k$ 在各点处的关联函数，等于同样这些可观测量在封住边界所得闭曲面上的关联函数，只需在封口的中心额外插入 $B_a$。由于场论的因子分解性质，只需验证 $\Sigma$ 为圆盘的情形。圆盘上的关联函数，通过圆环将 $\phi_1\cdots\phi_k\in\mathcal{C}_g$ 传播到 $\mathcal{H}(\varnothing)=\mathbb{C}$ 而得到；这个圆环中，不作为入边界的圆周受边界条件 $a$ 约束，沿该圆周的平行移动必为 $g$。按照我们的规则，结果为
\[
\theta_{\mathcal{O}_{aa}}\bigl(\iota_{g,a}(\phi_1\cdots\phi_k)\bigr).
\]
令它等于 $\theta_{\mathcal{C}_1}(\phi_1\cdots\phi_kB_a)$，就得到
\[
B_a=\sum_g\iota^{g,a}(1).
\]

\SourceParagraph{ch02-2.6.6-003}{99}
映射 $a\mapsto B_a$ 的像，显然位于 Turaev 代数的 $G$ 不变部分，也就是\BookRef{thm:2.14}{定理2.14}所确定的交叉积代数的中心。令 $h\in H_G^3(X,\mathbb{Z})$ 为开弦扭曲 $b_o=b^{-1}$ 所对应的类，这个映射延拓为同态
\[
K_{G,h}(X)\longrightarrow T(X,b,\theta)^G.
\]
于是有下面的定理。

\SourceParagraph{ch02-2.6.6-004}{99}
\BookStatementNumber{20}
\begin{theorem}\label{thm:2.20}
$G$ 不变边界态在 $T(X,b,\theta)^G$ 中生成一个整数格，它与扭曲等变 K 理论之间的关系，由复向量空间同构
\begin{equation}\label{eq:2.110}
K_{G,h}(X)\otimes_{\mathbb{Z}}\mathbb{C}\cong T(X,b,\theta)^G
\end{equation}
给出。
\end{theorem}

\SourceParagraph{ch02-2.6.6-005}{99}
\BookStatementNumber{21}
\begin{note}[2.21]\label{note:2.21}
\begin{enumerate}[label=(\arabic*)]
\item 方程\BookEqRef{eq:2.110}{2.110}与文献\BookCite{ref121}中早期的一个观察有关。如果 $X=G$，且 $G$ 以共轭方式作用于自身，那么 $T(X,0)^G$，就是只有一个表示的手征代数所对应的轨形共形场论中出现的 Verlinde 代数\BookCite{ref121}。不同的轨道，就是 $G$ 的共轭类。只考虑一个共轭类 $[g]$，就可以与上述结果比较：选取 $g$ 的中心化子的各个不可约特征标，给出一组态的基。这些正是上面得到的 $G$ 不变边界态。
\item 将上述结果译成轨形上的膜的语言，可以表述如下。与不同的 $b_o$ 不可约表示 $V_i$ 对应的边界态，就是文献\BookCite{ref141}中的“分数膜”。使用射影表示的建议，最初见于文献\BookCite{ref141}；文献\BookCite{ref132,ref136}论证了它们与离散挠率的对应关系。开弦扇区与闭弦扇区的二上闭链，在按边界取向作对应之后上同调等价。这一等价的另一证明，见文献\BookCite{ref015}。
\end{enumerate}
\end{note}

\SourceParagraph{ch02-2.6.6-006}{100}
本节最后，回过头来解释，用 $A(X)\rtimes_{b_o}G$ 模定义扭曲等变 K 理论的方法，与文献\BookCite{ref026}中的定义之间的关系。

\SourceParagraph{ch02-2.6.6-007}{100}
在文献\BookCite{ref026}中，扭曲等变理论的元素按如下方式描述。首先，用 $X$ 上的稳定射影 Hilbert 空间丛 $P$ 表示扭曲类 $h\in H_G^3(X,\mathbb{Z})$，并在 $P$ 上配备一个覆盖 $X$ 上 $G$ 作用的 $G$ 作用。然后，用与 $P$ 关联的 Fredholm 算符丛中的等变逐纤维算符族 $\{T_x\}_{x\in X}$，表示 $K_{G,P}(X)$ 的元素。下面说明，如何为有限生成 $A(X)\rtimes_{b_o}G$ 模关联这样的一对 $(P,\{T_x\})$。一个这样的模，等价于配备了相容的扩张群 $\widetilde{G}_o$ 作用的有限生成 $A(X)$ 模；其中，$\widetilde{G}_o$ 是前面描述的、与 $b_o$ 关联的扩张群。等价地，它是 $X$ 上的有限维向量丛 $E$，其总空间具有覆盖 $X$ 上 $G$ 作用的 $\widetilde{G}_o$ 作用。在同一上同调类中选取圆周值的二上闭链代表，并为 $E$ 选取相容的 Hermitian 度量。选取一个固定的可分无限维 Hilbert 空间 $\mathcal{H}$。先假设，$E_x$ 对每个稳定子群都包含该扭曲下的所有不可约射影表示。那么，$\widehat{E}=E\otimes\mathcal{H}$ 是 $X$ 上的稳定 Hilbert 丛，关联的射影空间丛 $P=\mathbb{P}(\widehat{E})$ 具有自然的 $G$ 作用，并代表 $H_G^3(X,\mathbb{Z})$ 中的类 $h$；参见文献\BookCite{ref026}中命题6.3的证明。如果 $T\colon\mathcal{H}\to\mathcal{H}$ 是一个固定的满射 Fredholm 算符，核为一维，那么 $\operatorname{id}_E\otimes T$ 是 $\widehat{E}$ 上逐纤维的 Fredholm 算符；它在 $P$ 所关联的 Fredholm 算符丛中定义等变截面，因此，按照文献\BookCite{ref026}的定义，代表 $K_{G,P}(X)$ 的一个元素。

\SourceParagraph{ch02-2.6.6-008}{100}
对于一般的 $E$，先加入一个由扭曲正则表示给出的满支集辅助丛 $F$，使每个稳定子群的所有不可约射影表示都出现。取 $\widehat{E}=(E\oplus F)\otimes\mathcal{H}$，并在 $E\otimes\mathcal{H}$ 上取 $\operatorname{id}_E\otimes T$，在 $F\otimes\mathcal{H}$ 上取恒等算符。所得等变 Fredholm 算符的核同构于 $E$，余核为零，故仍代表 $E$ 的 K 类；它同时保证了射影丛的稳定性，即使 $E$ 在某些点处的秩为零也没有问题。

\SourceParagraph{ch02-2.6.6-009}{100}
如果二上闭链 $b_o$ 是一个上边界，甚至只要 $(b_o)_x(g_1,g_2)$ 与 $x\in X$ 无关，这两种等变 K 理论定义的一致性就是显然的。再利用 Mayer--Vietoris 论证，就可以证明，对于任意 $b_o$，它们都一致。

\SourceParagraph{ch02-2.6.6-010}{100}
这里的关键在于：当 $X$ 与 $G$ 都有限时，扭曲类 $h$ 具有有限阶；于是，可以用核与余核的秩在各连通分支上保持不变的 Fredholm 算符族表示 K 类，从而用有限维的扭曲向量丛及其形式差表示这些类。

\section{附录：缝合定理的 Morse 理论证明}\label{sec:2.7}

\SourceParagraph{ch02-2.7-001}{100}
在这一附录中，用 Morse 理论统一证明四个定理。第一个是非常著名的结果：二维拓扑场论恰好由交换 Frobenius 代数编码。第二个是开闭弦理论的相应结果，也就是\BookSectionRef{sec:2.1}{2.1}中的\BookRef{thm:2.1}{定理2.1}。第三个与第四个，则是这两个结果的等变版本，即\BookSectionRef{sec:2.6}{2.6}中的\BookRef{thm:2.11}{定理2.11}与\BookRef{thm:2.15}{定理2.15}。

\subsection{经典定理}\label{sec:2.7.1}

\SourceParagraph{ch02-2.7.1-001}{101}
我们要证明：给定交换 Frobenius 代数 $\mathcal{C}$，就可以为从 $S_0$ 到 $S_1$ 的定向配边 $\Sigma$ 指定线性映射
\[
U_\Sigma\colon\mathcal{C}^{\otimes p}\longrightarrow\mathcal{C}^{\otimes q},
\]
其中，定向一维流形 $S_0,S_1$ 分别有 $p,q$ 个连通分支。

\SourceParagraph{ch02-2.7.1-002}{101}
\begin{proof}
总可以选取光滑函数 $f\colon\Sigma\to[0,1]\subset\mathbb{R}$，使 $f^{-1}(0)=S_0$、$f^{-1}(1)=S_1$，并且仅有“Morse”奇点。也就是说，微分 $\mathrm{d}f$ 只在有限多个点 $x_1,\ldots,x_n\in\Sigma$ 处为零，并且满足
\begin{enumerate}[label=(\roman*)]
\item 对于每个 $i$，Hessian $\mathrm{d}^2f(x_i)$ 是非退化二次型；
\item 临界值 $c_1=f(x_1),\ldots,c_n=f(x_n)$ 互不相同，且都不等于0或1。
\end{enumerate}
每个临界点都有一个指标，为0、1或2，等于 Hessian $\mathrm{d}^2f(x_i)$ 的负特征值的个数。

\SourceParagraph{ch02-2.7.1-003}{101}
函数 $f$ 的选择，将配边分解为“初等”配边。取
\[
0=t_0<c_1<t_1<c_2<t_2<\cdots<c_n<t_n=1,
\]
并令 $S_t=f^{-1}(t)$。那么，每个 $S_{t_i}$ 都是若干互不相交圆周的并，设其个数为 $m_i$，则 $m_i=m_{i-1}\pm1$；$\Sigma_i=f^{-1}([t_{i-1},t_i])$ 是从 $S_{t_{i-1}}$ 到 $S_{t_i}$ 的配边。除了一个连通分支具有\BookRef{fig:2.2}{图2}中四种形式之一外，其余部分都是平凡的，也就是圆柱的并。对于给定的 Frobenius 代数 $\mathcal{C}$，在每种情形下，都知道如何定义算符
\[
U_{\Sigma_i}\colon\mathcal{C}^{\otimes m_{i-1}}\longrightarrow\mathcal{C}^{\otimes m_i}.
\]
在第三种情形下，指定的映射为
\[
\phi\longmapsto\sum_i\phi\phi_i\otimes\phi^i,
\]
其中，$\{\phi_i\}$ 与 $\{\phi^i\}$ 是 $\mathcal{C}$ 的对偶基，满足 $\theta_{\mathcal{C}}(\phi^i\phi_j)=\delta_{ij}$。有两点需要注意。首先，必须要求 $\mathcal{C}$ 交换，否则就需要为裤面配边的两个入圆周指定顺序，而这里没有给定这样的顺序。其次，这些指定具有如下性质：反转配边中的时间方向，就会将算符替换为它关于态空间 Frobenius 内积的伴随算符。这一性质，将是后面所有构造遵循的基本原则；它也减少了以下繁复论证中需要检验的情形数目。

\SourceParagraph{ch02-2.7.1-004}{101}
现在，关键任务是证明，复合算符 $U_{\Sigma_n}\circ\cdots\circ U_{\Sigma_1}$ 与所选的 Morse 函数 $f$ 无关。

\begin{figure}[H]
\centering
\BookImage{01_Numbered_Figures/Chapter_02/Ch02_Fig31_PDFp102_Bookp092.pdf}
\BookNumberOnlyCaption\label{fig:2.31}
\end{figure}

\SourceParagraph{ch02-2.7.1-005}{102}
两个 Morse 函数 $f_0,f_1$，总可以用一条光滑路径 $\{f_s\}_{0\le s\le1}$ 连接。在有限多个参数值以外，$f_s$ 都是 Morse 函数；在这些例外参数值处，发生如下两种情形之一：
\begin{enumerate}[label=(\roman*)]
\item $f_s$ 有一个退化临界点，在该点附近的局部坐标 $(u,v)$ 中，具有形式 $f_s(u,v)=c\pm u^2+v^3$，其中 $c$ 为临界值；
\item $f_s$ 的两个不同临界点 $x_i,x_j$ 具有相同临界值 $f_s(x_i)=f_s(x_j)=c$。
\end{enumerate}
在第一种情形中，当参数越过这个非 Morse 参数值 $s$ 时，两个指标相邻的临界点产生或湮灭，配边按\BookRef{fig:2.31}{图31}所示方式改变，或按其反向方式改变，或按这些图示的时间反转改变。代数 $\mathcal{C}$ 中的恒等式 $1\cdot a=a$，保证了 $U_\Sigma$ 在这类改变下的良定义性。

\SourceParagraph{ch02-2.7.1-006}{102}
情形(ii)更棘手一些。由于形如 $U\otimes1$ 与 $1\otimes U'$ 的算符交换，容易看出，只有当两个临界点 $x_i,x_j$ 在“异常”的临界等值线 $S_c$ 中属于同一个连通分支时，才有需要证明的事情；在这种情况下，两者的指标都必须为1。

\SourceParagraph{ch02-2.7.1-007}{102}
考虑所得到的两步配边：在临界参数值 $s$ 的两侧，它以不同方式分解。这个配边只有一个非平凡连通分支；由于初等配边会使圆周的个数改变1，这个分支必定是从 $p$ 个圆周到 $q$ 个圆周的配边，其中 $(p,q)$ 为 $(1,1)$、$(2,2)$、$(1,3)$ 或 $(3,1)$。只需检验 $(1,3)$ 与 $(3,1)$ 中的一种，因为二者仅差时间反转。配边的 Euler 数，等于其 Morse 函数的临界点个数，按符号 $(-1)^{\text{指标}}$ 计数；因此，这个非平凡分支的 Euler 数为 $-2$。当 $(p,q)=(1,1)$ 时，它是有两个孔洞的环面；其余情形中，它是有四个孔洞的球面。

\SourceParagraph{ch02-2.7.1-008}{102}
在\BookRef{fig:2.32}{图32}所示的 $(1,1)$ 情形中，一个圆周分裂为两个，然后再次合并。这一情形无需另作检验：虽然有两个孔洞的环面，可以按许多不同同痕类的切法，分成两个裤面，但可能得到的复合配边只有一种，而可能得到的复合映射也只有 $\mathcal{C}\to\mathcal{C}\otimes\mathcal{C}\to\mathcal{C}$ 这一种。

\begin{figure}[H]
\centering
\BookImage{01_Numbered_Figures/Chapter_02/Ch02_Fig32_PDFp103_Bookp093.pdf}
\caption{图中表示被两个水平平面截去两端的环面，这些平面是 Morse 函数的等值面。两个临界点分别位于内圆圈的顶部与底部。如果把环面作为刚体倾斜，就可以使高度函数的这两个临界点处于同一水平面内。}\label{fig:2.32}
\end{figure}

\SourceParagraph{ch02-2.7.1-009}{102}
在 $(3,1)$ 情形中，三个圆周中的两个先合并，所得圆周再与第三个合并。图示见\BookRef{fig:2.33}{图33}。显然，$\mathcal{C}$ 中的结合律涵盖了这一情形。

\begin{figure}[H]
\centering
\BookImage{01_Numbered_Figures/Chapter_02/Ch02_Fig33_PDFp103_Bookp093.pdf}
\BookNumberOnlyCaption\label{fig:2.33}
\end{figure}

\begin{figure}[H]
\centering
\BookImage{01_Numbered_Figures/Chapter_02/Ch02_Fig34_PDFp103_Bookp093.pdf}
\caption{将具有两个临界点的配边嵌入 $\mathbb{R}^3$ 的不同方式。最右边的图示，对应于\BookRef{fig:2.35}{图35}(i)所绘等值线的情形。}\label{fig:2.34}
\end{figure}

\SourceParagraph{ch02-2.7.1-010}{103}
在 $(2,2)$ 情形中，再次将一个有四个孔洞的球面分解为两个初等配边。由\BookRef{fig:2.34}{图34}可见，这可以按许多方式完成。为确保没有遗漏任何可能性，最好的办法，是考虑略低于那个同时含有两个临界点的临界等值面的等值线。如果这条等值线由两个圆周组成，就必须具有\BookRef{fig:2.35}{图35}中 (i) 或 (ii) 两种形式之一。（考察虚线圆圈内部各“端点”的各种可能连接方式。）

\begin{figure}[H]
\centering
\BookImage{01_Numbered_Figures/Chapter_02/Ch02_Fig35_PDFp104_Bookp094.pdf}
\BookNumberOnlyCaption\label{fig:2.35}
\end{figure}

\SourceParagraph{ch02-2.7.1-011}{104}
不过，无论发生哪种情形，这个配边可能导出的代数映射，只有
\[
\mathcal{C}\otimes\mathcal{C}\longrightarrow\mathcal{C}\longrightarrow\mathcal{C}\otimes\mathcal{C}
\]
与
\[
\mathcal{C}\otimes\mathcal{C}\longrightarrow\mathcal{C}\otimes\mathcal{C}\otimes\mathcal{C}
\longrightarrow\mathcal{C}\otimes\mathcal{C},
\]
分别由
\[
\phi\otimes\phi'\longmapsto\phi\phi'\longmapsto\sum_i\phi\phi'\phi_i\otimes\phi^i
\]
与
\[
\phi\otimes\phi'\longmapsto\sum_i\phi\phi_i\otimes\phi^i\otimes\phi'
\longmapsto\sum_i\phi\phi_i\otimes\phi^i\phi'
\]
给出，其中，$\{\phi_i\}$ 与 $\{\phi^i\}$ 是 $\mathcal{C}$ 的对偶基，满足 $\theta_{\mathcal{C}}(\phi^i\phi_j)=\delta_{ij}$。这两个映射相等，是因为恒等式
\begin{equation}\label{eq:2.111}
\sum_i\phi'\phi_i\otimes\phi^i=\sum_i\phi_i\otimes\phi^i\phi'
\end{equation}
成立。这个恒等式在任意 Frobenius 代数中都成立；在这里，两端与 $\phi^j\otimes\phi_k$ 的内积，都是 $\theta_{\mathcal{C}}(\phi^j\phi'\phi_k)$。

\SourceParagraph{ch02-2.7.1-012}{104}
这就完成了定理的证明。注意，论证中用到了交换 Frobenius 代数的全部公理。
\end{proof}

\subsection{开闭弦理论}\label{sec:2.7.2}

\SourceParagraph{ch02-2.7.2-001}{104}
\begin{proof}
给定\BookSectionRef{sec:2.1.2}{2.1.2}中的开闭弦代数数据，包含 Frobenius 范畴 $\mathcal{B}$、闭弦代数 $\mathcal{C}$ 及满足该节条件的转变映射。与前面的论证一样，考虑从 $S_0$ 到 $S_1$ 的配边 $\Sigma$；但现在，$S_0,S_1$ 都是圆周与区间的并，而边界 $\partial\Sigma$ 有一个受约束部分 $\partial_{\mathrm{constr}}\Sigma$，简记为 $\partial'\Sigma$，它是从 $\partial S_0$ 到 $\partial S_1$ 的配边。仍选取函数 $f\colon\Sigma\to[0,1]$，但现在 $f$ 有两类临界点：$\Sigma$ 内部满足 $\mathrm{d}f=0$ 的点，以及 $\partial'\Sigma$ 上 $f$ 对边界的限制的微分为零的点。对于内部临界点，“非退化”具有通常的含义。
\begin{definition*}\label{def:2.boundary-critical}
边界上的临界点 $x$ 称为非退化的，如果它是 $f|_{\partial'\Sigma}$ 的非退化临界点，并且 $f$ 在 $x$ 处沿边界法向的导数不为零。
\end{definition*}

\SourceParagraph{ch02-2.7.2-002}{105}
\begin{definition*}\label{def:2.boundary-morse}
与前面一样，如果 $f$ 的所有临界点都非退化，所有临界值都互不相同，且均不等于0或1，就称 $f$ 为 Morse 函数。
\end{definition*}
总可以选取这样的函数。

\SourceParagraph{ch02-2.7.2-003}{105}
现在有四类边界临界点，记为 $0\pm$、$1\pm$，分别记录指标与朝内法向导数的符号。越过其中一个临界点时，可以发生六种过程。在 $0+$ 或 $1-$ 类型处，一条开弦产生或湮灭。在 $0-$ 类型处，或者两条开弦首尾相接，或者一条开弦变为闭弦。$1+$ 类型是 $0-$ 类型的时间反转。给定 Frobenius 范畴 $\mathcal{B}$，就知道在这六种情形中分别应当做什么。

\begin{figure}[H]
\centering
\BookImage{01_Numbered_Figures/Chapter_02/Ch02_Fig36_PDFp105_Bookp095.pdf}
\BookNumberOnlyCaption\label{fig:2.36}
\end{figure}

\SourceParagraph{ch02-2.7.2-004}{105}
与前面一样，内部临界点的指标为0、1或2。只有当指标为1时，相应配边才可能涉及开弦。在时间反转意义下，指标1的过程有三种：两条闭弦变为一条；一条开弦“吸收”一条闭弦；两条开弦“重新组合”，成为两条新的开弦，如\BookRef{fig:2.36}{图36}所示。

\SourceParagraph{ch02-2.7.2-005}{105}
对于给定的 Frobenius 范畴 $\mathcal{B}$，为“开弦 $+$ 闭弦 $\to$ 开弦”指定映射
\[
\mathcal{O}_{ab}\otimes\mathcal{C}\longrightarrow\mathcal{O}_{ab},
\]
由 $\phi\otimes\psi\mapsto\phi\psi$ 给出。这里与通常一样，将 $\mathcal{O}_{ab}$ 看作 $\mathcal{C}$ 模，并写为
\[
\phi\psi=\iota_a(\psi)\phi=\phi\iota_b(\psi).
\]
为“开弦 $+$ 开弦 $\to$ 开弦 $+$ 开弦”指定映射
\[
\mathcal{O}_{ab}\otimes\mathcal{O}_{cd}\longrightarrow\mathcal{O}_{ad}\otimes\mathcal{O}_{cb},
\]
由
\[
\psi\otimes\psi'\longmapsto\sum_i\psi\psi_i\otimes\psi'\psi^i
\]
给出，其中，$\psi_i$ 与 $\psi^i$ 分别是 $\mathcal{O}_{bd}$ 与 $\mathcal{O}_{db}$ 的对偶基。

\SourceParagraph{ch02-2.7.2-006}{105}
现在必须考虑，改变 Morse 函数时会发生什么。与前面一样，两个 Morse 函数可以通过路径 $\{f_s\}$ 连接，除有限多个参数值 $s$ 外，每个 $f_s$ 都是 Morse 函数；在这些例外值处，或者一个临界点退化，或者两个临界值相同。先考虑退化的情形。现在必须允许三类退化：除了与闭弦情形相同的内部退化，还有两类边界退化，即 $f|_{\partial'\Sigma}$ 具有一个三次拐点，或者某个边界临界点处的法向导数为零。

\begin{figure}[H]
\centering
\BookImage{01_Numbered_Figures/Chapter_02/Ch02_Fig37_PDFp106_Bookp096.pdf}
\BookNumberOnlyCaption\label{fig:2.37}
\end{figure}

\SourceParagraph{ch02-2.7.2-007}{106}
当 $s$ 越过边界拐点所对应的参数值时，两个指标分别为0与1、法向导数符号相同的非退化边界临界点产生或湮灭。这意味着，配边在\BookRef{fig:2.37}{图37}的两个图示之间改变，或者按其时间反转改变。范畴 $\mathcal{B}$ 具有恒等态射这一公理，涵盖了这些改变。

\SourceParagraph{ch02-2.7.2-008}{106}
当一个边界临界点处的法向导数为零时，发生的是一个内部临界点“穿越 $\Sigma$ 的边界”：它与一个边界临界点重合，并改变该点处法向导数的符号。共有四种情形：
\[
\begin{aligned}
(0-)+(\text{指标0})&\longrightarrow(0+),\\
(0+)+(\text{指标1})&\longrightarrow(0-),
\end{aligned}
\]
以及这两种情形的时间反转。在第一种情形中，先产生一条小闭弦、再将它切开的复合配边，被产生一条开弦的初等配边替代。这对应于映射 $\mathcal{C}\to\mathcal{O}_{aa}$ 将 $1_{\mathcal{C}}$ 映为 $1_a$ 的公理。在第二种情形中，复合配边先产生一条开弦，然后它或者“吸收”一条已有闭弦，或者与一条已有开弦“重新组合”；这两个复合配边，应当分别等价于将闭弦或开弦切开的初等配边。在上面的公式中令 $\psi=1_a$，就可以看出，Frobenius 范畴公理允许这些等价关系。

\SourceParagraph{ch02-2.7.2-009}{106}
当内部临界点退化时，在时间反转意义下，发生的是先产生一条闭弦，再与一条已有开弦或闭弦结合；这应当等价于平凡配边。单位公理再次涵盖了这些情形。

\SourceParagraph{ch02-2.7.2-010}{106}
最后，必须考虑两个临界值相互越过时会发生什么。它们可能来自两个边界临界点、两个内部临界点，或者各来自一类。

\SourceParagraph{ch02-2.7.2-011}{106}
如果两个边界临界点由一条临界等值线相连，那么，等值线具有\BookRef{fig:2.38}{图38}所示形式。这给出四种需要检验的情形，其中，临界等值面之下的等值线，如\BookRef{fig:2.39}{图39}所示。

\begin{figure}[H]
\centering
\BookImage{01_Numbered_Figures/Chapter_02/Ch02_Fig38_PDFp107_Bookp097.pdf}
\BookNumberOnlyCaption\label{fig:2.38}
\end{figure}

\begin{figure}[H]
\centering
\BookImage{01_Numbered_Figures/Chapter_02/Ch02_Fig39_PDFp107_Bookp097.pdf}
\BookNumberOnlyCaption\label{fig:2.39}
\end{figure}

\SourceParagraph{ch02-2.7.2-012}{106}
情形(i)a由范畴 $\mathcal{B}$ 中复合的结合律涵盖；情形(i)b由恒等式\BookEqRef{eq:2.111}{2.111}的开弦版本涵盖；情形(ii)a由迹关系 $\iota^a(\psi_1\psi_2)=\iota^b(\psi_2\psi_1)$ 涵盖，这一关系可由\BookEqRef{eq:2.5}{2.5}、\BookEqRef{eq:2.9}{2.9}与\BookEqRef{eq:2.10}{2.10}联合推出；情形(ii)b则由 Cardy 恒等式涵盖。

\SourceParagraph{ch02-2.7.2-013}{107}
如果一个边界临界点与一个内部临界点处于同一等值面上，不妨假设，边界临界点的类型为 $0-$，内部临界点的指标为1，并且二者在临界等值线中相连；这条等值线必定具有\BookRef{fig:2.40}{图40}中的四种形式之一。

\begin{figure}[H]
\centering
\BookImage{01_Numbered_Figures/Chapter_02/Ch02_Fig40_PDFp108_Bookp098.pdf}
\BookNumberOnlyCaption\label{fig:2.40}
\end{figure}

\SourceParagraph{ch02-2.7.2-014}{107}
在边界临界点处，或者一条开弦变为闭弦，或者两条开弦结合。依次考察这两种可能性。在第一种情形中，如果先遇到边界临界点，那么在内部临界点处可以发生三种过程：这条闭弦分裂为两条闭弦，或者与另一条闭弦或开弦结合。因此，可能性如下；这里以 $o,c$ 分别表示开弦与闭弦：
\[
\begin{aligned}
o&\longrightarrow c\longrightarrow c+c,\\
o+c&\longrightarrow c+c\longrightarrow c,\\
o+o&\longrightarrow c+o\longrightarrow o.
\end{aligned}
\]
如果先遇到内部临界点，每种情形都只有一种可能性，上面三个过程分别被替换为
\[
\begin{aligned}
o&\longrightarrow o+c\longrightarrow c+c,\\
o+c&\longrightarrow o\longrightarrow c,\\
o+o&\longrightarrow o+o\longrightarrow o.
\end{aligned}
\]
需要检验三个恒等式。前两个，归结为 $\iota^a\colon\mathcal{O}_{aa}\to\mathcal{C}$ 是 $\mathcal{C}$ 模映射这一事实。第三个，就是 Cardy 条件。

\SourceParagraph{ch02-2.7.2-015}{108}
现在考虑，两条开弦在边界临界点处结合的情形。如果先遇到边界临界点，那么在内部临界点处仍然可以发生三种过程：所得开弦发射一条闭弦，或者与另一条闭弦或开弦相互作用。可能性为
\[
\begin{aligned}
o+o&\longrightarrow o\longrightarrow o+c,\\
o+o+c&\longrightarrow o+c\longrightarrow o,\\
o+o+o&\longrightarrow o+o\longrightarrow o+o.
\end{aligned}
\]
在第二种与第三种情形中，如果反转临界点的顺序，只有一种过程可能发生；它们变为
\[
\begin{aligned}
o+o+c&\longrightarrow o+o\longrightarrow o,\\
o+o+o&\longrightarrow o+o+o\longrightarrow o+o.
\end{aligned}
\]
这两种情形分别给出类型为 $\mathcal{O}_{ab}\otimes\mathcal{O}_{bc}\otimes\mathcal{C}\to\mathcal{O}_{ac}$ 与 $\mathcal{O}_{ab}\otimes\mathcal{O}_{bc}\otimes\mathcal{O}_{de}\to\mathcal{O}_{ae}\otimes\mathcal{O}_{dc}$ 的代数映射。对每种情形，比较两种顺序所需的恒等式，都立即成立。

\SourceParagraph{ch02-2.7.2-016}{108}
不过，第一个过程可能变为
\[
o+o\longrightarrow o+o+c\longrightarrow o+c,
\]
或者
\[
o+o\longrightarrow o+o\longrightarrow o+c.
\]
第一种过程在代数上没有需要特别讨论的内容；为了处理第二种过程，则需要检验
\[
\sum_i\psi\psi'\phi_i\otimes\phi^i
=\sum_k\psi\psi^k\otimes\iota^b(\psi'\psi_k),
\]
其中，$\psi\in\mathcal{O}_{ab}$、$\psi'\in\mathcal{O}_{bc}$；$\phi^i,\phi_i$ 是 $\mathcal{C}$ 的对偶基，而 $\psi^k,\psi_k$ 分别是 $\mathcal{O}_{bc},\mathcal{O}_{cb}$ 的对偶基。这一关系成立，因为左端与 $\psi_m\otimes\phi_j$ 的内积为 $\theta_a(\psi\psi'\phi_j\psi_m)$，其中 $\psi_m\in\mathcal{O}_{ca}$；右端与 $\psi_m\otimes\phi_j$ 的内积则为
\begin{equation}\label{eq:2.112}
\begin{aligned}
&\sum_k\theta_a(\psi\psi^k\psi_m)\theta_{\mathcal{C}}\bigl(\iota^b(\psi'\psi_k)\phi_j\bigr)\\
&\qquad=\sum_k\theta_c(\psi_m\psi\psi^k)\theta_c(\psi_k\phi_j\psi')\\
&\qquad=\theta_c(\psi_m\psi\phi_j\psi')
=\theta_a(\psi\psi'\phi_j\psi_m).
\end{aligned}
\end{equation}

\begin{figure}[H]
\centering
\BookImage{01_Numbered_Figures/Chapter_02/Ch02_Fig41_PDFp109_Bookp099.pdf}
\BookNumberOnlyCaption\label{fig:2.41}
\end{figure}

\SourceParagraph{ch02-2.7.2-017}{109}
最后，必须考虑两个内部临界点处于同一等值面上的情形。这里有闭弦情形中已经讨论过的可能性，但还必须允许所涉及的弦中任意若干条，乃至全部，都是开弦。可以按照略低于同时含有两个临界点的等值面、并经过临界点附近的那部分等值线的连通分支数目，分析这一情形。这样的分支必定有一、二或三个。如果有三个，除三条弦全是闭弦的情形外，可以构成\BookRef{fig:2.41}{图41}所示的五种构型。在所有这些情形中，Frobenius 范畴中复合的结合律，立即保证了复合映射的良定义性。

\begin{figure}[H]
\centering
\BookImage{01_Numbered_Figures/Chapter_02/Ch02_Fig42_PDFp110_Bookp100.pdf}
\BookNumberOnlyCaption\label{fig:2.42}
\end{figure}

\SourceParagraph{ch02-2.7.2-018}{110}
如果临界等值面之下有两个分支，那么仍然可以构成五种构型，如\BookRef{fig:2.42}{图42}所示：或者这两个分支相遇两次，或者它们相遇一次，而其中一个分支还与自身相互作用。不过，只需检验三种情形，因为第二种是\BookRef{fig:2.41}{图41}中某种情形的时间反转，而最后两种互为时间反转。\BookRef{fig:2.42}{图42}(i)对应于如下事实：复合
\[
\mathcal{O}_{ab}\otimes\mathcal{C}\longrightarrow\mathcal{O}_{ab}
\longrightarrow\mathcal{O}_{ab}\otimes\mathcal{C}
\]
可以通过以不同方式切开复合配边而实现；但这里没有需要检验的事情，因为可能的代数映射只有一种。

\SourceParagraph{ch02-2.7.2-019}{111}
在\BookRef{fig:2.42}{图42}的情形(iii)中，临界点的一种顺序给出与前面相同的复合
\[
\mathcal{O}_{ab}\otimes\mathcal{C}\longrightarrow\mathcal{O}_{ab}
\longrightarrow\mathcal{O}_{ab}\otimes\mathcal{C},
\]
另一种顺序则给出
\[
\mathcal{O}_{ab}\otimes\mathcal{C}\longrightarrow
\mathcal{O}_{ab}\otimes\mathcal{C}\otimes\mathcal{C}
\longrightarrow\mathcal{O}_{ab}\otimes\mathcal{C}.
\]
不过，使用前面的记号，很容易验证，两个映射都将 $\psi\otimes\phi$ 映为 $\sum_i\psi\phi\phi_i\otimes\phi^i$。

\SourceParagraph{ch02-2.7.2-020}{111}
在\BookRef{fig:2.42}{图42}的情形(iv)中，需要再次比较复合
\[
\mathcal{O}_{ab}\otimes\mathcal{C}\longrightarrow
\mathcal{O}_{ab}\otimes\mathcal{C}\otimes\mathcal{C}
\longrightarrow\mathcal{O}_{ab}\otimes\mathcal{C}
\]
与
\[
\mathcal{O}_{ab}\otimes\mathcal{C}\longrightarrow\mathcal{O}_{ab}
\longrightarrow\mathcal{O}_{ab}\otimes\mathcal{C}.
\]
这一次，需要检验
\[
\sum_i\psi\phi_i\otimes\phi^i\phi=\sum_i\psi\phi\phi_i\otimes\phi^i.
\]
这正是闭弦理论讨论结束时遇到的同一个公式。

\SourceParagraph{ch02-2.7.2-021}{111}
最后，假设临界等值面之下的等值线只有一个连通分支。有三种可能构型，对应于将区间上四个点两两配对的三种方式，如\BookRef{fig:2.43}{图43}所示。前两种，是已经处理过的情形的时间反转。最后一种，无论临界点采用哪种顺序，都会导出分解
\[
\mathcal{O}_{ab}\longrightarrow\mathcal{O}_{ab}\otimes\mathcal{C}
\longrightarrow\mathcal{O}_{ab}.
\]
这只有一种可能性，因此，没有需要检验的事情。

\begin{figure}[H]
\centering
\BookImage{01_Numbered_Figures/Chapter_02/Ch02_Fig43_PDFp111_Bookp101.pdf}
\BookNumberOnlyCaption\label{fig:2.43}
\end{figure}

\SourceParagraph{ch02-2.7.2-022}{112}
这就完成了开闭弦理论的定理的证明。
\end{proof}

\subsection{等变闭弦理论}\label{sec:2.7.3}

\SourceParagraph{ch02-2.7.3-001}{113}
现在必须重新进行本附录第一部分的讨论，但要为曲面与圆周配备主 $G$ 丛，其中 $G$ 是给定的有限群。

\SourceParagraph{ch02-2.7.3-002}{113}
\begin{proof}
首先注意，任意配备丛的圆周，都同构于标准对象 $S_g$：它是标准圆周 $S^1$ 上平行移动为 $g\in G$ 的标准丛。此外，从 $S_g$ 到 $S_{g'}$ 的态射集合，是 $\{h\in G:hgh^{-1}=g'\}$。换言之，$S^1$ 上的丛范畴，等价于 $G$ 以共轭方式作用于自身所形成的范畴 $G//G$。因此，等变理论为每个 $g$ 给出向量空间 $\mathcal{C}_g$，这些空间一起构成 $G$ 上的 $G$ 等变向量丛。反过来，给定 $G$ 等变向量丛 $\{\mathcal{C}_g\}$，以及带有丛 $P$ 的圆周 $S$，理论就给出向量空间 $\mathcal{H}(S,P)$。它的元素是如下法则：为每个 $x\in S$ 与平凡化 $t\colon P_x\to G$，指定 $\psi_{x,t}\in\mathcal{C}_{g_{x,t}}$，其中 $g_{x,t}$ 是以 $(x,t)$ 为基点数据的 $P$ 的平行移动。要求
\[
\psi_{x',t'}=g\cdot\psi_{x,t},
\]
其中，$g$ 是 $P$ 沿从 $(x,t)$ 到 $(x',t')$ 的正向路径的平行移动。为了使这个定义良好，必须要求 $g_{x,t}$ 在 $\mathcal{C}_{g_{x,t}}$ 上的作用平凡；这一条件的必要性，已经在\BookSectionRef{sec:2.6}{2.6}中说明。

\SourceParagraph{ch02-2.7.3-003}{113}
接着，考虑从 $S_g$ 到 $S_{g'}$ 的平凡配边。将两端的丛延拓到圆柱上的各种可能方式，对应于从入基点到出基点的各种可能平行移动，也就是 $G//G$ 中的态射集合 $\{h\in G:hgh^{-1}=g'\}$。显然，这些圆柱诱导已经知道的同构 $\mathcal{C}_g\to\mathcal{C}_{g'}$。不过，如果两个这样的配边之间，存在一个保持丛、并在两端为恒等映射的圆柱微分同胚，就将它们视为等价。圆柱的映射类群，由绕圆柱的 Dehn 扭转生成；因此，与 $h$ 对应的态射，等价于与 $hg=g'h$ 对应的态射。这意味着，$g$ 在 $\mathcal{C}_g$ 上的作用必须平凡，正如前面已经知道的那样。

\SourceParagraph{ch02-2.7.3-004}{113}
现在考察\BookRef{fig:2.2}{图2}中四种初等配边上可能的丛。帽面上的丛当然必须平凡。这里所需的裤面配边，是由一个临界等值面分开的、相邻等值线之间的区域。可以将它们画成\BookRef{fig:2.44}{图44}的形式，其中实线等值线位于临界等值面之下，虚线等值线位于其上。在临界点的邻域，也就是着色区域内，可以将 $G$ 丛平凡化；随后，沿带状区域，也就是曲面的未着色部分，给定平行移动 $g_1,g_2$，就确定了配边上的丛。为左图的情形(i)指定的算符，是\BookEqRef{eq:2.80}{2.80}中的乘法映射
\[
m_{g_1,g_2}\colon\mathcal{C}_{g_1}\otimes\mathcal{C}_{g_2}\longrightarrow\mathcal{C}_{g_1g_2}.
\]
这样书写时，就为两条带状区域选择了顺序，也就是在出环路上选择了一个基点。两种顺序通过共轭映射
\[
\alpha_{g_2}\colon\mathcal{C}_{g_1g_2}\longrightarrow\mathcal{C}_{g_2g_1}
\]
关联。因此，为使指定的算符良好定义，所需的相容性条件为
\[
m_{g_2,g_1}(\psi_2\otimes\psi_1)
=\alpha_{g_2}\bigl(m_{g_1,g_2}(\psi_1\otimes\psi_2)\bigr).
\]
将\BookEqRef{eq:2.82}{2.82}与如下事实结合，就可看出，它在任意 Turaev 代数中都成立：$G$ 通过代数自同构作用，且 $\alpha_{g_2}$ 在 $\mathcal{C}_{g_2}$ 上的作用平凡。裤面的映射类群，由平行于其三个边界圆周的三次 Dehn 扭转生成，所以，为使裤面算符的指定良好定义，无需再添加新的条件。

\begin{figure}[H]
\centering
\BookImage{01_Numbered_Figures/Chapter_02/Ch02_Fig44_PDFp112_Bookp102.pdf}
\BookNumberOnlyCaption\label{fig:2.44}
\end{figure}

\SourceParagraph{ch02-2.7.3-005}{114}
与\BookRef{fig:2.15}{图15}的时间反转配边对应的线性映射
\[
c_{g_1,g_2}\colon\mathcal{C}_{g_1g_2}\longrightarrow\mathcal{C}_{g_1}\otimes\mathcal{C}_{g_2},
\]
由伴随要求确定；这里要记住，$\mathcal{C}_g$ 的对偶空间是 $\mathcal{C}_{g^{-1}}$。这个映射为
\[
c_{g_1,g_2}(\phi)=\sum_i\phi\phi^i\otimes\phi_i,
\]
其中，$\{\phi_i\}$ 是 $\mathcal{C}_{g_2}$ 的一组基，$\{\phi^i\}$ 是 $\mathcal{C}_{g_2^{-1}}$ 中的对偶基。

\SourceParagraph{ch02-2.7.3-006}{114}
任意带有丛的配边，都可以像前面一样，利用 Morse 理论分解；初等配边继承相应的丛。讨论的困难部分，是改变 Morse 函数时会发生什么。不过，事实上，唯一需要实质性检验的步骤，是交换位于同一等值面上的两个指标1临界点，也就是\BookRef{fig:2.32}{图32}、\BookRef{fig:2.33}{图33}、\BookRef{fig:2.34}{图34}中的配边。

\SourceParagraph{ch02-2.7.3-007}{114}
先考虑\BookRef{fig:2.32}{图32}的情形：一条弦分裂后再合并，也就是一个有两个孔洞的环面，一个孔洞为入边界，另一个为出边界。采用\BookRef{fig:2.45}{图45}的画法。没有采用看似更清楚的\BookRef{fig:2.46}{图46}形式，是因为那样会使两个临界点的邻域在平面内具有相反定向。

\begin{figure}[H]
\centering
\BookImage{01_Numbered_Figures/Chapter_02/Ch02_Fig45_PDFp112_Bookp102.pdf}
\BookNumberOnlyCaption\label{fig:2.45}
\end{figure}

\begin{figure}[H]
\centering
\BookImage{01_Numbered_Figures/Chapter_02/Ch02_Fig46_PDFp112_Bookp102.pdf}
\BookNumberOnlyCaption\label{fig:2.46}
\end{figure}

\SourceParagraph{ch02-2.7.3-008}{114}
这个配边对应于映射 $\mathcal{C}_{4321}\to\mathcal{C}_{2341}$；这里与下文一样，将 $\mathcal{C}_{g_4g_3g_2g_1}$ 简记为 $\mathcal{C}_{4321}$，并将 $\alpha_{g_3}$ 简记为 $\alpha_3$，依此类推。如果先遇到左边的临界点，得到的映射为
\[
\mathcal{C}_{4321}\longrightarrow\mathcal{C}_{43}\otimes\mathcal{C}_{21}
\cong\mathcal{C}_{34}\otimes\mathcal{C}_{12}\longrightarrow\mathcal{C}_{3412}
\cong\mathcal{C}_{2341},
\]
\[
\begin{aligned}
\phi&\longmapsto\sum_i\phi\phi^i\otimes\phi_i
\longmapsto\sum_i\alpha_3(\phi\phi^i)\otimes\alpha_1(\phi_i)\\
&\longmapsto\sum_i\alpha_3(\phi\phi^i)\alpha_1(\phi_i)
\longmapsto\sum_i\alpha_2\bigl(\alpha_3(\phi\phi^i)\alpha_1(\phi_i)\bigr),
\end{aligned}
\]
其中，$\phi_i$ 遍历 $\mathcal{C}_{21}$ 的一组基。（前一行中标为 $\cong$ 的映射，对应于移动各条弦上所选的基点。）

\SourceParagraph{ch02-2.7.3-009}{115}
采用另一种顺序，则得到
\[
\mathcal{C}_{4321}\cong\mathcal{C}_{3214}\longrightarrow\mathcal{C}_{32}\otimes\mathcal{C}_{14}
\cong\mathcal{C}_{23}\otimes\mathcal{C}_{41}\longrightarrow\mathcal{C}_{2341},
\]
\[
\begin{aligned}
\phi&\longmapsto\alpha_4^{-1}(\phi)
\longmapsto\sum_i\alpha_4^{-1}(\phi)\psi^i\otimes\psi_i\\
&\longmapsto\sum_i\alpha_2\bigl(\alpha_4^{-1}(\phi)\psi^i\bigr)\otimes\alpha_4(\psi_i)\\
&\longmapsto\sum_i\alpha_2\bigl(\alpha_4^{-1}(\phi)\psi^i\bigr)\alpha_4(\psi_i),
\end{aligned}
\]
其中，$\psi_i$ 遍历 $\mathcal{C}_{14}$ 的一组基。

\SourceParagraph{ch02-2.7.3-010}{115}
因此，必须证明
\[
\sum_i\alpha_{23}(\phi\phi^i)\phi_i
=\sum_i\alpha_{24^{-1}}(\phi)\alpha_2(\psi^i)\alpha_4(\psi_i).
\]
可以从\BookSectionRef{sec:2.6}{2.6}的公理\BookEqRef{eq:2.84}{2.84}推出这一点，取 $h=g_2g_4^{-1}g_1^{-1}g_2^{-1}$、$g=g_1^{-1}g_2^{-1}$，论证如下。将方程右端改写为
\[
\sum_i\alpha_{24^{-1}}(\phi)\eta_i\alpha_g(\eta^i),
\]
其中，$\eta_i$ 是 $\mathcal{C}_h$ 中的基 $\alpha_2(\psi^i)$，从而 $\eta^i=\alpha_2(\psi_i)$，且
\[
\alpha_g(\eta^i)=\alpha_1^{-1}(\psi_i)=\alpha_4(\psi_i).
\]
由该公理，上式等于
\[
\sum_i\alpha_{24^{-1}}(\phi)\alpha_h(\xi_i)\xi^i
=\sum_i\alpha_{24^{-1}}(\phi)\alpha_h(\phi^i)\phi_i.
\]
最后，
\[
\alpha_{24^{-1}}(\phi)\alpha_h(\phi^i)
=\alpha_{24^{-1}}(\phi\phi^i)=\alpha_{23}(\phi\phi^i),
\]
因为 $\phi\phi^i\in\mathcal{C}_{43}$，所以
\[
\alpha_{24^{-1}}(\phi\phi^i)
=\alpha_{24^{-1}}\alpha_{43}(\phi\phi^i)=\alpha_{23}(\phi\phi^i).
\]
由此，\BookRef{fig:2.32}{图32}的情形就处理完了。

\SourceParagraph{ch02-2.7.3-011}{115}
事实上，这一情形显然是所有情形中最复杂的。再处理一种，即\BookRef{fig:2.35}{图35}的情形(i)，其中，两条弦先结合，再分裂。画成\BookRef{fig:2.47}{图47}的形式，它对应于两个复合：
\[
\begin{aligned}
\mathcal{C}_{43}\otimes\mathcal{C}_{21}
&\longrightarrow\mathcal{C}_{4321}\cong\mathcal{C}_{1432}
\longrightarrow\mathcal{C}_{14}\otimes\mathcal{C}_{32}\cong\mathcal{C}_{41}\otimes\mathcal{C}_{23},\\
\mathcal{C}_{43}\otimes\mathcal{C}_{21}
&\cong\mathcal{C}_{34}\otimes\mathcal{C}_{12}
\longrightarrow\mathcal{C}_{3412}\cong\mathcal{C}_{4123}
\longrightarrow\mathcal{C}_{41}\otimes\mathcal{C}_{23}.
\end{aligned}
\]

\begin{figure}[H]
\centering
\BookImage{01_Numbered_Figures/Chapter_02/Ch02_Fig47_PDFp112_Bookp102.pdf}
\BookNumberOnlyCaption\label{fig:2.47}
\end{figure}

\SourceParagraph{ch02-2.7.3-012}{115}
第一个过程给出
\[
\begin{aligned}
\psi\otimes\psi'&\longmapsto\psi\psi'
\longmapsto\alpha_1(\psi\psi')
\longmapsto\sum_i\alpha_1(\psi\psi')\phi^i\otimes\phi_i\\
&\longmapsto\sum_i\alpha_4\bigl(\alpha_1(\psi\psi')\phi^i\bigr)\otimes\alpha_2(\phi_i),
\end{aligned}
\]
其中，$\phi_i$ 是 $\mathcal{C}_{32}$ 的一组基。第二个过程给出
\[
\begin{aligned}
\psi\otimes\psi'&\longmapsto\alpha_3(\psi)\otimes\alpha_1(\psi')
\longmapsto\alpha_3(\psi)\alpha_1(\psi')\\
&\longmapsto\psi\alpha_{3^{-1}1}(\psi')
\longmapsto\sum_i\psi\alpha_{3^{-1}1}(\psi')\psi^i\otimes\psi_i,
\end{aligned}
\]
其中，$\psi_i$ 是 $\mathcal{C}_{23}$ 的一组基。

\SourceParagraph{ch02-2.7.3-013}{116}
不过，可以假设 $\psi_i=\alpha_2(\phi_i)$，从而 $\psi^i=\alpha_2(\phi^i)$。注意，$\alpha_1(\psi\psi')\phi^i\in\mathcal{C}_{14}$，因此
\[
\alpha_4\bigl(\alpha_1(\psi\psi')\phi^i\bigr)
=\alpha_1^{-1}\bigl(\alpha_1(\psi\psi')\phi^i\bigr).
\]
于是，只需证明
\[
\psi'\alpha_1^{-1}(\phi^i)
=\alpha_3^{-1}\bigl(\alpha_1(\psi')\phi^i\bigr).
\]
这一点成立，是因为 $\alpha_1(\psi')\phi^i\in\mathcal{C}_{13^{-1}}$，所以它被 $\alpha_{13^{-1}}$ 固定。

\SourceParagraph{ch02-2.7.3-014}{116}
其余验证留给读者。
\end{proof}

\subsection{等变开闭弦理论}\label{sec:2.7.4}

\SourceParagraph{ch02-2.7.4-001}{116}
现在必须重新处理开闭弦情形，并考虑配边上的 $G$ 丛。

\SourceParagraph{ch02-2.7.4-002}{116}
\begin{proof}
取\BookSectionRef{sec:2.6.4}{2.6.4}中的等变开闭弦代数数据。为一条从 $b$ 到 $a$、且所配丛带有平凡化的开弦，指定向量空间 $\mathcal{O}_{ab}$。通过元素 $g\in G$ 改变平凡化，对应于 $g$ 在 $\mathcal{O}_{ab}$ 上的作用 $\rho_g$；它也对应于一个矩形配边诱导的映射，沿矩形的受约束边的平行移动为 $g$。

\SourceParagraph{ch02-2.7.4-003}{117}
必须考虑，为 Morse 函数临界点所对应的初等配边指定哪些映射。在时间反转意义下，边界临界点处有两种值得讨论的过程：两条开弦首尾相接，或者一条开弦变为闭弦，如\BookRef{fig:2.48}{图48}所示。与前面一样，实线表示临界点之下的等值线，虚线表示其上的等值线。在\BookRef{fig:2.48}{图48}(i)中，$g_a,g_b,g_c$ 是各个 D 膜上相邻点之间的平行移动，用弦上选定的平凡化来表示；它们满足 $g_cg_b=g_a$。为这一情形指定的映射 $\mathcal{O}_{ab}\otimes\mathcal{O}_{bc}\to\mathcal{O}_{ac}$ 为
\[
\psi\otimes\psi'\longmapsto\rho_{g_a}(\psi)\rho_{g_b}(\psi').
\]
对偶操作 $\mathcal{O}_{ac}\to\mathcal{O}_{ab}\otimes\mathcal{O}_{bc}$ 为
\[
\psi\longmapsto\sum_i\rho_{g_1}(\psi\xi^i)\otimes\rho_{g_2}(\xi_i),
\]
其中，$\xi_i,\xi^i$ 分别是 $\mathcal{O}_{bc},\mathcal{O}_{cb}$ 的对偶基；$g_1,g_2$ 表示相应分裂配边中，两条出弦所选平凡化相对于局部参考平凡化的转换元素。

\begin{figure}[H]
\centering
\BookImage{01_Numbered_Figures/Chapter_02/Ch02_Fig48_PDFp116_Bookp106.pdf}
\BookNumberOnlyCaption\label{fig:2.48}
\end{figure}

\SourceParagraph{ch02-2.7.4-004}{117}
在\BookRef{fig:2.48}{图48}的情形(ii)中，开弦变为一条闭弦；相对于图中指定的基点，以及来自开弦起点的平凡化，它的平行移动为 $g$。相应映射是 $\iota^g$；按照\BookEqRef{eq:2.104}{2.104}的配对，其伴随映射为 $\iota_{g^{-1}}$。

\SourceParagraph{ch02-2.7.4-005}{117}
内部临界点还给出两类涉及开弦的操作，如\BookRef{fig:2.49}{图49}所示。与\BookRef{fig:2.49}{图49}(i)对应的映射 $\mathcal{C}_g\otimes\mathcal{O}_{ab}\to\mathcal{O}_{ab}$，为 $\phi\otimes\psi\mapsto\rho_{g_a}\bigl(\iota_g(\phi)\psi\bigr)$；与\BookRef{fig:2.49}{图49}(ii)对应的映射 $\mathcal{O}_{ab}\otimes\mathcal{O}_{cd}\to\mathcal{O}_{ad}\otimes\mathcal{O}_{cb}$，则为
\[
\psi\otimes\psi'\longmapsto
\sum_i\rho_{g_a}(\psi)\psi_i\otimes\rho_{g_c}(\psi')\rho_{g_bg_a^{-1}}(\psi^i),
\]
其中，$\{\psi_i\}$ 是 $\mathcal{O}_{bd}$ 的一组基。

\begin{figure}[H]
\centering
\BookImage{01_Numbered_Figures/Chapter_02/Ch02_Fig49_PDFp116_Bookp106.pdf}
\BookNumberOnlyCaption\label{fig:2.49}
\end{figure}

\SourceParagraph{ch02-2.7.4-006}{117}
现在，需要进行与非等变情形完全相同的各项验证。这些验证十分繁复，但与已经完成的验证一一对应，没有新的内容。作为所需修改的一个例子，指出经常使用的公式\BookEqRef{eq:2.111}{2.111}的范畴版本：当 $\phi'\in\mathcal{O}_{bb}$，且 $\phi^i,\phi_i$ 分别是 $\mathcal{O}_{ab},\mathcal{O}_{ba}$ 的对偶基时，它在任意 Frobenius 范畴中成立。如果范畴带有 $G$ 作用，使用同一个证明，就能将它推广为
\[
\sum_i\phi'\phi_i\otimes\rho_g(\phi^i)
=\sum_i\phi_i\otimes\rho_g(\phi^i\phi'),
\]
对任意 $g\in G$ 都成立。

\SourceParagraph{ch02-2.7.4-007}{117}
关于证明，这里不再多说。
\end{proof}

\section{笔记}\label{sec:2.8}

\SourceParagraph{ch02-2.8-001}{118}
\begin{note}
有关二维 TFT 的文献非常庞大，不可能给出全面的参考文献。这里仅指出一些紧密相关的工作。二维闭弦缝合定理是一个十分古老的结果，它已经隐含在弦论最早的一批论文中。代数表述或许最早由 Friedan 提出。文献\BookCite{ref117,ref410,ref401}，以及 Segal 在斯坦福的讲义\BookCite{ref413}，都给出了相应论述。二维开闭弦理论中的缝合约束，最早由文献\BookCite{ref332}研究。推广到不可定向世界面的情形，见文献\BookCite{ref074,ref007,ref071,ref244}。

\SourceParagraph{ch02-2.8-002}{118}
本章的工作，最初在 Strings 2000 会议上介绍\BookCite{ref361}，并在文献\BookCite{ref362}中作了简要概述。2001年在 KITP 的讲座，以及2002年 Clay 讲习班的讲座\BookCite{ref360}，给出了更完整的说明。文献\BookCite{ref359}提供了不同于这里的、更多借助计算的证明和例子，以及原始图件。我们的一些结果，也由 C. Lazaroiu 在文献\BookCite{ref325,ref323,ref322,ref324}中独立得到；不过，那些论文着重将这些结果用于低能超引力中的圆盘瞬子修正。关于 $G$ 等变理论，D 膜与轨形的文献非常多，并未反映在以上参考文献中。在二维 TFT 的框架下，两个相关参考文献是\BookCite{ref299,ref341}。关于 TFT 中轨形 B 场含义的其他讨论，见文献\BookCite{ref152,ref155,ref428,ref427,ref426}。我们对上链层面理论与 $A_\infty$ 代数的讨论，已经由 Costello 在文献\BookCite{ref100}中作了大幅发展。
\end{note}
